\chapter{The Monster Lie algebra: string states, the no-ghost theorem and the Borcherds--Kac--Moody structure}\label{ch:18}

\providecommand{\chXVIIIm}{\mathfrak{m}}
\providecommand{\chXVIIIg}{\mathfrak{g}}
\providecommand{\chXVIIIu}{\mathfrak{u}}
\providecommand{\chXVIIIp}{\mathfrak{p}}
\providecommand{\chXVIIIl}{\mathfrak{l}}
\providecommand{\chXVIIIh}{\mathfrak{h}}
\providecommand{\chXVIIIII}{\mathrm{II}_{1,1}}
\providecommand{\chXVIIIS}{\mathcal{S}}
\providecommand{\chXVIIIP}{\mathcal{P}}
\providecommand{\chXVIIIK}{\mathcal{K}}
\providecommand{\chXVIIIT}{\mathcal{T}}
\providecommand{\chXVIIIU}{\mathcal{U}}
\providecommand{\chXVIIIN}{\mathcal{N}}

This chapter documents the release of the library named in the front matter; every Lean quotation is checked by
\texttt{scripts/textbook/sync.py} to occur verbatim (up to whitespace and comments) in the source files of that
release.
It is the third chapter of the part ``Monstrous moonshine'' and the central chapter of Borcherds' proof
of the Conway--Norton conjecture.  Chapter~\ref{ch:16} stated the conjecture; Chapter~\ref{ch:17}
built the lattice vertex algebra $V_{\chXVIIIII}$ of the even unimodular Lorentzian lattice of rank two and
the tensor product $\Vnat\otimes V_{\chXVIIIII}$, a vertex algebra of central charge $24+2=26$.  Here we
do three things.
\begin{enumerate}[label=(\Roman*)]
\item We prove the \emph{no-ghost theorem} of Goddard and Thorn in an abstract form (``string
  modules''), and apply it to the explicit momentum sectors $\Vnat\otimes\mathbb C[a_j,b_j]$ of
  $\Vnat\otimes V_{\chXVIIIII}$: for every nonzero momentum $r=(x,y)\in\chXVIIIII$ the space of
  physical states of weight one modulo the radical of the form is isomorphic, $\Aut(\Vnat)$-equivariantly and
  isometrically \emph{for the sector form} (the form of the explicit Fock model, not the vertex-algebra form; see the
  warning after Theorem~\ref{ch18:thm:noghost}), to the weight space $\Vnat_{1-xy}$.
\item We define the \emph{Monster Lie algebra} $\chXVIIIm$ as the quotient of the weight-one
  primary states of $\Vnat\otimes V_{\chXVIIIII}$ by the primary states that are Virasoro
  descendants, with bracket $[\bar u,\bar v]=\overline{u_{(0)}v}$; we prove that it is graded by
  $\chXVIIIII$, that its root spaces are $\chXVIIIm_{(x,y)}\cong\Vnat_{1-xy}$ for $(x,y)\neq 0$ and
  $\chXVIIIm_0\cong\mathbb C^2$, that $\Aut(\Vnat)$ acts on it, and that it carries a symmetric,
  invariant, nondegenerate bilinear form.
\item We develop the Lie theory of an abstract ``Borcherds datum'' (a rank-two graded Lie algebra
  with the root support, form and simple roots of the Monster Lie algebra), and prove the structure
  theorem that Chapter~\ref{ch:19} needs: the length-one resolution of the trivial module by
  generalized Verma modules, the isomorphism $T(U)\cong \mathcal U(\chXVIIIu^+)$, and the freeness of the
  positive cone $\chXVIIIu^+$ on the space $U$ of real-root strings of the imaginary simple root
  spaces.  We prove that $\chXVIIIm$ is such a datum, \emph{given} that it is generated by its simple
  root spaces; that last input is a theorem of Chapter~\ref{ch:20} (it rests on the product formula
  of Chapter~\ref{ch:19}).
\end{enumerate}

As everywhere in this book, our concern is the \emph{meaning} of the Lean statements.  Proofs are
explained mathematically (and, where it helps, physically); we do not describe how the Lean proof
scripts are organised.

\paragraph{The results in one paragraph.}  The no-ghost isomorphism, the Lie algebra $\chXVIIIm$, its root-space
isomorphisms, its invariant nondegenerate form and the whole abstract Lie theory are proved in Lean with \emph{no
hypotheses} beyond the data they are about.  The single exception is the statement ``$\chXVIIIm$ is a Borcherds datum''
(\S\ref{ch18:sec:monsterbd}) and its consequences for $\chXVIIIm$, which take generation of $\chXVIIIm$ by its simple root spaces
as an explicit argument; that argument is itself a theorem with no hypotheses, proved in Chapter~\ref{ch:20}.  The Lean
names are collected in the referee's checklist (\S\ref{ch18:sec:checklist}).

\paragraph{Reading the Lean in this chapter.}  (Chapter~\ref{ch:01} explains the conventions in detail; we recall the
ones used here.)  \texttt{\#print axioms X} lists the logical axioms a declaration \lean{X} ultimately depends on;
``standard axioms'' means exactly Lean's \lean{propext}, \lean{Classical.choice} and \lean{Quot.sound}, with no
\lean{sorryAx} (an unfinished proof) and no project axiom.  Every status box names the audit module in which the
\texttt{\#print axioms} line for the declaration occurs; the outputs are those reported by the maintainers' targeted
builds, they were not re-run for this chapter, and no aggregate build of the whole library is claimed.  In the
quotations, a statement ending in \texttt{:=} or \texttt{:= by} has its proof omitted; a line \texttt{...} marks
omitted lines (the omission is always described in the text).  \lean{def} introduces a definition, \lean{abbrev} a
definition that Lean unfolds automatically, \lean{theorem} a proposition with proof, \lean{structure} a record type
whose fields are listed; \lean{noncomputable} only says the definition uses choice and has no executable code, and
\lean{instance} registers a structure (here: a Lie bracket) for automatic use.  A proof written \lean{rfl} means the
two sides are equal by unfolding definitions.  Lean type variables are named freely; in particular, in the induced-module
files of \S\ref{ch18:sec:verma} the Lean variable \lean{U} denotes an arbitrary module (called $W$ in the text), not the
generating space $U$.

\paragraph{What is \emph{not} in this chapter.}  (a) No BRST complex: the physical states are
defined in the ``old covariant'' way (Virasoro primaries modulo spurious states), and the BRST
description is a physical comment only.  (b) No positivity: the form on $\Vnat$ is never shown to be
positive definite, and nothing here needs it.  (c) No generalized Kac--Moody \emph{characterisation}
theorem: we do not prove that $\chXVIIIm$ satisfies Borcherds' axioms for a generalized Kac--Moody
algebra, nor any Weyl--Kac--Borcherds character formula; we prove exactly the structural statements
listed in (III).  (d) The statement ``$E=\bigoplus_{m>0}\chXVIIIm_{(m,n)}$ minus the real simple root
is free on $\bigoplus_n\chXVIIIm_{(1,n)}$'' in its literal form is not what is proved; see
\S\ref{ch18:sec:free} for the exact freeness statement and why it is the right one.

\section{Orientation: what is being built, and the chain of results}\label{ch18:sec:intro}

\subsection{Why a Lie algebra?}

Borcherds' proof of the Conway--Norton conjecture runs through a Lie algebra.  The McKay--Thompson
series $T_g(q)=q^{-1}+\sum_{n\ge1}\tr(g\mid\Vnat_{n+1})q^n$ of Chapter~\ref{ch:16} are to be shown to be
Hauptmoduln.  Borcherds' strategy: (1) build a Lie algebra $\chXVIIIm$ graded by the lattice
$\chXVIIIII$ whose root spaces are, as $\Aut(\Vnat)$-modules, the homogeneous pieces of $\Vnat$; (2)
show that $\chXVIIIm$ has a ``Kac--Moody-like'' structure with known simple roots, so that it satisfies
a \emph{denominator identity}; (3) take $g$-traces of the denominator identity (the \emph{twisted}
denominator identities), which yield the \emph{replication formulas} for the $T_g$; (4) show that
completely replicable functions with the right leading coefficients are Hauptmoduln.  This chapter
does (1) and the Lie-theoretic heart of (2); Chapter~\ref{ch:19} does the rest of (2) and (3); the
genus-zero step (4) is in Chapter~\ref{ch:19} as well, and the final assembly in Chapter~\ref{ch:20}.

\subsection{The chain of results at a glance}\label{ch18:sec:chain}

The symbols in the following remark are defined in \S\S\ref{ch18:sec:stringmodule}--\ref{ch18:sec:abstractnoghost}; informally:
a \emph{string module} $H$ is a graded space with a Virasoro action, an abelian family of ``longitudinal'' operators
$K_n$ and an invariant form; $\chXVIIIS$ is the span of all Virasoro descendants $L_{-n}H$ (\emph{spurious} states);
$\chXVIIIP^1$ the weight-one states killed by all $L_n$, $n\ge1$ (\emph{physical} states); $\chXVIIIK$ the states killed by all
$K_n$, $n\ge1$ (\emph{longitudinal vacua}); $\chXVIIIT^1=\chXVIIIP^1\cap\chXVIIIK$ (\emph{transverse} states); and ``rad'' is the radical
of the form restricted to the space in question.

\begin{remark}[The chain]\label{ch18:rem:chain}
Each arrow is a Lean theorem presented below; the hypotheses in parentheses are discharged where
indicated.
\begin{align*}
  &\text{abstract string module } H,\ c=26,\ H=\chXVIIIS\oplus\chXVIIIK
   \;\Longrightarrow\; \chXVIIIP^1/\mathrm{rad}\cong \chXVIIIT^1\cong \chXVIIIK^1/\mathrm{rad}
   && \text{(\S\ref{ch18:sec:abstractnoghost})}\\
  &\Vnat\otimes\mathbb C[a,b]\text{ at momentum }(x,y)\neq0
   \text{ is such a module, with } \chXVIIIS\oplus\chXVIIIK \text{ proved}
   && \text{(\S\ref{ch18:sec:sector})}\\
  &\chXVIIIK^1/\mathrm{rad}\cong\Vnat_{1-xy}
   \;\Longrightarrow\;
   \chXVIIIP^1_{(x,y)}/\mathrm{rad}\cong \Vnat_{1-xy}
   && \text{(\lean{stringSectorNoGhostEquiv}, \S\ref{ch18:sec:vnatnoghost})}\\
  &\chXVIIIm := P^1(\Vnat\otimes V_{\chXVIIIII})/(P^1\cap\text{spurious}),\
   [\bar u,\bar v]=\overline{u_{(0)}v}
   && \text{(\lean{monsterLieAlgebra}, \S\ref{ch18:sec:lie})}\\
  &\chXVIIIm_{(x,y)}\cong\Vnat_{1-xy}\ ((x,y)\neq0),\quad \chXVIIIm_0\cong\mathbb C^2
   && \text{(\lean{monsterLieRootSpaceEquiv}, \S\ref{ch18:sec:roots})}\\
  &\text{form on }\chXVIIIm\text{: symmetric, invariant, nondegenerate}
   && \text{(\S\ref{ch18:sec:form})}\\
  &\text{abstract Borcherds datum }\Longrightarrow\ 0\to \mathrm{Ind}(U)\to \mathrm{Ind}(\mathbb C)\to\mathbb C\to0,\
   T(U)\cong \mathcal U(\chXVIIIu^+),\ \chXVIIIu^+ \text{ free on } U
   && \text{(\S\S\ref{ch18:sec:bd}--\ref{ch18:sec:free})}\\
  &\chXVIIIm\text{ is a Borcherds datum } (\text{given } \lean{MonsterLieSimpleRootGeneration})
   && \text{(\lean{monsterLieBorcherdsData}, \S\ref{ch18:sec:monsterbd})}.
\end{align*}
\end{remark}

\subsection{Conventions that are fixed once and for all}\label{ch18:sec:conventions}

\paragraph{The lattice.}  $\chXVIIIII$ is $\Z^2$ with the bilinear form
\begin{equation}\label{ch18:eq:repform}
  \bigl((x,y),(x',y')\bigr) \;=\; xy'+x'y,\qquad (x,y)^2=2xy,
\end{equation}
the Gram matrix $\left(\begin{smallmatrix}0&1\\1&0\end{smallmatrix}\right)$.  In Lean this is
\lean{lorentzianLattice} (Chapter~\ref{ch:17}); a lattice vector is written \lean{![x, y]} and the
type of lattice vectors is \lean{lorentzianLattice.carrier}.  We call $(x,y)$ the \emph{repository
momentum}.  For a momentum $r=(x,y)$ the half-norm is $r^2/2=xy$, an integer because the lattice is
even.

\paragraph{Borcherds coordinates.}  Borcherds (and the Lie-theoretic part of this chapter) uses the
form $((m,n),(m',n'))_{\mathrm{Bor}}=-(mn'+m'n)$, so $(m,n)^2_{\mathrm{Bor}}=-2mn$.  The dictionary is
\begin{equation}\label{ch18:eq:dictionary}
  \text{Borcherds root }(m,n)\;\longleftrightarrow\;\text{repository momentum }(x,y)=(m,-n),
\end{equation}
which is an isometry: $-(mn'+m'n)=m(-n')+m'(-n)$.  In Lean the map is \lean{borcherdsMomentum}
(\S\ref{ch18:sec:monsterbd}).  Under \eqref{ch18:eq:dictionary} the root-space theorem reads
\begin{equation*}
  \chXVIIIm_{(x,y)}\cong\Vnat_{1-xy}\quad\text{(repository)}\qquad\Longleftrightarrow\qquad
  \chXVIIIm_{(m,n)}\cong\Vnat_{1+mn}\quad\text{(Borcherds)}.
\end{equation*}
The real simple root is $\alpha_r=(1,-1)$ (norm $2$; repository momentum $(1,1)$) and the imaginary
simple roots are $\alpha_j=(1,j)$, $j\ge1$ (norm $-2j$; repository momentum $(1,-j)$), with root
space $\chXVIIIm_{\alpha_j}\cong\Vnat_{1+j}$, of dimension $c(j)$, the $j$-th coefficient of
$J(q)=q^{-1}+\sum_{n\ge1}c(n)q^n$.

\begin{warning}
\textbf{The sign trap $(m,n)\leftrightarrow(m,-n)$.}  A product over repository momenta
$x,y\ge1$ (a tempting first transcription of Borcherds' formulas) is \emph{wrong}: for $x,y\ge 1$ the root space is $\Vnat_{1-xy}$, which vanishes as soon as
$xy\ge2$ (and is $\Vnat_0=\mathbb C$ only for $x=y=1$, the real root).  The positive cone of
Borcherds' theory, where all the imaginary roots live, is $\{(m,n):m,n\ge1\}$ in \emph{Borcherds}
coordinates, i.e.\ repository momenta $(x,y)$ with $x\ge1$, $y\le-1$.  Every statement of this
chapter is labelled with the coordinates it uses, and the bridge \lean{monsterLieBorcherdsRootEquiv}
(\S\ref{ch18:sec:monsterbd}) states the Borcherds form of the root-space theorem once, in Lean.
\end{warning}

\paragraph{Modes.}  The library convention \lean{mode f n = f (-n - 1)} (Chapter~\ref{ch:L01},
Definition~\ref{l01:def:mode}) is in force: \lean{mode (V.Y u) n} is the mathematician's $u_{(n)}$, the
coefficient of $z^{-n-1}$ in $Y(u,z)$.  For a near vertex operator algebra \lean{V} (the structure
\lean{VOA.NearVOA} of Chapter~\ref{ch:L03}: a vertex operator algebra without finite-dimensional
weight spaces and without a lower bound on weights), \lean{V.smul u n v} abbreviates
\lean{mode (V.Y u) n v}, i.e.\ $u_{(n)}v$, and \lean{V.L n} is the Virasoro mode
$L_n=\omega_{(n+1)}$:
\begin{leancode}
noncomputable def smul (V : NearVOA 𝒱) (u : 𝒱) (n : ℤ) (v : 𝒱) : 𝒱 := mode (V.Y u) n v
\end{leancode}
(\leanpath{VOA/Moonshine/MoonshineCNNearGoddard.lean}; the same file defines \lean{NearVOA.L}.)

\paragraph{The moonshine module and its form.}  $\Vnat$ is the Lean type \lean{Vnat}, the vertex
operator algebra is \lean{VnatVOA} (Chapter~\ref{ch:14}; \leanpath{VOA/Moonshine/AutVnat.lean}), its
weight-$n$ space is \lean{VnatVOA.grading n}, and $\Aut(\Vnat)$ is \lean{AutVnat := Aut VnatVOA}.  The
invariant bilinear form on $\Vnat$ is the block form \lean{natForm}
(\leanpath{VOA/Moonshine/SimpleForm.lean}):
\begin{leancode}
noncomputable def natForm : LinearMap.BilinForm ℂ Vnat :=
  leechBV.compl₁₂ (LinearMap.fst ℂ SkewVP TP) (LinearMap.fst ℂ SkewVP TP)
    + TPForm.compl₁₂ (LinearMap.snd ℂ SkewVP TP) (LinearMap.snd ℂ SkewVP TP)
\end{leancode}
that is, $B^\natural((u,t),(u',t'))=B_V(u,u')+B_T(t,t')$ on $\Vnat=V_\Leech^+\oplus(V_\Leech^T)^+$.  It is
symmetric, satisfies $B^\natural(L_nu,v)=B^\natural(u,L_{-n}v)$, is nondegenerate, and its weight
spaces are mutually orthogonal (Chapters~\ref{ch:09}--\ref{ch:10}).  The restriction of $g\in\Aut\Vnat$
to $\Vnat_n$ is \lean{GriessAut.autRestrict g n} (\leanpath{VOA/Core/GriessAutomorphism.lean}):
\begin{leancode}
noncomputable def autRestrict (α : Aut W) (n : ℤ) : W.grading n ≃ₗ[ℂ] W.grading n :=
  α.toLinearEquiv.ofSubmodules _ _ (aut_map_grading α n)
\end{leancode}

\subsection{Notation}\label{ch18:sec:notation}

Many letters are traditionally used for several objects in this subject.  The following table fixes the meaning of
each symbol in this chapter; where a Lean name uses a letter differently, the text says so.

\begin{longtable}{p{0.22\textwidth}p{0.72\textwidth}}
\toprule
symbol & meaning \\ \midrule
\endhead
$H$, $H^N$, $N\in\Z$ & a string module and its weight-$N$ space (\S\ref{ch18:sec:stringmodule}); $N$ always denotes a weight\\
$L_n$, $K_n$ & Virasoro operators; longitudinal operators (one Heisenberg field along a null direction)\\
$B$ & the form of a string module; $B^\natural$ the form \lean{natForm} on $\Vnat$; $B_F$ the monomial form on $F$\\
$(\cdot,\cdot)$ & the lattice form \eqref{ch18:eq:repform} on $\chXVIIIII$ (repository coordinates)\\
$(\cdot,\cdot)_{\mathrm{Bor}}$ & the form $-(mn'+m'n)$ in Borcherds coordinates; $\chXVIIIm^{\mathrm{Bor}}_\gamma$ the root space in Borcherds coordinates\\
$\chXVIIIS,\chXVIIIS_K,\chXVIIIP,\chXVIIIK,\chXVIIIT$ & spurious, longitudinal spurious, physical, longitudinal-vacuum, transverse subspaces\\
$F=\mathbb C[a_j,b_j]$ & the polynomial (Fock) space of the two light-cone bosons; $a_j,b_j$ its variables\\
$w_0=(0,1)$, $w_1=(1,0)$ & the isotropic basis of $\chXVIIIII$ whose oscillators are $a$ and $b$\\
$\kappa,\ell$ & the two charges (light-cone momenta) of a sector; $\kappa=x$, $\ell=y$ at momentum $(x,y)$\\
$e^r$ & the lattice state of momentum $r$ in $V_{\chXVIIIII}$\\
$P^1(V)$, $\mathrm{Sp}(V)$, $\mathrm{Null}^1(V)$ & weight-one primaries, spurious states, and their intersection, in a near-VOA $V$\\
$\chXVIIIm$, $\chXVIIIg$ & the Monster Lie algebra; an abstract Lie algebra with Borcherds datum\\
$\alpha_r$, $\alpha_j$ & the real simple root $(1,-1)$ and the imaginary simple roots $(1,j)$ (Borcherds coordinates)\\
$e,f$ & a normalized pair of real root vectors, $e\in\chXVIIIg_{\alpha_r}$, $f\in\chXVIIIg_{-\alpha_r}$, $(e,f)=1$\\
$\delta=(1,0)$ & the Weyl vector\\
$\chXVIIIu^\pm,\chXVIIIp^-,\chXVIIIl$ & positive cone, opposite parabolic, Levi subalgebra (\S\ref{ch18:sec:strings})\\
$U$, $U_{j,k}$ & the generating space of real-root strings and its pieces (\S\ref{ch18:sec:strings})\\
$\mathcal \mathcal U(L)$, $T(S)$ & universal enveloping algebra of a Lie algebra $L$; tensor algebra of a vector space $S$\\
$\mathrm{Ind}(W)$ & the module induced from a $\chXVIIIp^-$-module $W$ (generalized Verma module)\\
$\rho$ & a representation $\chXVIIIg\to\End M$ (never the Weyl vector)\\
$\Omega$, $Q(\gamma)$ & the Casimir operator and its scalar $(\gamma,\gamma)_{\mathrm{Bor}}-2(\delta,\gamma)_{\mathrm{Bor}}$\\
\bottomrule
\end{longtable}

\section{Physical background: the bosonic string and its physical states}\label{ch18:sec:physics}

This section is entirely physics; nothing in it is used as a mathematical input.  It explains why the
definitions of the following sections are the natural ones, and where the number $26$ comes from.

\subsection{The bosonic string in 26 dimensions}

A closed bosonic string moving in $D$-dimensional Minkowski space is described, after fixing the
conformal gauge on its world-sheet, by $D$ free bosons $X^\mu(z,\bar z)$, $\mu=0,\dots,D-1$.  Each
chiral half is a Heisenberg vertex algebra of rank $D$ with the Lorentzian form $\eta_{\mu\nu}$ of
signature $(1,D-1)$, and its Sugawara stress tensor has central charge $c=D$.  The world-sheet metric
was gauge-fixed, and the remnant of its equations of motion is the constraint that the stress tensor
annihilates physical states.  In the \emph{old covariant quantization} this is imposed as
\begin{equation}\label{ch18:eq:virasoroconstraints}
  L_n\,\psi = 0\ (n\ge1),\qquad (L_0-1)\,\psi=0,
\end{equation}
the ``$-1$'' being the normal-ordering constant of the $bc$ ghosts.  The Fock space carries an
indefinite inner product (the timelike oscillators $\alpha^0_{-n}$ create negative-norm states, the
``ghosts'').  The states \eqref{ch18:eq:virasoroconstraints} are the \emph{physical states}; among them
are the \emph{null} physical states, which are also of the form $L_{-n}\chi$ (``spurious'') and have
zero inner product with every physical state.  The \emph{no-ghost theorem} of Goddard and Thorn (1972)
asserts that when $D=26$ the inner product is positive semi-definite on physical states and that its
radical is exactly the null physical states; moreover the quotient is isomorphic to the space of
\emph{transverse} states, built from the $D-2=24$ transverse oscillators only (the light-cone gauge
spectrum).  For $D>26$ there are negative-norm physical states; for $D<26$ the quotient is bigger than
the light-cone count.

The BRST formulation (Kato--Ogawa 1983, Frenkel--Garland--Zuckerman 1986) packages the same statement
as the vanishing of relative semi-infinite cohomology outside degree one: the BRST charge
$Q=\sum c_{-n}L_n+\dots$ satisfies $Q^2=0$ exactly when the matter central charge is $26$, and
$H^1_{\mathrm{rel}}\cong\chXVIIIP^1/\text{null}$.  The library does \emph{not} formalize the BRST complex
for this chapter; it formalizes the old covariant version, with Borcherds' simplifications.

\subsection{\texorpdfstring{Compactification on $\Vnat\otimes V_{\chXVIIIII}$}{Compactification on V-natural tensor V(II11)}}

The no-ghost argument does not need $D$ flat directions.  It needs: (i) a two-dimensional
Lorentzian lattice factor $V_{\chXVIIIII}$ (the ``light-cone directions'' $X^\pm$, whose momenta $r\in\chXVIIIII$
are the string's momenta in those two directions), and (ii) an arbitrary ``internal'' chiral CFT $V$ of
central charge $24$ with a nondegenerate invariant form.  The total central charge is $24+2=26$.  Taking
$V=\Vnat$ means: a bosonic string whose internal $24$ dimensions are replaced by the Monster CFT and
whose two remaining directions are compactified on the even unimodular lattice $\chXVIIIII$.  Its
physical states of momentum $r=(x,y)$ satisfy $(L_0-1)\psi=0$, where $L_0$ is the internal weight plus
oscillator level plus $r^2/2=xy$.  The no-ghost theorem identifies them, modulo null states, with the
internal states of weight $1-xy$ \emph{with no light-cone oscillators}: $\Vnat_{1-xy}$.

\subsection{The Lie algebra of physical states}

Physical states of weight one are \emph{conserved currents}: a weight-one primary $u$ has a zero mode
$u_{(0)}$ that commutes with all $L_n$ and therefore maps physical states to physical states.  In
string language, $u_{(0)}$ is the charge of the vertex operator $\oint V_u(z)\,dz$, and $[\bar u,\bar
v]=\overline{u_{(0)}v}$ is the commutator of two such charges (the ``algebra of string vertex
operators'' of Frenkel, and of Goddard--Olive; in the BRST language the bracket on $H^1$).  Null
states decouple, so the bracket descends to physical states modulo null states.  The resulting Lie
algebra, for $V=\Vnat$, is the \emph{Monster Lie algebra}: the spacetime gauge symmetry algebra of the
compactified string, with the Monster acting through its action on the internal CFT.  Its Cartan
subalgebra $\chXVIIIm_0\cong\chXVIIIII\otimes\mathbb C$ is spanned by the two momentum currents
$\partial X^\pm$; the root space at momentum $r$ consists of the physical states of momentum $r$.  In
Borcherds' words, $\chXVIIIm$ is ``the space of physical states of a bosonic string moving on the
$26$-dimensional orbifold $\Vnat\times(\mathbb R^{1,1}/\chXVIIIII)$''.

\subsection{Simple roots as BPS-like states}

The real simple root $\alpha_r$ (repository momentum $(1,1)$, $r^2=2$, root space $\Vnat_0=\mathbb
C\mathbb 1$) is the \emph{tachyon} of momentum $(1,1)$: the state $\mathbb 1\otimes e^{(1,1)}$, with
no oscillators.  It generates, with its opposite, an $\mathfrak{sl}_2$ (a ``W-boson'' pair).  The
imaginary simple roots $\alpha_j$ (repository momenta $(1,-j)$, $r^2=-2j<0$) are the states
$v\otimes e^{(1,-j)}$ with $v\in\Vnat_{1+j}$: massive internal states balanced by a timelike momentum.
That the positive part of $\chXVIIIm$ is generated freely by the $\mathfrak{sl}_2$-strings of these
states (\S\ref{ch18:sec:free}) is, physically, the statement that string scattering in this background
produces no relations among the simple-root states beyond those forced by the $\mathfrak{sl}_2$
symmetry.  This is a heuristic description, not an input.

\section{Abstract string modules}\label{ch18:sec:stringmodule}

The no-ghost theorem is, at its core, linear algebra in a graded vector space carrying two commuting
structures: a Virasoro action, and an abelian family of ``longitudinal'' operators that the Virasoro
algebra normalises.  Following Borcherds' formulation of Goddard--Thorn (Borcherds 1992, §5), the library abstracts
exactly these data into a structure \lean{StringModule}.  Nothing in this section uses a vertex algebra; the file
depends on Mathlib only.

\paragraph{Why the operators $K_n$.}  In the string Fock space at momentum $r$ there are two light-cone bosons.  Pick a
\emph{null} lattice direction $w$ (so $(w,w)=0$) with $(w,r)\neq0$, and let $K_n:=w(n)$ be the modes of the Heisenberg field
along $w$.  Because $w$ is null, the $K_n$ commute with each other; because $w(z)$ is a weight-one primary current,
$[L_m,K_n]=-nK_{m+n}$.  The point of these operators is the \emph{light-cone splitting}: when $(w,r)\neq0$, every state is
uniquely a Virasoro descendant plus a state killed by all $K_n$, $n\ge1$ (hypothesis (F1) below, proved for the concrete
sector in Theorem~\ref{ch18:thm:F1}, where $K_n$ is the explicit oscillator $a(n)$).  The states killed by the $K_n$ are
built only from the other light-cone boson and the internal CFT, and modulo null states they are the light-cone
(transverse) states.  So the $K_n$ are the bookkeeping device that turns ``covariant'' into ``light-cone'' states.

\subsection{The definition}

\begin{definition}[String module of central charge $c$]\label{ch18:def:stringmodule}
A \emph{string module} is a complex vector space $H$ with
\begin{enumerate}[label=(S\arabic*)]
\item a grading $H=\bigoplus_{N\in\Z}H^N$ (internal direct sum) with every $H^N$ finite-dimensional and
  $H^N=0$ for $N<N_0$;
\item operators $L_n$ ($n\in\Z$) with $L_nH^N\subseteq H^{N-n}$, $L_0=N$ on $H^N$, and
  $[L_m,L_n]=(m-n)L_{m+n}+\frac{c}{12}(m^3-m)\delta_{m+n,0}$;
\item operators $K_n$ ($n\in\Z$; $K_0$ is never used) with $K_nH^N\subseteq H^{N-n}$ for $n\neq0$,
  $[K_m,K_n]=0$, and $[L_m,K_n]=-nK_{m+n}$ whenever $n\neq0$ and $m+n\neq0$;
\item a symmetric bilinear form $B$ with $B(H^N,H^M)=0$ for $N\ne M$, nondegenerate on each $H^N$, and
  $B(L_nx,y)=B(x,L_{-n}y)$, $B(K_nx,y)=B(x,K_{-n}y)$ for all $n,x,y$.
\end{enumerate}
\end{definition}

\begin{leancode}
structure StringModule (H : Type u) [AddCommGroup H] [Module ℂ H] where
  wt : ℤ → Submodule ℂ H
  wt_isInternal : DirectSum.IsInternal wt
  wt_finite : ∀ N, FiniteDimensional ℂ (wt N)
  wt_bdd : ∃ N₀, ∀ N < N₀, wt N = ⊥
  c : ℂ
  L : ℤ → Module.End ℂ H
  K : ℤ → Module.End ℂ H
  L_wt : ∀ n N x, x ∈ wt N → L n x ∈ wt (N - n)
  K_wt : ∀ n N x, n ≠ 0 → x ∈ wt N → K n x ∈ wt (N - n)
  L_zero : ∀ N x, x ∈ wt N → L 0 x = (N : ℂ) • x
  L_comm : ∀ m n, L m * L n - L n * L m =
    ((m : ℂ) - n) • L (m + n) +
      if m + n = 0 then (((m : ℂ) ^ 3 - m) / 12 * c) • 1 else 0
  K_comm : ∀ m n, K m * K n = K n * K m
  LK_comm : ∀ m n, n ≠ 0 → m + n ≠ 0 →
    L m * K n - K n * L m = (-(n : ℂ)) • K (m + n)
  B : LinearMap.BilinForm ℂ H
  B_symm : B.IsSymm
  B_orth : ∀ N M x y, N ≠ M → x ∈ wt N → y ∈ wt M → B x y = 0
  B_nondeg : ∀ N x, x ∈ wt N → (∀ y ∈ wt N, B x y = 0) → x = 0
  L_adj : ∀ n x y, B (L n x) y = B x (L (-n) y)
  K_adj : ∀ n x y, B (K n x) y = B x (K (-n) y)
\end{leancode}
(\leanpath{VOA/Moonshine/MoonshineCNStringModule.lean}, \lean{VOA.Moonshine.ConwayNorton.StringModule}.)

\begin{dictionary}
\textbf{Decoding.}  \lean{wt N} is the subspace $H^N$; \lean{DirectSum.IsInternal wt} says the
canonical map $\bigoplus_N H^N\to H$ is bijective (every vector is a unique finite sum of homogeneous
components).  \lean{Module.End ℂ H} is the algebra of $\mathbb C$-linear endomorphisms, with
\lean{*} the composition; \lean{L m * L n - L n * L m} is $[L_m,L_n]$.  The central term is written
\lean{(((m : ℂ) ^ 3 - m) / 12 * c) • 1} when \lean{m + n = 0}, i.e.\ $\frac{c}{12}(m^3-m)\,\mathrm{id}$,
the standard normalisation (central charge $c$ means $[L_2,L_{-2}]=4L_0+c/2$).
\lean{LinearMap.BilinForm ℂ H} is a bilinear form $H\times H\to\mathbb C$; \lean{B.IsSymm} is symmetry.

\textbf{Hypotheses, field by field.}  (S1) is \lean{wt_isInternal}, \lean{wt_finite}, \lean{wt_bdd};
(S2) is \lean{L_wt}, \lean{L_zero}, \lean{L_comm}; (S3) is \lean{K_wt}, \lean{K_comm},
\lean{LK_comm}; (S4) is \lean{B_symm}, \lean{B_orth}, \lean{B_nondeg}, \lean{L_adj}, \lean{K_adj}.
Two places deserve comment.  (i) \lean{K_wt} is required only for $n\neq0$ and \lean{LK_comm} only
for $n\neq0$, $m+n\neq0$: the value of $K_0$ is immaterial (in the concrete sector it is the scalar
momentum, and $[L_m,K_{-m}]=mK_0$ would hold too, but nothing needs it).  The module docstring
says so (``The value of \lean{K 0} is immaterial: mixed commutators exclude it'').  (ii) \lean{B_nondeg}
is nondegeneracy \emph{on each weight space}, which, together with \lean{B_orth} and finite dimension,
is what the finite-dimensional linear algebra below needs.

\textbf{Two things are deliberately not fields.}  The optional group action is a separate structure
(\lean{StringGroupAction}, below), and the light-cone hypothesis (F1) is a separate \lean{Prop}, so that every theorem
displays exactly where it is used.

\textbf{Could it be vacuous?}  No: the structure is instantiated in \S\ref{ch18:sec:sector} by the
concrete sectors \lean{vnatStringModule x y a}, whose fields are all proved.  The axioms are
consistent and are satisfied by a nonzero example.
\end{dictionary}

\begin{physics}
$H$ is the Fock space of the string at fixed momentum $r$: the matter CFT tensored with the oscillators
of the two light-cone bosons.  $L_n$ are the total Virasoro generators (matter plus light-cone), $K_n$
are the modes $w\cdot\alpha_n$ of one light-cone boson along an isotropic direction $w$ with $(w,r)\neq0$
(Borcherds' choice; in light-cone gauge $\alpha^+_n$).  $[L_m,K_n]=-nK_{m+n}$ says that $K$ is a
weight-one primary current, and $B$ is the (indefinite!) string inner product, for which $L_n^\dagger=L_{-n}$.
\end{physics}

\subsection{The derived subspaces}

\begin{definition}[Spurious, physical, longitudinal and transverse states]\label{ch18:def:derived}
In a string module put
\begin{align*}
  \chXVIIIS &:= \sum_{n\ge1}L_{-n}H &&\text{(spurious states)}, &
  \chXVIIIS_K &:= \sum_{n\ge1}K_{-n}H &&\text{(longitudinal spurious states)},\\
  \chXVIIIP &:= \bigcap_{n\ge1}\ker L_n &&\text{(physical states)}, &
  \chXVIIIK &:= \bigcap_{n\ge1}\ker K_n &&\text{(longitudinal vacua)},
\end{align*}
$\chXVIIIT:=\chXVIIIP\cap\chXVIIIK$ (transverse states), and say that $H$ satisfies the \emph{light-cone splitting}
(F1) if $H=\chXVIIIS\oplus\chXVIIIK$.  For a subspace $U\subseteq H$, $U^N:=U\cap H^N$ viewed inside $H^N$, and $B^N$ is
the restriction of $B$ to $H^N$.
\end{definition}

\begin{leancode}
def spurious : Submodule ℂ H := ⨆ n : ℕ, LinearMap.range (M.L (-(n + 1 : ℤ)))
def kSpurious : Submodule ℂ H := ⨆ n : ℕ, LinearMap.range (M.K (-(n + 1 : ℤ)))
def physical : Submodule ℂ H := ⨅ n : ℕ, LinearMap.ker (M.L (n + 1 : ℤ))
def kVacuum : Submodule ℂ H := ⨅ n : ℕ, LinearMap.ker (M.K (n + 1 : ℤ))
def transverse : Submodule ℂ H := M.physical ⊓ M.kVacuum
def LightConeSplit : Prop := IsCompl M.spurious M.kVacuum
def weightSlice (U : Submodule ℂ H) (N : ℤ) : Submodule ℂ (M.wt N) :=
  U.comap (M.wt N).subtype
def stringWeightForm (N : ℤ) : LinearMap.BilinForm ℂ (M.wt N) :=
  M.B.restrict (M.wt N)
\end{leancode}
(\leanpath{VOA/Moonshine/MoonshineCNStringModule.lean}, namespace
\lean{VOA.Moonshine.ConwayNorton.StringModule}, with \lean{M : StringModule H} a section variable.)

\begin{dictionary}
\textbf{Decoding.}  \lean{⨆ n : ℕ, ...} is the supremum (sum) of subspaces over $n\in\N$, and
\lean{-(n + 1 : ℤ)} runs over $-1,-2,\dots$; so \lean{spurious} is $\sum_{n\ge1}\mathrm{im}\,L_{-n}$.
\lean{⨅} is the intersection; \lean{ker (M.L (n + 1))} over $n\in\N$ is $\bigcap_{n\ge1}\ker L_n$.
\lean{⊓} is intersection.  \lean{IsCompl A B} means $A\cap B=0$ and $A+B=H$.  \lean{U.comap (M.wt
N).subtype} is the preimage of $U$ under the inclusion $H^N\hookrightarrow H$, i.e.\ $U\cap H^N$ as a
subspace \emph{of} $H^N$ (so that finite-dimensional linear algebra applies).  \lean{B.restrict}
restricts the form.

\textbf{Note.}  These are the subspaces of Definition~\ref{ch18:def:derived}.  Note that
\lean{physical} is \emph{not} restricted to weight one; the weight-one physical space is
\lean{M.weightSlice M.physical 1}.  (The $L_0$-condition $(L_0-1)\psi=0$ of
\eqref{ch18:eq:virasoroconstraints} is imposed by taking the weight-one slice.)
\end{dictionary}

\begin{physics}
$\chXVIIIP^1$ is the space of Virasoro-primary states of weight one, the physical states of old covariant
quantization.  $\chXVIIIS$ are the spurious states $L_{-n}\chi$; $\chXVIIIP^1\cap\chXVIIIS$ are the null physical states.
$\chXVIIIK$ consists of the states with no excitation of the longitudinal oscillator $\alpha^+$ conjugate to
$K$; $\chXVIIIT=\chXVIIIP\cap\chXVIIIK$ are the DDF-like ``transverse'' states.  (F1) is the precise form of ``every state
is uniquely a longitudinal-vacuum state dressed by Virasoro descendants''.
\end{physics}

\begin{definition}[Symmetries]\label{ch18:def:stringgroup}
A \emph{symmetry} of a string module by a group $G$ is a homomorphism $G\to\mathrm{GL}(H)$ whose
elements preserve each $H^N$, commute with every $L_n$ and every $K_n$, and preserve $B$.
\end{definition}
\begin{leancode}
structure StringGroupAction {H : Type u} [AddCommGroup H] [Module ℂ H]
    (M : StringModule H) (G : Type*) [Group G] where
  action : G →* (H ≃ₗ[ℂ] H)
  preserves_wt : ∀ g N x, x ∈ M.wt N → action g x ∈ M.wt N
  commutes_L : ∀ g n x, action g (M.L n x) = M.L n (action g x)
  commutes_K : ∀ g n x, action g (M.K n x) = M.K n (action g x)
  isometry : ∀ g x y, M.B (action g x) (action g y) = M.B x y
\end{leancode}
(\leanpath{VOA/Moonshine/MoonshineCNStringModule.lean}, \lean{VOA.Moonshine.ConwayNorton.StringGroupAction}.)
\begin{dictionary}
\lean{G →* (H ≃ₗ[ℂ] H)} is a group homomorphism into the group of linear automorphisms of $H$.  The
four laws say that $g$ preserves weights, commutes with all $L_n$ and $K_n$, and is an isometry.  This is the optional structure used to state equivariance; for
$\Vnat$ the group is \lean{AutVnat} acting on the matter factor (\S\ref{ch18:sec:sector}).
\end{dictionary}
\begin{physics}
An internal global symmetry of the matter CFT that does not touch the light-cone directions: it
commutes with the total stress tensor and with the light-cone currents, and is unitary for the string
inner product.  For $\Vnat$ this is the Monster.
\end{physics}
\begin{status}
Definitions; no axioms involved.  Instantiated, with all fields proved, by \lean{vnatStringModule} and
\lean{vnatStringGroupAction} (\S\ref{ch18:sec:sector}).
\end{status}

\section{Linear algebra in a string module, before any splitting}\label{ch18:sec:linalg}

This section collects what holds in \emph{every} string module, at \emph{every} central charge, with no
light-cone hypothesis.  It is the analogue of the elementary part of the Kac--Shapovalov theory of
Verma modules.

\subsection{Graded orthogonals}

\begin{theorem}[Orthogonals inside a weight space]\label{ch18:thm:orthogonals}
For every $N\in\Z$, inside the finite-dimensional space $H^N$ with the nondegenerate form $B^N$:
\begin{equation*}
  (\chXVIIIS^N)^{\perp}=\chXVIIIP^N,\quad (\chXVIIIP^N)^{\perp}=\chXVIIIS^N,\quad
  (\chXVIIIS_K^N)^{\perp}=\chXVIIIK^N,\quad (\chXVIIIK^N)^{\perp}=\chXVIIIS_K^N,
\end{equation*}
and consequently $\dim\chXVIIIS^N+\dim\chXVIIIP^N=\dim H^N=\dim\chXVIIIS_K^N+\dim\chXVIIIK^N$.
\end{theorem}
\begin{leancode}
theorem orthogonal_spurious_eq (N : ℤ) :
    (M.stringWeightForm N).orthogonal (M.weightSlice M.spurious N) =
      M.weightSlice M.physical N := by
theorem orthogonal_kSpurious_eq (N : ℤ) :
    (M.stringWeightForm N).orthogonal (M.weightSlice M.kSpurious N) =
      M.weightSlice M.kVacuum N := by
theorem orthogonal_physical_eq (N : ℤ) :
    (M.stringWeightForm N).orthogonal (M.weightSlice M.physical N) =
      M.weightSlice M.spurious N := by
theorem orthogonal_kVacuum_eq (N : ℤ) :
    (M.stringWeightForm N).orthogonal (M.weightSlice M.kVacuum N) =
      M.weightSlice M.kSpurious N := by
theorem physical_spurious_finrank (N : ℤ) :
    Module.finrank ℂ (M.weightSlice M.spurious N) +
      Module.finrank ℂ (M.weightSlice M.physical N) = Module.finrank ℂ (M.wt N) := by
\end{leancode}
(\leanpath{VOA/Moonshine/MoonshineCNStringOrthogonals.lean}; also \lean{kVacuum_kSpurious_finrank}.)

\begin{dictionary}
\lean{(M.stringWeightForm N).orthogonal W} is Mathlib's orthogonal $\{x\in H^N: B(y,x)=0\ \forall
y\in W\}$ of a subspace $W\subseteq H^N$ (for a symmetric form the side does not matter).
\lean{Module.finrank ℂ} is the dimension.  No hypothesis beyond the string module $M$.
\end{dictionary}
\begin{physics}
The physical states of weight $N$ are exactly the states orthogonal to all spurious states of the same
weight: being annihilated by $L_n$ is, by $L_n^\dagger=L_{-n}$, the same as being orthogonal to the image of
$L_{-n}$.  This is why null states decouple from physical amplitudes.
\end{physics}

\paragraph{Proof on paper.}  The general statement (\lean{positiveKernel_weight_orthogonal}) is: if
operators $A_n$ ($n\ge1$) lower the weight by $n$, operators $D_n$ raise it by $n$, and $B(A_nx,y)=B(x,D_ny)$,
then at weight $N$ the common kernel of the $A_n$ is the orthogonal of $\sum_n D_nH$.  Indeed, for
$x\in H^N$, $A_nx\in H^{N-n}$ is zero iff it is orthogonal to all of $H^{N-n}$ (nondegeneracy there), iff
$B(x,D_ny)=0$ for all $y\in H^{N-n}$; and $D_n$ applied to the other weight spaces lands outside $H^N$, where
$x$ is orthogonal anyway by \lean{B_orth}.  The two remaining equalities are the double-orthogonal
theorem for a nondegenerate form on a finite-dimensional space, and the dimension formula is
$\dim W+\dim W^\perp=\dim H^N$. \qed

\begin{status}
No hypotheses.  \texttt{\#print axioms}: standard axioms (printed in \leanpath{VOA/Moonshine/MoonshineCNAudit4.lean}).
\end{status}

\subsection{The two radicals}

The \emph{radical} of the form restricted to a subspace $U$ at weight $N$ is
$U^N\cap(U^N)^\perp$, in Lean \lean{formRadicalIn}:
\begin{leancode}
def formRadicalIn (U : Submodule ℂ H) (N : ℤ) : Submodule ℂ (M.wt N) :=
  M.weightSlice U N ⊓ (M.stringWeightForm N).orthogonal (M.weightSlice U N)
\end{leancode}
(\leanpath{VOA/Moonshine/MoonshineCNStringPhysical.lean}.)

\begin{theorem}[Radicals of the physical and longitudinal spaces]\label{ch18:thm:radicals}
In every string module, at every central charge: the radical of $B$ on $\chXVIIIP^1$ is $\chXVIIIP^1\cap\chXVIIIS$, and for every
$N$ the radical of $B$ on $\chXVIIIK^N$ is $\chXVIIIK^N\cap\chXVIIIS_K$.
\end{theorem}
\begin{leancode}
theorem physicalOne_radical_eq :
    M.formRadicalIn M.physical 1 = M.weightSlice (M.physical ⊓ M.spurious) 1 := by
theorem kVacuum_radical_eq (N : ℤ) :
    M.formRadicalIn M.kVacuum N = M.weightSlice (M.kVacuum ⊓ M.kSpurious) N := by
\end{leancode}
(\leanpath{VOA/Moonshine/MoonshineCNStringPhysical.lean}.)
\begin{dictionary}
Both are equalities of subspaces of $H^1$, resp.\ $H^N$.  The first is \emph{weight one only} and
needs neither $c=26$ nor (F1), as the docstring says: ``without a critical-charge or light-cone
assumption''.  (In Borcherds' sketch the identification of the radical with the null states comes only after the
splitting; here it is proved first, at every central charge.)
\end{dictionary}
\begin{physics}
``Null physical states are exactly the physical states of zero norm against all physical states.''
\end{physics}
\paragraph{Proof on paper.}  By Theorem~\ref{ch18:thm:orthogonals}, $(\chXVIIIP^1)^\perp=\chXVIIIS^1$, so the radical
$\chXVIIIP^1\cap(\chXVIIIP^1)^\perp$ is $\chXVIIIP^1\cap\chXVIIIS^1$.  Similarly $(\chXVIIIK^N)^\perp=\chXVIIIS_K^N$.  (The only subtlety: $\chXVIIIS\cap H^N=\chXVIIIS^N$
uses that the spurious space is graded, \S\ref{ch18:sec:grading} below.) \qed
\begin{status}
No hypotheses; standard axioms (\leanpath{VOA/Moonshine/MoonshineCNAudit4.lean}).
\end{status}

\subsection{Two generators suffice}

\begin{theorem}[Virasoro generation]\label{ch18:thm:twogenerators}
(a) $\chXVIIIS=L_{-1}H+L_{-2}H$.  (b) If $L_1x=L_2x=0$ then $x\in\chXVIIIP$.  (c) $L_n\chXVIIIK\subseteq\chXVIIIK$ for $n\ge1$
and $K_n\chXVIIIK\subseteq\chXVIIIK$ for all $n$.
\end{theorem}
\begin{leancode}
theorem spurious_eq_sup_two :
    M.spurious = LinearMap.range (M.L (-1)) ⊔ LinearMap.range (M.L (-2)) := by
theorem physical_of_L_one_two (x : H) (h₁ : M.L 1 x = 0)
    (h₂ : M.L 2 x = 0) : x ∈ M.physical :=
theorem kVacuum_L_stable (n : ℕ) (x : H) (hx : x ∈ M.kVacuum) :
    M.L (n + 1 : ℤ) x ∈ M.kVacuum := by
theorem kVacuum_K_stable (n : ℤ) (x : H) (hx : x ∈ M.kVacuum) :
    M.K n x ∈ M.kVacuum := by
\end{leancode}
(\leanpath{VOA/Moonshine/MoonshineCNStringGenerators.lean}.)
\begin{dictionary}
\lean{⊔} is the sum of subspaces.  No hypotheses beyond $M$; valid at every central charge.
\end{dictionary}
\begin{physics}
The Virasoro algebra's negative part is generated by $L_{-1}$ and $L_{-2}$, its positive part by $L_1$
and $L_2$: to check that a state is primary it suffices to check the two lowest modes.
\end{physics}
\paragraph{Proof on paper.}  (a) $[L_{-1},L_{-(n-1)}]=(n-2)L_{-n}$ for $n\ge3$, so $L_{-n}H\subseteq L_{-1}H+L_{-(n-1)}H$;
induct.  (b) $[L_{n-1},L_1]=(n-2)L_n$; induct.  (c) For $k\in\chXVIIIK$ and $m,n\ge1$:
$K_mL_nk=L_nK_mk-[L_n,K_m]k=0+mK_{m+n}k=0$; and the $K$'s commute. \qed
\begin{status}No hypotheses; standard axioms.\end{status}

\subsection{Homogeneity}\label{ch18:sec:grading}

All four derived spaces are \emph{graded}: if a vector lies in $\chXVIIIS$ (resp.\ $\chXVIIIS_K$, $\chXVIIIP$, $\chXVIIIK$), so does each of
its homogeneous components (\lean{spurious_weightComponent}, \lean{kSpurious_weightComponent},
\lean{physical_weightComponent}, \lean{kVacuum_weightComponent}, in
\leanpath{VOA/Moonshine/MoonshineCNStringGrading.lean}, via the projection \lean{weightComponent N} onto
$H^N$ built from \lean{wt_isInternal}).  Consequently the global splitting (F1) restricts to every
weight space:
\begin{leancode}
theorem lightConeSplit_weight (hF : M.LightConeSplit) (N : ℤ) :
    IsCompl (M.weightSlice M.spurious N) (M.weightSlice M.kVacuum N) := by
\end{leancode}
(\leanpath{VOA/Moonshine/MoonshineCNStringGrading.lean}.)  This is the precise sense in which an
\emph{ungraded} hypothesis $H=\chXVIIIS\oplus\chXVIIIK$ is enough: it does not have to be stated weight by weight.  The
file also records that a spurious vector of weight $N$ is $L_{-1}a+L_{-2}b$ with $a\in H^{N-1}$,
$b\in H^{N-2}$ (\lean{spurious_weight_two}).
\begin{physics}
The light-cone decomposition respects the energy: one can decompose level by level.
\end{physics}

\subsection{Straightening of Virasoro words}

The light-cone splitting for the concrete sector (proved in Chapter~\ref{ch:17}, used in
\S\ref{ch18:sec:sector}) needs the classical fact that every word in the negative Virasoro modes can be
rewritten as a combination of \emph{sorted} words $L_{-m_1}\cdots L_{-m_t}$ with $m_1\ge\dots\ge m_t$, at
the same total energy, using only the commutators among negative modes.  The library proves it for an
arbitrary family of operators:

\begin{theorem}[Straightening]\label{ch18:thm:straightening}
Let $D_i$ ($i\in\N$; $D_i$ stands for $L_{-(i+1)}$) be operators on $H$ with
$[D_i,D_j]=(j-i)D_{i+j+1}$.  For every $x\in H$ and every word $w=(i_1,\dots,i_s)$, the vector
$D_{i_1}\cdots D_{i_s}x$ lies in the span of the sorted words applied to $x$ of the same energy
$\sum(i_k+1)$; if $w$ is nonempty, only nonempty sorted words occur.
\end{theorem}
\begin{leancode}
theorem stringNegativeWord_mem_sortedSpan (D : ℕ → Module.End ℂ H)
    (hcomm : ∀ i j, D i * D j - D j * D i = ((j : ℂ) - i) • D (i + j + 1))
    (x : H) (w : List ℕ) :
    stringNegativeWord D x w ∈ stringSortedSpan D x (stringWordEnergy w) := by
theorem stringNegativeWord_mem_nonempty_sortedSpan (D : ℕ → Module.End ℂ H)
    (hcomm : ∀ i j, D i * D j - D j * D i = ((j : ℂ) - i) • D (i + j + 1))
    (x : H) (w : List ℕ) (hw : w ≠ []) :
    stringNegativeWord D x w ∈ Submodule.span ℂ {y | ∃ v : List ℕ,
      v ≠ [] ∧ stringSortedWord v ∧ stringWordEnergy v = stringWordEnergy w ∧
        stringNegativeWord D x v = y} := by
\end{leancode}
(\leanpath{VOA/Moonshine/MoonshineCNStringStraightening.lean}, where
\lean{stringNegativeWord D x [i₁, …, iₛ] = D i₁ (… (D iₛ x))}, \lean{stringWordEnergy w} is
$\sum(i_k+1)$, \lean{stringSortedWord w} is ``weakly decreasing'', and \lean{stringSortedSpan D x N} is
the span of sorted words of energy $N$ applied to $x$.)
\begin{dictionary}
The index shift is $D_i=L_{-(i+1)}$.  With $a=i+1$, $b=j+1$ the Virasoro relation gives
$[L_{-a},L_{-b}]=(-a+b)L_{-a-b}=(j-i)L_{-(i+j+2)}=(j-i)D_{i+j+1}$ (no central term, since $-a-b\neq0$),
which is exactly \lean{hcomm}.  The hypothesis is only the negative-negative commutator, so the theorem applies to any
Virasoro representation and any starting vector (no highest-weight condition).  It is a structural port
of the repository's Verma-module lemma \lean{negWord_mem_sortedSpan}.
\end{dictionary}
\begin{physics}
Every Virasoro descendant can be written in the standard PBW basis $L_{-m_1}\cdots L_{-m_t}|\chi\rangle$,
$m_1\ge\dots\ge m_t$: the familiar basis of a Verma module, here without assuming one has a Verma module.
\end{physics}
\paragraph{Proof on paper.}  Induction on the energy, inserting one letter at a time into a sorted word:
if $D_i$ is to be placed in front of a sorted word whose first letter is $j>i$, use
$D_iD_j=D_jD_i+(j-i)D_{i+j+1}$; the first term is handled by induction on the word length, the second has
the same energy and a shorter word. \qed
\begin{status}No hypotheses beyond \lean{hcomm}; standard axioms (\leanpath{VOA/Moonshine/MoonshineCNAudit5.lean} prints
\lean{stringNegativeWord_mem_sortedSpan}).\end{status}

\section{The critical dimension: the \texorpdfstring{$c=26$}{c = 26} lemma}\label{ch18:sec:critical}

This is the only place in the whole construction where the value of the central charge matters.

\begin{definition}[Modified modes]\label{ch18:def:tilde}
$\widetilde L_{-2}:=L_{-2}+\tfrac32L_{-1}^2$ and $\widetilde L_2:=L_2+\tfrac32L_1^2$.  They are adjoint:
$B(\widetilde L_2x,y)=B(x,\widetilde L_{-2}y)$.
\end{definition}
\begin{leancode}
noncomputable def tildeNegativeTwo : Module.End ℂ H := M.L (-2) + (3 / 2 : ℂ) • (M.L (-1) * M.L (-1))
noncomputable def tildePositiveTwo : Module.End ℂ H := M.L 2 + (3 / 2 : ℂ) • (M.L 1 * M.L 1)
theorem tildeTwo_adjoint (x y : H) :
    M.B (M.tildePositiveTwo x) y = M.B x (M.tildeNegativeTwo y) := by
\end{leancode}
(\leanpath{VOA/Moonshine/MoonshineCNStringCritical.lean}.)
\begin{dictionary}
\lean{M.L (-1) * M.L (-1)} is $L_{-1}^2$ (composition).  Literal transcription.
\end{dictionary}
\begin{physics}
$\widetilde L_{-2}\chi$, for $\chi$ of weight $-1$ annihilated by $L_1$ and $L_2$, is the famous null
vector of the Verma module $M(c,h=-1)$ at level $2$ when $c=26$; the combination $L_{-2}+\frac32L_{-1}^2$ is
the level-two singular vector for $h=-1$ in Kac's table (the list of pairs $(c,h)$ for which the Virasoro Verma module $M(c,h)$ has a singular vector).
\end{physics}

\begin{theorem}[The four identities and the defect $c/2-13$]\label{ch18:thm:critical}
For $x\in H^0$ and $y\in H^{-1}$, at every central charge $c$:
\begin{align*}
 &\text{(i) } L_1L_{-1}x=L_{-1}L_1x, &&\text{(ii) } \widetilde L_2L_{-1}x=L_{-1}\widetilde L_2x,\\
 &\text{(iii) } L_1\widetilde L_{-2}y=\widetilde L_{-2}L_1y, &&\text{(iv) } \widetilde L_2\widetilde L_{-2}y=\widetilde L_{-2}\widetilde L_2y+\bigl(\tfrac c2-13\bigr)y.
\end{align*}
Consequently, if $c=26$, then for every spurious $s$ of weight one, $L_1s\in\chXVIIIS$ (at every $c$) and
$\widetilde L_2s\in\chXVIIIS$.
\end{theorem}
\begin{leancode}
theorem L_one_negative_one_weight_zero (x : H) (hx : x ∈ M.wt 0) :
    M.L 1 (M.L (-1) x) = M.L (-1) (M.L 1 x) := by
theorem tildePositiveTwo_negative_one_weight_zero (x : H) (hx : x ∈ M.wt 0) :
    M.tildePositiveTwo (M.L (-1) x) = M.L (-1) (M.tildePositiveTwo x) := by
theorem L_one_tildeNegativeTwo_weight_neg_one (x : H) (hx : x ∈ M.wt (-1)) :
    M.L 1 (M.tildeNegativeTwo x) = M.tildeNegativeTwo (M.L 1 x) := by
theorem tildeTwo_commutator_weight_neg_one (x : H) (hx : x ∈ M.wt (-1)) :
    M.tildePositiveTwo (M.tildeNegativeTwo x) =
      M.tildeNegativeTwo (M.tildePositiveTwo x) + (M.c / 2 - 13) • x := by
theorem tildeTwo_critical_commute (hc : M.c = 26) (x : H) (hx : x ∈ M.wt (-1)) :
    M.tildePositiveTwo (M.tildeNegativeTwo x) =
      M.tildeNegativeTwo (M.tildePositiveTwo x) := by
theorem criticalSpurious_L_one (x : H) (hx : x ∈ M.spurious) (hw : x ∈ M.wt 1) :
    M.L 1 x ∈ M.spurious := by
theorem criticalSpurious_L_two_tilde (hc : M.c = 26) (x : H)
    (hx : x ∈ M.spurious) (hw : x ∈ M.wt 1) : M.tildePositiveTwo x ∈ M.spurious := by
\end{leancode}
(\leanpath{VOA/Moonshine/MoonshineCNStringCritical.lean} for the identities,
\leanpath{VOA/Moonshine/MoonshineCNStringPhysical.lean} for the two \lean{criticalSpurious} statements.)

\begin{dictionary}
\textbf{Hypotheses.}  Identities (i)--(iv) carry only a weight hypothesis (\lean{x ∈ M.wt 0} or
\lean{M.wt (-1)}); no highest-weight condition is imposed on $y$ (the file's docstring: ``No
highest-weight condition is imposed on the weight-minus-one vector'').  Only \lean{hc : M.c = 26}
enters \lean{tildeTwo_critical_commute} and \lean{criticalSpurious_L_two_tilde}.

\textbf{The defect is explicit} in \lean{tildeTwo_commutator_weight_neg_one}.  The explicit defect is what makes the role of $26$
checkable: the identity holds as an exact polynomial identity in $c$.

\textbf{How it could have been wrong.}  A sign or factor error in $\frac32$ would shift the critical
value away from $26$.  The statement in Lean is an exact identity with the defect $\frac c2-13$, so no such error is possible.
(As \emph{evidence only}: an exact-arithmetic computation in small Verma-module models, done before the formalization,
found the conclusion (A5) below true at $c=26$ and false at $c=25$ and $c=27$.)
\end{dictionary}
\begin{physics}
This is the critical dimension of the bosonic string.  The norm of the spurious state
$\widetilde L_{-2}\chi$ built on a weight-$(-1)$ state $\chi$ is proportional to $c/2-13$: at $c=26$ it is
simultaneously spurious \emph{and} physical, i.e.\ null, so it decouples.  (The further facts that for $c<26$ these states
have nonzero norm and for $c>26$ some physical states have negative norm are the classical results of Goddard--Thorn and
Brower; they are quoted here as physics and are not proved in this chapter.)
\end{physics}

\paragraph{Proof on paper.}  Only the Virasoro relations are used, together with $L_0=0$ on $H^0$,
$-1$ on $H^{-1}$, $-2$ on $H^{-2}$.  (i): $L_1L_{-1}x=L_{-1}L_1x+2L_0x$ and $L_0x=0$.  (ii):
$L_2L_{-1}x=L_{-1}L_2x+3L_1x$ and $L_1^2L_{-1}x=L_{-1}L_1^2x-2L_1x$ (two applications of the $[L_1,L_{-1}]$
relation on weights $0$ and $-1$), so the $3L_1x$ terms cancel in
$\widetilde L_2L_{-1}x=L_{-1}L_2x+3L_1x+\frac32(L_{-1}L_1^2x-2L_1x)$.  (iii): $[L_1,L_{-2}]=3L_{-1}$ and
$[L_1,L_{-1}^2]=2L_{-1}(2L_0+1)$, which is $-2L_{-1}$ on $H^{-1}$; so $[L_1,\widetilde L_{-2}]y=3L_{-1}y-3L_{-1}y=0$.
(iv): on $y\in H^{-1}$,
\begin{align*}
  [L_2,L_{-2}]y&=(4L_0+\tfrac c2)y=(\tfrac c2-4)y, &
  [L_2,L_{-1}^2]y&=6L_{-1}L_1y-6y,\\
  [L_1^2,L_{-2}]y&=6L_{-1}L_1y-6y, &
  [L_1^2,L_{-1}^2]y&=-8L_{-1}L_1y+4y,
\end{align*}
hence $[\widetilde L_2,\widetilde L_{-2}]y=(\frac c2-4)y+\frac32\cdot2(6L_{-1}L_1y-6y)+\frac94(-8L_{-1}L_1y+4y)
=(\frac c2-13)y$: the $L_{-1}L_1y$ terms cancel ($\frac32\cdot2\cdot6=18$ against $\frac94\cdot(-8)=-18$) and the scalars give
$\frac c2-4-18+9=\frac c2-13$.
For the consequence: by Theorem~\ref{ch18:thm:twogenerators}(a) and \lean{spurious_weight_two}, every spurious $s$ of weight
one is $L_{-1}x+L_{-2}y'$ with $x\in H^0$, $y'\in H^{-1}$, and $L_{-2}y'=\widetilde L_{-2}y'-\frac32L_{-1}(L_{-1}y')$ with $L_{-1}y'\in H^0$;
so $s=L_{-1}x+\widetilde L_{-2}y$ with $x\in H^0$, $y\in H^{-1}$.  Now apply the identities term by term:
\begin{center}
\begin{tabular}{lll}
\toprule
 & on $L_{-1}x$ ($x\in H^0$) & on $\widetilde L_{-2}y$ ($y\in H^{-1}$)\\ \midrule
$L_1$ & (i): $=L_{-1}(L_1x)\in\chXVIIIS$ & (iii): $=\widetilde L_{-2}(L_1y)\in\chXVIIIS$\\
$\widetilde L_2$ & (ii): $=L_{-1}(\widetilde L_2x)\in\chXVIIIS$ & (iv), $c=26$: $=\widetilde L_{-2}(\widetilde L_2y)\in\chXVIIIS$\\
\bottomrule
\end{tabular}
\end{center}
Every entry is $L_{-1}$ or $\widetilde L_{-2}$ applied to some vector (of weight $-1$, $-2$ or $-3$), hence spurious, since
$\widetilde L_{-2}z=L_{-2}z+\frac32L_{-1}(L_{-1}z)\in\chXVIIIS$ for any $z$.  Only the bottom-right entry needs $c=26$. \qed

\begin{status}
No hypotheses beyond the displayed weight hypotheses and, for the critical statements, $c=26$.  Standard
axioms (\leanpath{VOA/Moonshine/MoonshineCNAudit4.lean} prints \lean{tildeTwo_commutator_weight_neg_one},
\lean{tildeTwo_critical_commute}, \lean{criticalSpurious_L_two_tilde}).  The supporting identities
\lean{L_two_negative_two}, \lean{L_one_square_negative_square_weight_neg_one}, etc.\ in the same file are
the four bracket computations displayed above.
\end{status}

\section{The abstract no-ghost theorem}\label{ch18:sec:abstractnoghost}

From now on in this part we assume the light-cone splitting (F1) $H=\chXVIIIS\oplus\chXVIIIK$, as an explicit
hypothesis \lean{hF : M.LightConeSplit}.  It is \emph{proved} for the concrete sectors in
\S\ref{ch18:sec:sector}; here it is a hypothesis.

\subsection{The longitudinal splitting}

\begin{theorem}[Longitudinal spurious states]\label{ch18:thm:longitudinal}
Under (F1): $\chXVIIIS_K\subseteq\chXVIIIS+(\chXVIIIK\cap\chXVIIIS_K)$, and hence $\chXVIIIS_K=(\chXVIIIS_K\cap\chXVIIIS)\oplus(\chXVIIIS_K\cap\chXVIIIK)$, at every weight.
\end{theorem}
\begin{leancode}
theorem kSpurious_le_spurious_sup (hF : M.LightConeSplit) :
    M.kSpurious ≤ M.spurious ⊔ (M.kVacuum ⊓ M.kSpurious) := by
theorem kSpurious_eq_intersections_sup (hF : M.LightConeSplit) :
    M.kSpurious = (M.kSpurious ⊓ M.spurious) ⊔ (M.kSpurious ⊓ M.kVacuum) := by
\end{leancode}
(\leanpath{VOA/Moonshine/MoonshineCNStringLongitudinal.lean}; the weightwise, direct form is
\lean{kSpurious_weight_split} in \leanpath{VOA/Moonshine/MoonshineCNStringTransverse.lean}.)
\begin{dictionary}
Only \lean{hF}.  The lower bound \lean{wt_bdd} of the string module is used, as the docstring says, ``in
the strong induction'' — it is part of the structure, not an extra hypothesis.
\end{dictionary}
\begin{physics}
A state created by a longitudinal oscillator $\alpha^+_{-m}$ is, up to spurious states, a longitudinal
vacuum state: in light-cone gauge, $\alpha^+$ excitations are pure gauge.
\end{physics}
\paragraph{Proof on paper.}  Strong induction on the weight (legitimate because $H^N=0$ below $N_0$).  For
$x\in H^N$ and $m\ge1$, write $x=s+k$ by (F1).  Then $K_{-m}k\in\chXVIIIK\cap\chXVIIIS_K$ ($K$'s commute, so $K_{-m}$ preserves
$\chXVIIIK$).  For $s=\sum_nL_{-n}x_n$, the commutator gives $K_{-m}L_{-n}x_n=L_{-n}K_{-m}x_n-mK_{-m-n}x_n$
(\lean{negative_K_negative_L}); the first term is spurious, the second is handled by induction since $x_n$
has smaller weight. \qed
\begin{status}Only \lean{hF}; standard axioms.\end{status}

\subsection{The decompositions in every weight}

\begin{theorem}[Abstract no-ghost theorem, parts (A1)--(A4)]\label{ch18:thm:A14}
Under (F1), for every $N$:
\begin{enumerate}[label=(A\arabic*)]
\item $\dim\chXVIIIS^N=\dim\chXVIIIS_K^N$ and $H^N=\chXVIIIP^N\oplus\chXVIIIS_K^N$; globally $H=\chXVIIIP\oplus\chXVIIIS_K$;
\item $\chXVIIIS_K^N=(\chXVIIIS_K^N\cap\chXVIIIS^N)\oplus(\chXVIIIS_K^N\cap\chXVIIIK^N)$;
\item $\chXVIIIK^N=\chXVIIIT^N\oplus\mathrm{rad}(B|_{\chXVIIIK^N})$, where $\mathrm{rad}(B|_{\chXVIIIK^N})=\chXVIIIK^N\cap\chXVIIIS_K$;
\item $B$ is nondegenerate on $\chXVIIIT^N$.
\end{enumerate}
\end{theorem}
\begin{leancode}
theorem spurious_kSpurious_finrank (hF : M.LightConeSplit) (N : ℤ) :
    Module.finrank ℂ (M.weightSlice M.spurious N) =
      Module.finrank ℂ (M.weightSlice M.kSpurious N) := by
theorem physical_kSpurious_isCompl (hF : M.LightConeSplit) :
    IsCompl M.physical M.kSpurious := by
theorem transverse_radical_disjoint (hF : M.LightConeSplit) (N : ℤ) :
    Disjoint (M.weightSlice M.transverse N) (M.formRadicalIn M.kVacuum N) := by
theorem kVacuum_eq_transverse_sup_radical (hF : M.LightConeSplit) (N : ℤ) :
    M.weightSlice M.kVacuum N =
      M.weightSlice M.transverse N ⊔ M.formRadicalIn M.kVacuum N := by
theorem transverse_nondegenerate (hF : M.LightConeSplit) (N : ℤ) (x : H)
    (hx : x ∈ M.transverse) (hw : x ∈ M.wt N)
    (hB : ∀ y ∈ M.transverse, y ∈ M.wt N → M.B x y = 0) : x = 0 := by
\end{leancode}
(\leanpath{VOA/Moonshine/MoonshineCNStringTransverse.lean}; also \lean{physical_kSpurious_isCompl_weight}
and \lean{kSpurious_weight_split}.)
\begin{dictionary}
\lean{Disjoint A B} means $A\cap B=0$.  (A3) is stated as two statements (sum, \lean{kVacuum_eq_transverse_sup_radical};
directness, \lean{transverse_radical_disjoint}) together with \lean{kVacuum_radical_eq} of
Theorem~\ref{ch18:thm:radicals}.  (A4) is stated elementwise: a transverse $x$ of weight $N$ orthogonal to
all transverse $y$ of weight $N$ is $0$.  Only \lean{hF}; no central-charge hypothesis.
\end{dictionary}
\begin{physics}
At every level, the longitudinal vacua (no $\alpha^+$ excitations) are the transverse states plus null
states, and the transverse states carry a nondegenerate inner product.  This is the ``light-cone gauge
is equivalent to covariant gauge'' statement at the level of Fock spaces, before imposing $L_0=1$.
\end{physics}
\paragraph{Proof on paper.}  (A1): by (F1) and Theorem~\ref{ch18:thm:orthogonals},
$\dim\chXVIIIS^N=\dim H^N-\dim\chXVIIIK^N=\dim\chXVIIIS_K^N$.  If $x\in\chXVIIIP^N\cap\chXVIIIS_K^N$ then $x\perp\chXVIIIS^N$ (as $x\in\chXVIIIP$) and $x\perp\chXVIIIK^N$ (as
$x\in\chXVIIIS_K=(\chXVIIIK)^\perp$), so $x\perp H^N=\chXVIIIS^N+\chXVIIIK^N$ and $x=0$; dimensions add up.  (A2): Theorem~\ref{ch18:thm:longitudinal} and the
modular law, directness from $\chXVIIIS\cap\chXVIIIK=0$.  (A3): $\chXVIIIT^N$ is the kernel of $\beta:k\mapsto B(k,\cdot)|_{\chXVIIIS^N}$ on
$\chXVIIIK^N$ (a $k\in\chXVIIIK^N$ is physical iff it is orthogonal to $\chXVIIIS^N$, Theorem~\ref{ch18:thm:orthogonals}).  The image of $\beta$
is the annihilator in $(\chXVIIIS^N)^*$ of the right kernel $\{s\in\chXVIIIS^N:B(\chXVIIIK^N,s)=0\}=\chXVIIIS^N\cap(\chXVIIIK^N)^\perp=\chXVIIIS^N\cap\chXVIIIS_K^N$
(Theorem~\ref{ch18:thm:orthogonals} again), so it has dimension $\dim\chXVIIIS^N-\dim(\chXVIIIS^N\cap\chXVIIIS_K^N)$, and
$\dim\chXVIIIT^N=\dim\chXVIIIK^N-\dim\chXVIIIS_K^N+\dim(\chXVIIIS^N\cap\chXVIIIS_K^N)=\dim\chXVIIIK^N-\dim(\chXVIIIK^N\cap\chXVIIIS_K^N)$ by (A1), (A2); and $\chXVIIIT\cap\chXVIIIS_K\subseteq\chXVIIIP\cap\chXVIIIS_K=0$.
(A4): a $t\in\chXVIIIT^N$ orthogonal to $\chXVIIIT^N$ is also orthogonal to the radical of $\chXVIIIK^N$, hence to $\chXVIIIK^N$ by (A3),
hence lies in that radical, hence is $0$. \qed
\begin{status}Only \lean{hF}; standard axioms (\leanpath{VOA/Moonshine/MoonshineCNAudit4.lean}).\end{status}

\subsection{The weight-one physical splitting (A5)}

\begin{theorem}[Abstract no-ghost theorem, part (A5)]\label{ch18:thm:A5}
If $c=26$ and (F1) holds, then $\chXVIIIP^1=\chXVIIIT^1\oplus(\chXVIIIP^1\cap\chXVIIIS)$.
\end{theorem}
\begin{leancode}
theorem physicalOne_eq_transverse_sup_null (hc : M.c = 26) (hF : M.LightConeSplit) :
    M.weightSlice M.physical 1 = M.weightSlice M.transverse 1 ⊔
      M.weightSlice (M.physical ⊓ M.spurious) 1 := by
theorem physicalOne_transverse_null_disjoint (hF : M.LightConeSplit) :
    Disjoint (M.weightSlice M.transverse 1) (M.weightSlice (M.physical ⊓ M.spurious) 1) := by
\end{leancode}
(\leanpath{VOA/Moonshine/MoonshineCNStringPhysical.lean}.)
\begin{dictionary}
The spanning half needs $c=26$ and (F1); the directness half needs only (F1).  Together with
Theorem~\ref{ch18:thm:radicals} ($\mathrm{rad}(B|_{\chXVIIIP^1})=\chXVIIIP^1\cap\chXVIIIS$) this is the full weight-one statement: $\chXVIIIP^1$ is the
transverse space plus the radical.
\end{dictionary}
\begin{physics}
The no-ghost theorem proper: every physical state of a $26$-dimensional string is a transverse
(light-cone) state plus a null state.
\end{physics}
\paragraph{Proof on paper.}  Let $p\in\chXVIIIP^1$ and write $p=k+s$ with $k\in\chXVIIIK^1$, $s\in\chXVIIIS^1$ (F1).  For
$A\in\{L_1,\widetilde L_2\}$: $Ap=0$, $Ak\in\chXVIIIK$ (Theorem~\ref{ch18:thm:twogenerators}(c)) and $As\in\chXVIIIS$
(Theorem~\ref{ch18:thm:critical}, $c=26$); $Ak$ and $As$ have the same weight and sum to zero, and
$\chXVIIIS\cap\chXVIIIK=0$, so $Ak=As=0$.  From $L_1k=0=\widetilde L_2k$ we get $L_2k=0$, so $k\in\chXVIIIP$ by
Theorem~\ref{ch18:thm:twogenerators}(b), i.e.\ $k\in\chXVIIIT^1$, and $s=p-k\in\chXVIIIP^1\cap\chXVIIIS$.  Directness: $\chXVIIIT\cap\chXVIIIS\subseteq\chXVIIIK\cap\chXVIIIS=0$. \qed
\begin{status}Hypotheses \lean{hc}, \lean{hF}; standard axioms (\leanpath{VOA/Moonshine/MoonshineCNAudit4.lean}).\end{status}

\subsection{The quotient isometries}

The quotient of a weight slice by the radical of its form, with the descended form, is defined
generally:
\begin{leancode}
def weightRadical (U : Submodule ℂ H) (N : ℤ) :
    Submodule ℂ (M.weightSlice U N) :=
  (M.formRadicalIn U N).comap (M.weightSlice U N).subtype
abbrev weightRadicalQuotient (U : Submodule ℂ H) (N : ℤ) :=
  M.weightSlice U N ⧸ M.weightRadical U N
def radicalQuotientForm (U : Submodule ℂ H) (N : ℤ) :
    LinearMap.BilinForm ℂ (M.weightRadicalQuotient U N) :=
\end{leancode}
(\leanpath{VOA/Moonshine/MoonshineCNStringQuotient.lean}; \lean{radicalQuotientForm} is characterised by
\lean{radicalQuotientForm_mk}: its value on classes $[x],[y]$ is $B(x,y)$, and it is symmetric and
nondegenerate, \lean{radicalQuotientForm_symm}, \lean{radicalQuotientForm_nondegenerate}.)

\begin{theorem}[The no-ghost isometries]\label{ch18:thm:quotientequivs}
Under (F1), the inclusion $\chXVIIIT^N\hookrightarrow\chXVIIIK^N$ induces an isometric isomorphism
$\chXVIIIT^N\cong\chXVIIIK^N/\mathrm{rad}$; if moreover $c=26$, the inclusion $\chXVIIIT^1\hookrightarrow\chXVIIIP^1$ induces an isometric
isomorphism $\chXVIIIT^1\cong\chXVIIIP^1/\mathrm{rad}$, whose inverse sends $[p]$ to the $\chXVIIIK$-component of $p$ in
$H=\chXVIIIS\oplus\chXVIIIK$.  Composing, $\chXVIIIP^1/\mathrm{rad}\cong\chXVIIIK^1/\mathrm{rad}$ isometrically.  All three isomorphisms are equivariant for
every symmetry $A$ of the string module.
\end{theorem}
\begin{leancode}
def transverseKQuotientEquiv (hF : M.LightConeSplit) (N : ℤ) :
    M.weightSlice M.transverse N ≃ₗ[ℂ] M.weightRadicalQuotient M.kVacuum N :=
def transversePhysicalQuotientEquiv (hc : M.c = 26) (hF : M.LightConeSplit) :
    M.weightSlice M.transverse 1 ≃ₗ[ℂ] M.weightRadicalQuotient M.physical 1 :=
theorem transversePhysicalQuotientEquiv_isometry (hc : M.c = 26) (hF : M.LightConeSplit)
    (x y : M.weightSlice M.transverse 1) :
    M.radicalQuotientForm M.physical 1 (M.transversePhysicalQuotientEquiv hc hF x)
      (M.transversePhysicalQuotientEquiv hc hF y) = M.B x.val y.val := rfl
theorem physicalQuotientRepresentative_kComponent (hc : M.c = 26) (hF : M.LightConeSplit)
    (p : M.weightSlice M.physical 1) :
    (((M.transversePhysicalQuotientEquiv hc hF).symm
      (Submodule.Quotient.mk p)).val : H) =
      M.kVacuum.projection M.spurious hF.symm p.val := by
\end{leancode}
(\leanpath{VOA/Moonshine/MoonshineCNStringQuotient.lean}.)
\begin{leancode}
def physicalKQuotientEquiv (hc : M.c = 26) (hF : M.LightConeSplit) :
    M.weightRadicalQuotient M.physical 1 ≃ₗ[ℂ] M.weightRadicalQuotient M.kVacuum 1 :=
  (M.transversePhysicalQuotientEquiv hc hF).symm.trans (M.transverseKQuotientEquiv hF 1)
theorem physicalKQuotientEquiv_isometry (hc : M.c = 26) (hF : M.LightConeSplit)
    (p q : M.weightRadicalQuotient M.physical 1) :
    M.radicalQuotientForm M.kVacuum 1 (M.physicalKQuotientEquiv hc hF p)
      (M.physicalKQuotientEquiv hc hF q) = M.radicalQuotientForm M.physical 1 p q := by
\end{leancode}
(\leanpath{VOA/Moonshine/MoonshineCNStringComparison.lean}; equivariance:
\lean{transverseKQuotientEquiv_equivariant}, \lean{transversePhysicalQuotientEquiv_equivariant} in
\leanpath{VOA/Moonshine/MoonshineCNStringSymmetry.lean} and
\lean{StringGroupAction.physicalKQuotientEquiv_equivariant} in
\leanpath{VOA/Moonshine/MoonshineCNStringComparison.lean}.)

\begin{dictionary}
\textbf{Decoding.}  \lean{≃ₗ[ℂ]} is a linear isomorphism.  \lean{Submodule.Quotient.mk p} is the class
$[p]$.  \lean{.symm} is the inverse.  \lean{M.kVacuum.projection M.spurious hF.symm} is Mathlib's
projection onto $\chXVIIIK$ along $\chXVIIIS$ for the complementary pair given by (F1).  The isometry statement for
\lean{transversePhysicalQuotientEquiv} is proved by \lean{rfl}: the map is literally $t\mapsto[t]$, so the
descended form on $[t],[t']$ is $B(t,t')$ by definition.  The content is in the \emph{bijectivity}
(\lean{transverseQuotientMap_bijective}), which uses (A3), resp.\ (A5) and Theorem~\ref{ch18:thm:radicals}.

\textbf{Equivariance.}  For a \lean{StringGroupAction} $A$, the file \lean{StringSymmetry} first shows that
$A$ preserves $\chXVIIIS,\chXVIIIS_K,\chXVIIIP,\chXVIIIK,\chXVIIIT$ and the null space, then restricts it to weight slices
(\lean{stableWeightAction}) and descends it to radical quotients (\lean{stableWeightQuotientAction}),
isometrically; the equivariance of the three isomorphisms is then a statement about these actual
quotient actions (two of them hold by \lean{rfl}).

\textbf{Could it be vacuous?}  If (F1) were never satisfiable, every statement here would be empty.  It
is satisfied by the concrete sectors (\S\ref{ch18:sec:sector}).  If the radical were the whole space, the
quotients would be zero; the isometry with $\Vnat_{1-xy}$ proved later (which has the nondegenerate
\lean{natForm}) rules this out whenever $\Vnat_{1-xy}\neq0$.
\end{dictionary}
\begin{physics}
Covariant physical states modulo null states $=$ light-cone states $=$ longitudinal vacua modulo null
states, as inner-product spaces and as representations of every internal symmetry.  The inverse map
``take the $\chXVIIIK$-component'' is the passage to light-cone gauge.
\end{physics}
\begin{status}
Hypotheses \lean{hF} (and \lean{hc} for the weight-one physical statements).  Standard axioms
(\leanpath{VOA/Moonshine/MoonshineCNAudit4.lean}, \leanpath{VOA/Moonshine/MoonshineCNAudit5.lean},
\leanpath{VOA/Moonshine/MoonshineCNAudit8.lean}).
\end{status}

\section{The moonshine string sector \texorpdfstring{$\Vnat\otimes\mathbb C[a_j,b_j]$}{V-natural tensor C[a,b]}}\label{ch18:sec:sector}

We now build, for each momentum $(x,y)\in\chXVIIIII$, an explicit string module whose underlying space is
$\Vnat\otimes\mathbb C[a_j,b_j:j\ge1]$, and \emph{prove} the light-cone splitting for it when
$(x,y)\neq0$.  In \S\ref{ch18:sec:roots} this model is identified with the momentum-$(x,y)$ sector of the
vertex algebra $\Vnat\otimes V_{\chXVIIIII}$ of Chapter~\ref{ch:17}.

\subsection{The hyperbolic Fock space (recalled from Chapter~\ref{ch:17})}

Let $w_0=(0,1)$, $w_1=(1,0)$, an isotropic basis of $\chXVIIIII$: $(w_0,w_0)=(w_1,w_1)=0$, $(w_0,w_1)=1$.  The Fock
space of the rank-two Heisenberg algebra at momentum $r$ is modelled on the polynomial ring
$F=\mathbb C[a_j,b_j:j\ge1]$, where $a_j$ creates $w_0(-j)$ and $b_j$ creates $w_1(-j)$.  In Lean
\begin{leancode}
abbrev hypFock := FockN 2
def hypOscillator (a : Fin 2) (q : ℂ) (n : ℤ) : Module.End ℂ hypFock :=
  if n < 0 then oscN 2 1 a n
  else if 0 < n then oscN 2 1 a.rev n
  else q • LinearMap.id
def hypLevel (N : ℤ) : Submodule ℂ hypFock := gradingZN 2 N
\end{leancode}
(\leanpath{VOA/Moonshine/MoonshineCNHyperbolicFock.lean}.)  Here \lean{FockN 2 = MvPolynomial (ℕ × Fin 2) ℂ}
and, by the file's docstring, the variable \lean{(j,0)} is $a_{j+1}$ and \lean{(j,1)} is $b_{j+1}$.  The
oscillator of family \lean{a} acts for $n<0$ by multiplication by the variable \lean{(-n-1, a)} and for
$n>0$ by $n\,\partial$ with respect to the variable of the \emph{opposite} family \lean{a.rev} at level
$n$ (\lean{hypOscillator_negative_apply}, \lean{hypOscillator_positive_apply}); its zero mode is the
scalar \lean{q}.  So, with $a(n):=$\lean{hypOscillator 0 κ n} and $b(n):=$\lean{hypOscillator 1 ell n},
\begin{equation*}
  a(-j)=a_j\cdot,\quad a(j)=j\,\partial/\partial b_j,\quad a(0)=\kappa,\qquad
  b(-j)=b_j\cdot,\quad b(j)=j\,\partial/\partial a_j,\quad b(0)=\ell,
\end{equation*}
with $[a(m),b(n)]=m\delta_{m+n,0}$ and $[a(m),a(n)]=[b(m),b(n)]=0$ (\lean{hypOscillator_comm}).
\lean{hypLevel N} is the span of monomials of total level $N$ ($a_j,b_j$ have level $j$); it is
finite-dimensional.

The Sugawara (Virasoro) operators are given by explicit finite formulas, avoiding infinite sums:
\begin{leancode}
def hypRaisingDerivation (n : ℕ) : Derivation ℂ hypFock hypFock :=
  mkDerivation ℂ fun v => (v.1 + 1 : ℂ) • X (v.1 + n, v.2)
def hypNegativePotential (κ ell : ℂ) (j : ℕ) : hypFock :=
  κ • X (j, 1) + ell • X (j, 0) +
    ∑ k ∈ Finset.range j, X (j - 1 - k, 0) * X (k, 1)
def hypSugawara (κ ell : ℂ) (n : ℤ) : Module.End ℂ hypFock :=
  if n < 0 then
    LinearMap.mul ℂ hypFock (hypNegativePotential κ ell (-n - 1).toNat) +
      (hypRaisingDerivation (-n).toNat).toLinearMap
  else if 0 < n then
    (hypLoweringDerivation κ ell n.toNat).toLinearMap +
      hypPositiveSecondOrder (n - 1).toNat
  else (κ * ell) • LinearMap.id + hypEnergyDerivation.toLinearMap
\end{leancode}
(\leanpath{VOA/Moonshine/MoonshineCNHyperbolicSugawara.lean}.)  For $n\ge1$ this is
\begin{equation}\label{ch18:eq:sugawara}
  L^F_{-n}=\kappa\,b_n+\ell\,a_n+\sum_{i=1}^{n-1}a_{n-i}b_i+\sum_{k\ge1}k\Bigl(a_{n+k}\frac{\partial}{\partial a_k}
  +b_{n+k}\frac{\partial}{\partial b_k}\Bigr),\qquad L^F_0=\kappa\ell+\text{level},
\end{equation}
the explicit form of $L^F_n=\sum_k{:}a(n-k)b(k){:}$; the positive modes are a first-order derivation plus a
finite second-order part.  The contravariant form is the monomial pairing
$B_F(a^\alpha b^\beta,a^{\alpha'}b^{\beta'})=\delta_{\alpha\beta'}\delta_{\beta\alpha'}\prod_jj^{\alpha_j}\alpha_j!\prod_jj^{\beta_j}\beta_j!$,
in Lean \lean{hypForm} (\leanpath{VOA/Moonshine/MoonshineCNHyperbolicForm.lean}):
\begin{leancode}
def hypForm : LinearMap.BilinForm ℂ hypFock :=
  Finsupp.lsum ℂ fun d =>
    LinearMap.smulRight LinearMap.id (hypMonomialNorm d • lcoeff ℂ (hypSwapExponent d))
\end{leancode}
where \lean{hypSwapExponent} exchanges the two families and \lean{hypMonomialNorm d} is
$\prod_v(\text{level of }v)^{d_v}\,d_v!$.  Chapter~\ref{ch:17} proves that $L^F$ satisfies the
Virasoro relations with $c=2$, that $L^F_n$ and $L^F_{-n}$ are $B_F$-adjoint, and that $B_F$ is
nondegenerate on each level.

\begin{physics}
$F$ is the Fock space of the two light-cone bosons $X^\pm$ at momentum $r$, written in null coordinates:
$a_j$ creates $\alpha^+_{-j}$ and $b_j$ creates $\alpha^-_{-j}$, and $\kappa$, $\ell$ are the
light-cone momenta $p^\pm$ (here $\kappa=(w_0,r)$, $\ell=(w_1,r)$).  The pairing of $a$ with $b$ (and not with $a$) is the null-coordinate form of
the Minkowski metric.
\end{physics}

\paragraph{The momentum labels.}  In the Lean sector at momentum $(x,y)$ the two charges are
$\kappa=x$, $\ell=y$, i.e.\ \lean{hypSugawara x y}.  The frame file
\leanpath{VOA/Moonshine/MoonshineCNSectorFrame.lean} explains: ``family zero [is] assigned to the
lattice vector \lean{(0,1)}.  Thus the sector charges \lean{(κ,ell) = (x,y)} match momentum \lean{(x,y)}'':
with $w_0=(0,1)$, $w_1=(1,0)$ one has $(w_0,(x,y))=x$ and $(w_1,(x,y))=y$ in the form
\eqref{ch18:eq:repform} (\lean{lorentzianIsotropicVector_charge}).  Hence $\kappa\ell=xy=r^2/2$.

\subsection{The sector operators and form}

\begin{definition}[The sector Virasoro operators and form]\label{ch18:def:sectorops}
For $\kappa,\ell\in\mathbb C$ let $L_n:=L^\natural_n\otimes1+1\otimes L^F_n$ on $\Vnat\otimes F$, where $L^\natural_n$
is the Virasoro action of $\Vnat$, and let $B:=B^\natural\otimes B_F$.  For $a\in\{0,1\}$ the longitudinal operators are
$K_n:=1\otimes(\text{family-}a\text{ oscillator})(n)$, with zero mode the charge of that family.
\end{definition}
\begin{leancode}
def vnatHypMode (κ ell : ℂ) (n : ℤ) : Module.End ℂ (Vnat ⊗[ℂ] hypFock) :=
  tensorTranslation (nearVirasoroMode VnatVOA.toNearVOA n) (hypSugawara κ ell n)
def vnatHypForm : LinearMap.BilinForm ℂ (Vnat ⊗[ℂ] hypFock) := natForm.tmul hypForm
\end{leancode}
(\leanpath{VOA/Moonshine/MoonshineCNVnatSectorOperators.lean}.)
\begin{leancode}
def vnatHypLongitudinal (a : Fin 2) (κ ell : ℂ) (n : ℤ) :
    Module.End ℂ (Vnat ⊗[ℂ] hypFock) :=
  (hypOscillator a (if a = 0 then κ else ell) n).lTensor Vnat
\end{leancode}
(\leanpath{VOA/Moonshine/MoonshineCNVnatSectorLaws.lean}.)
\begin{dictionary}
\lean{tensorTranslation TA TB} is $T_A\otimes1+1\otimes T_B$ (\leanpath{VOA/Core/TensorFieldLift.lean}:
\lean{LinearMap.rTensor B TA + LinearMap.lTensor A TB}); \lean{nearVirasoroMode VnatVOA.toNearVOA n} is
$L^\natural_n=\omega^\natural_{(n+1)}$ (\leanpath{VOA/Moonshine/MoonshineCNTensorConformal.lean}:
\lean{mode (V.Y V.conformalVector) (n + 1)}).  \lean{natForm.tmul hypForm} is the tensor product of
the two bilinear forms, $B(v\otimes p,v'\otimes p')=B^\natural(v,v')B_F(p,p')$.  \lean{.lTensor Vnat} is
$1_{\Vnat}\otimes(\cdot)$.  These are the actual $\Vnat$ Virasoro modes and the actual $\Vnat$ form, not
abstract data.
\end{dictionary}
\begin{physics}
Total stress tensor (internal Monster CFT plus light-cone bosons), the string
inner product, and the light-cone current $\alpha^+(z)$ (or $\alpha^-(z)$).
\end{physics}

\begin{theorem}[Central charge 26, adjointness, nondegeneracy]\label{ch18:thm:sectorlaws}
For all $\kappa,\ell$: the $L_n$ satisfy the Virasoro relations with central charge $26$; $B$ is symmetric,
$L_n$ and $L_{-n}$ are $B$-adjoint, and $B$ is nondegenerate on the whole of $\Vnat\otimes F$.
\end{theorem}
\begin{leancode}
theorem vnatHypMode_virasoro (κ ell : ℂ) (m n : ℤ) :
    vnatHypMode κ ell m * vnatHypMode κ ell n -
      vnatHypMode κ ell n * vnatHypMode κ ell m =
      ((m : ℂ) - n) • vnatHypMode κ ell (m + n) +
        if m + n = 0 then (((m : ℂ) ^ 3 - m) / 12 * 26) • 1 else 0 := by
theorem vnatHypForm_symm (x y : Vnat ⊗[ℂ] hypFock) : vnatHypForm x y = vnatHypForm y x := by
theorem vnatHypForm_adjoint (κ ell : ℂ) (n : ℤ) (x y : Vnat ⊗[ℂ] hypFock) :
    vnatHypForm (vnatHypMode κ ell n x) y = vnatHypForm x (vnatHypMode κ ell (-n) y) := by
theorem vnatHypForm_nondegenerate (x : Vnat ⊗[ℂ] hypFock) (hx : x ≠ 0) :
    ∃ y, vnatHypForm x y ≠ 0 := by
\end{leancode}
(\leanpath{VOA/Moonshine/MoonshineCNVnatSectorOperators.lean}.)
\begin{dictionary}
The central term \lean{(((m : ℂ) ^ 3 - m) / 12 * 26) • 1} is literally $\frac{26}{12}(m^3-m)$: the
number $26$ appears in the statement, with no hypothesis.  It comes from $c(\Vnat)=24$ (Chapter~\ref{ch:14},
\lean{VOA.Moonshine.moonshine}) plus $c(F)=2$ (Chapter~\ref{ch:17}).  Adjointness combines
\lean{natForm_virasoro_adj} with the $F$-adjointness of $L^F$.  Nondegeneracy on the whole space
follows from the coefficient formula \lean{vnatHypForm_coefficient}: pairing with $v\otimes(\text{swapped
monomial})$ extracts, up to the nonzero factor \lean{hypMonomialNorm e}, the $\Vnat$-coefficient of a
monomial, and \lean{natForm} is nondegenerate.
\end{dictionary}
\begin{physics}
$c_{\text{matter}}=24+2=26$: the Monster CFT plus two light-cone bosons is a critical bosonic string.
\end{physics}
\begin{status}No hypotheses; standard axioms (\leanpath{VOA/Moonshine/MoonshineCNAudit7.lean}).\end{status}

\subsection{The weight grading}

\begin{definition}[Sector weight spaces]\label{ch18:def:sectorweight}
$H^N_{(x,y)}:=\operatorname{span}\{v\otimes p: v\in\Vnat_i,\ p\in F_j,\ i+j+xy=N\}$: matter weight plus
oscillator level plus half the momentum norm.
\end{definition}
\begin{leancode}
def vnatSectorWeight (x y N : ℤ) : Submodule ℂ (Vnat ⊗[ℂ] hypFock) :=
  Submodule.span ℂ {t | ∃ (i j : ℤ) (v : Vnat) (p : hypFock),
    v ∈ VnatVOA.grading i ∧ p ∈ hypLevel j ∧ i + j + x * y = N ∧ v ⊗ₜ[ℂ] p = t}
theorem vnatSectorWeight_eq_eigenspace (x y N : ℤ) :
    vnatSectorWeight x y N = (vnatHypMode x y 0).eigenspace (N : ℂ) := by
theorem vnatSectorWeight_isInternal (x y : ℤ) :
    DirectSum.IsInternal (vnatSectorWeight x y) := by
theorem vnatSectorWeight_negative (x y N : ℤ) (hN : N < x * y) :
    vnatSectorWeight x y N = ⊥ := by
theorem vnatSectorWeight_finiteDimensional (x y N : ℤ) :
    FiniteDimensional ℂ (vnatSectorWeight x y N) := by
\end{leancode}
(\leanpath{VOA/Moonshine/MoonshineCNVnatSectorGrading.lean}.)
\begin{dictionary}
The weight space is exactly the $L_0$-eigenspace with eigenvalue $N$ (\lean{eigenspace (N : ℂ)}); this
pins the definition down.  The lower bound is $N\ge xy$ because $\Vnat$ has no negative weights and $F$ no
negative levels.  Finite-dimensionality uses that each $\Vnat_i$ and each $F_j$ is finite-dimensional and
only finitely many $(i,j)$ with $i,j\ge0$ have $i+j=N-xy$.  Note that the momentum $(x,y)$ enters only
through $xy$ and through the charges of the Sugawara operator.
\end{dictionary}
\begin{physics}
$L_0=L_0^\natural+N^{\pm}+\frac12r^2$: the mass-shell condition $L_0=1$ reads
$\text{(internal level)}+\text{(light-cone level)}=1-xy$.
\end{physics}
\begin{status}No hypotheses; standard axioms (\leanpath{VOA/Moonshine/MoonshineCNAudit7.lean}).\end{status}

The remaining string-module laws are in \leanpath{VOA/Moonshine/MoonshineCNVnatSectorLaws.lean}:
\lean{vnatSectorWeight_mode} ($L_n$ lowers weight by $n$), \lean{vnatSectorWeight_oscillator} (same for
$K_n$), \lean{vnatSectorWeight_orthogonal} and \lean{vnatSectorWeight_nondegenerate} (orthogonality of
distinct weights and nondegeneracy on each weight), \lean{vnatHypLongitudinal_comm},
\lean{vnatHypLongitudinal_mixed} ($[L_m,K_n]=-nK_{m+n}$, here for \emph{all} $m,n$, which is stronger
than the structure requires) and \lean{vnatHypLongitudinal_adjoint}.

\subsection{The concrete string module}

\begin{definition}[The moonshine string sector]\label{ch18:def:vnatstringmodule}
For $(x,y)\in\Z^2$ and a family $a\in\{0,1\}$, the string module $H_{(x,y),a}$ has carrier
$\Vnat\otimes F$, weights $H^N_{(x,y)}$, central charge $26$, Virasoro operators $L_n$ at charges $(x,y)$, longitudinal
operators $K_n$ of family $a$, and form $B^\natural\otimes B_F$.  For $(x,y)\neq0$, the \emph{string sector}
$H_{(x,y)}$ is $H_{(x,y),a}$ for the family $a$ furnished by the light-cone theorem below.
\end{definition}
\begin{leancode}
def vnatStringModule (x y : ℤ) (a : Fin 2) : StringModule (Vnat ⊗[ℂ] hypFock) where
  wt := vnatSectorWeight x y
  wt_isInternal := vnatSectorWeight_isInternal x y
  wt_finite := vnatSectorWeight_finiteDimensional x y
  wt_bdd := ⟨x * y, fun N hN => vnatSectorWeight_negative x y N hN⟩
  c := 26
  L := vnatHypMode x y
  K := vnatHypLongitudinal a x y
  L_wt := vnatSectorWeight_mode x y
  K_wt := fun n N u _ hu => vnatSectorWeight_oscillator x y n N a u hu
  L_zero := vnatSectorWeight_energy x y
  L_comm := vnatHypMode_virasoro x y
  K_comm := vnatHypLongitudinal_comm a x y
  LK_comm := fun m n _ _ => vnatHypLongitudinal_mixed a x y m n
  B := vnatHypForm
  B_symm := ⟨vnatHypForm_symm⟩
  B_orth := vnatSectorWeight_orthogonal x y
  B_nondeg := vnatSectorWeight_nondegenerate x y
  L_adj := vnatHypForm_adjoint x y
  K_adj := vnatHypLongitudinal_adjoint a x y
def stringSectorFamily (x y : ℤ) (h : x ≠ 0 ∨ y ≠ 0) : Fin 2 :=
  (vnatHypMode_lightConeSplit x y (vnatNonzeroMomentum_cast x y h)).choose
def stringSector (x y : ℤ) (h : x ≠ 0 ∨ y ≠ 0) :
    StringModule (Vnat ⊗[ℂ] hypFock) := vnatStringModule x y (stringSectorFamily x y h)
theorem stringSector_charge (x y : ℤ) (h : x ≠ 0 ∨ y ≠ 0) :
    (stringSector x y h).c = 26 := rfl
\end{leancode}
(\leanpath{VOA/Moonshine/MoonshineCNVnatStringModule.lean}.)
\begin{dictionary}
Every field is a proved theorem; there is no hypothesis.  The integer momenta are coerced to $\mathbb C$
in \lean{vnatHypMode x y}.  \lean{stringSectorFamily} is defined by \lean{.choose} from an existence
statement: it is \emph{some} family for which the light-cone splitting holds.  (The proof of the existence statement
exhibits family $0$ when $x\neq0$ and family $1$ otherwise, but \lean{.choose} does not pin the witness, so Lean does not
assert which family is chosen.)  The downstream statements do not depend on which.  The zero-momentum
module \lean{vnatStringModule 0 0 0} is also used (for $\chXVIIIm_0$), but no light-cone splitting is claimed for it.
\end{dictionary}
\begin{physics}
The Fock space of the critical string at momentum $(x,y)$ with internal CFT $\Vnat$, together with the
choice of a light-cone direction $w$ with $(w,r)\neq0$.
\end{physics}

\subsection{The light-cone splitting}

\begin{theorem}[Light-cone splitting (F1) for the sector]\label{ch18:thm:F1}
For every $(\kappa,\ell)\neq(0,0)$ there is a family $a$ such that
$\Vnat\otimes F=\sum_{j\ge0}L_{-(j+1)}(\Vnat\otimes F)\ \oplus\ \bigcap_{j\ge0}\ker K_{j+1}$.  In particular every
string sector $H_{(x,y)}$, $(x,y)\neq0$, satisfies \lean{LightConeSplit}.
\end{theorem}
\begin{leancode}
theorem vnatHypTensor_lightConeSplit (κ ell : ℂ) (h : κ ≠ 0 ∨ ell ≠ 0) :
    ∃ a : Fin 2, IsCompl (⨆ j : ℕ, LinearMap.range (vnatHypNegative κ ell j))
      (⨅ j : ℕ, LinearMap.ker
        ((hypOscillator a (if a = 0 then κ else ell) (j + 1 : ℤ)).lTensor Vnat)) :=
\end{leancode}
(\leanpath{VOA/Moonshine/MoonshineCNVnatLightCone.lean}, with \lean{vnatHypNegative κ ell j} the operator
\lean{vnatHypMode κ ell (-(j + 1))}, by \lean{vnatHypMode_negative}, a \lean{rfl} lemma.)
\begin{leancode}
theorem vnatHypMode_lightConeSplit (κ ell : ℂ) (h : κ ≠ 0 ∨ ell ≠ 0) :
    ∃ a : Fin 2, IsCompl (⨆ j : ℕ, LinearMap.range (vnatHypMode κ ell (-(j + 1 : ℤ))))
      (⨅ j : ℕ, LinearMap.ker
        ((hypOscillator a (if a = 0 then κ else ell) (j + 1 : ℤ)).lTensor Vnat)) := by
theorem stringSector_lightConeSplit (x y : ℤ) (h : x ≠ 0 ∨ y ≠ 0) :
    (stringSector x y h).LightConeSplit :=
\end{leancode}
(\leanpath{VOA/Moonshine/MoonshineCNVnatSectorOperators.lean} and
\leanpath{VOA/Moonshine/MoonshineCNVnatStringModule.lean}.)
\begin{dictionary}
\textbf{The statement} is exactly (F1) for the string module \lean{vnatStringModule x y a}: the left
summand is \lean{spurious} ($\sum_{n\ge1}\mathrm{im}\,L_{-n}$), the right summand is \lean{kVacuum}
($\bigcap_{n\ge1}\ker K_n$).  The docstring of \lean{vnatHypTensor_lightConeSplit} reads ``No
unestablished form, representation, or splitting datum is assumed''.  The general theorem for an
arbitrary space $V$ with an operator family satisfying the negative Virasoro commutators,
\lean{hypTensor_lightConeSplit_nonzero}, is in
\leanpath{VOA/Moonshine/MoonshineCNTensorLightConeKernels.lean} (Chapter~\ref{ch:17}); the present file
specialises it to $\Vnat$ using \lean{nearVirasoroNegative_comm}.

\textbf{Why $(x,y)\neq0$ is necessary.}  At $(x,y)=(0,0)$ the leading coefficient $\kappa$ of
\eqref{ch18:eq:sugawara} vanishes and the splitting genuinely fails: at weight one,
$H^1_{(0,0)}=\operatorname{span}(\mathbb 1\otimes a_1,\mathbb 1\otimes b_1)$ (because $\Vnat_1=0$), the spurious part is
$L_{-1}H^0=\mathbb C\,\mathbb 1\otimes(\kappa b_1+\ell a_1)=0$, and the family-$0$ longitudinal vacua are
$\mathbb C\,\mathbb 1\otimes a_1$; so $\chXVIIIS^1\oplus\chXVIIIK^1$ misses $\mathbb 1\otimes b_1$ (and symmetrically for family $1$).  This is why zero momentum is treated separately (\S\ref{ch18:sec:zero}).
\end{dictionary}
\begin{physics}
``Every Fock state is uniquely a light-cone state dressed by Virasoro descendants'': the Virasoro
generators $L_{-n}$ act, to leading order, as $p^+\alpha^-_{-n}$ (since $L_{-n}=\kappa b_n+\dots$), so
when the light-cone momentum $p^+$ is nonzero the Virasoro descendants can be traded for $\alpha^-$
excitations.  This is the algebraic core of the light-cone gauge.
\end{physics}

\paragraph{Proof on paper.}  (The Lean proof of the general statement is in the files of Chapter~\ref{ch:17}.)  Assume $\kappa\neq0$ and take $K_n=a(n)$, so $\chXVIIIK=\Vnat\otimes\mathbb C[a]$ (polynomials free of $b$'s:
$a(j)=j\partial/\partial b_j$).  Filter $\Vnat\otimes F$ by the pair $(d,w)$ = (number of $b$'s, total $b$-level) of
the monomials, ordered lexicographically.  By \eqref{ch18:eq:sugawara}, for a vector of filtration degree
$(d,w)$, $L_{-n}x-\kappa(1\otimes b_n)x$ has strictly smaller degree than $(d+1,w+n)$ (all other terms of
$L_{-n}$ either keep $d$, or raise $d$ by one but $w$ by less than $n$).  Iterating, for $k\in\chXVIIIK$ and a sorted word
$\lambda=(n_1\ge\dots\ge n_s)$, $L_{-\lambda}k=\kappa^s(1\otimes b_\lambda)k+(\text{lower})$.  Hence the map
$\Phi:\mathbb C[b]\otimes\chXVIIIK\to H$, $b^\beta\otimes k\mapsto L_{-\lambda(\beta)}k$, is ``upper triangular with nonzero
diagonal $\kappa^d$'' with respect to the isomorphism $\mu:\mathbb C[b]\otimes\chXVIIIK\cong H$, $b^\beta\otimes k\mapsto
b^\beta k$; so $\Phi$ is injective, and bijective on each (finite-dimensional) weight space.
$\Phi(1\otimes\chXVIIIK)=\chXVIIIK$, and by straightening (Theorem~\ref{ch18:thm:straightening}) the spurious space is
$\Phi(\mathbb C[b]_+\otimes\chXVIIIK)$.  So $H=\chXVIIIK\oplus\chXVIIIS$.  If $\kappa=0\neq\ell$, exchange the roles of $a$ and $b$. \qed

\begin{status}
No hypotheses beyond $(\kappa,\ell)\neq0$; standard axioms (\leanpath{VOA/Moonshine/MoonshineCNAudit7.lean},
\leanpath{VOA/Moonshine/MoonshineCNAudit8.lean}).  Unlike most no-ghost presentations, no Kac determinant,
no Shapovalov form for $\mathrm{Vir}\ltimes\text{Heisenberg}$ and no character count is used.
\end{status}

\subsection{The automorphism group acts}

\begin{theorem}[$\Aut(\Vnat)$ preserves the form and acts on every sector]\label{ch18:thm:sectoraction}
Every $g\in\Aut(\Vnat)$ satisfies $B^\natural(gv,gw)=B^\natural(v,w)$.  Consequently $g\otimes1$ defines a
symmetry (Definition~\ref{ch18:def:stringgroup}) of every string module $H_{(x,y),a}$.
\end{theorem}
\begin{leancode}
theorem vnatAut_preserves_natForm (g : AutVnat) (v w : Vnat) :
    natForm (g v) (g w) = natForm v w := by
def vnatSectorAction : AutVnat →* ((Vnat ⊗[ℂ] hypFock) ≃ₗ[ℂ] (Vnat ⊗[ℂ] hypFock)) where
theorem vnatSectorAction_tmul (g : AutVnat) (v : Vnat) (p : hypFock) :
    vnatSectorAction g (v ⊗ₜ[ℂ] p) = g v ⊗ₜ[ℂ] p := rfl
def vnatStringGroupAction (x y : ℤ) (a : Fin 2) :
    StringGroupAction (vnatStringModule x y a) AutVnat where
\end{leancode}
(\leanpath{VOA/Moonshine/MoonshineCNVnatFormAction.lean} for the first,
\leanpath{VOA/Moonshine/MoonshineCNVnatSectorAction.lean} for the rest; \lean{stringSectorAction} is the
specialisation to \lean{stringSector x y h}.)
\begin{dictionary}
The docstring stresses: ``The actual \lean{natForm}, not an unspecified invariant form, is
Aut($\Vnat$)-invariant.''  Proof idea: $v\mapsto B^\natural(gv,gw)$ defines another invariant form
\lean{vnatAutPullbackForm g}, which is fully (FHL-)invariant (\lean{vnatAutPullbackForm_fhl}); by the
uniqueness of invariant forms on a simple CFT-type VOA with $L_1V_1=0$ (Chapter~\ref{ch:L04}) it is a
multiple of $B^\natural$, and the multiple is $1$ because $g\mathbb 1=\mathbb 1$.  The four laws of
\lean{StringGroupAction} are \lean{vnatSectorAction_weight}, \lean{vnatSectorAction_mode},
\lean{vnatSectorAction_longitudinal}, \lean{vnatSectorAction_isometry}.
\end{dictionary}
\begin{physics}
The Monster is a global symmetry of the internal CFT; it commutes with the light-cone sector and is
unitary for the string inner product.
\end{physics}
\begin{status}No hypotheses; standard axioms (\leanpath{VOA/Moonshine/MoonshineCNAudit8.lean}).\end{status}

\section{The no-ghost theorem for the moonshine module}\label{ch18:sec:vnatnoghost}

\subsection{The longitudinal quotient is the matter weight space}

For family $a$, the longitudinal vacua are the vectors whose polynomial parts involve only family-$a$
variables (\lean{vnatStringModule_kVacuum}: \lean{kVacuum = hypTensorOneFamily a}).  Let
$\mathrm{ev}_0:\Vnat\otimes F\to\Vnat$ be ``take the coefficient of the constant monomial'':
\begin{leancode}
def vnatSectorEval : (Vnat ⊗[ℂ] hypFock) →ₗ[ℂ] Vnat :=
  (Finsupp.lapply 0).comp hypTensorCoordinates.toLinearMap
theorem vnatSectorEval_tmul (v : Vnat) (p : hypFock) :
    vnatSectorEval (v ⊗ₜ[ℂ] p) = coeff 0 p • v := hypTensorCoordinates_tmul v p 0
\end{leancode}
(\leanpath{VOA/Moonshine/MoonshineCNVnatSectorEvaluation.lean}.)

\begin{theorem}[Longitudinal radical quotient]\label{ch18:thm:kquotient}
For every $(x,y)$, family $a$ and weight $N$, evaluation induces an isometric, $\Aut(\Vnat)$-equivariant
isomorphism $\chXVIIIK^N/\mathrm{rad}\cong\Vnat_{N-xy}$; its kernel on $\chXVIIIK^N$ is the radical, which equals
$\chXVIIIK^N\cap\chXVIIIS_K$.
\end{theorem}
\begin{leancode}
theorem vnatKVacuum_augmentation (x y : ℤ) (a : Fin 2) (N : ℤ)
    (u : (vnatStringModule x y a).weightSlice (vnatStringModule x y a).kVacuum N) :
    (u.val : Vnat ⊗[ℂ] hypFock) ∈ (vnatStringModule x y a).kSpurious ↔
      vnatKVacuumEvaluation x y a N u = 0 := by
def vnatKVacuumQuotientEquiv (x y : ℤ) (a : Fin 2) (N : ℤ) :
    (vnatStringModule x y a).weightRadicalQuotient (vnatStringModule x y a).kVacuum N ≃ₗ[ℂ]
      VnatVOA.grading (N - x * y) :=
theorem vnatKVacuumQuotientEquiv_isometry (x y : ℤ) (a : Fin 2) (N : ℤ)
    (u v : (vnatStringModule x y a).weightRadicalQuotient (vnatStringModule x y a).kVacuum N) :
    natForm (vnatKVacuumQuotientEquiv x y a N u) (vnatKVacuumQuotientEquiv x y a N v) =
      (vnatStringModule x y a).radicalQuotientForm (vnatStringModule x y a).kVacuum N u v := by
theorem vnatKVacuumQuotientEquiv_equivariant (x y : ℤ) (a : Fin 2) (N : ℤ) (g : AutVnat)
    (u : (vnatStringModule x y a).weightRadicalQuotient (vnatStringModule x y a).kVacuum N) :
    vnatKVacuumQuotientEquiv x y a N
      ((vnatStringGroupAction x y a).stableWeightQuotientAction _
        (vnatStringGroupAction x y a).action_preserves_kVacuum N g u) =
      GriessAut.autRestrict g (N - x * y) (vnatKVacuumQuotientEquiv x y a N u) := by
\end{leancode}
(\leanpath{VOA/Moonshine/MoonshineCNVnatLongitudinalQuotient.lean}.)
\begin{dictionary}
\lean{vnatKVacuumEvaluation x y a N} is $\mathrm{ev}_0$ restricted to $\chXVIIIK^N$, with values in
$\Vnat_{N-xy}$ (\lean{vnatSectorEval_weight}).  The equivariance statement compares the action of $g$ on the
quotient (descended from $g\otimes1$) with $g|_{\Vnat_{N-xy}}$.  No hypotheses: in particular this holds
at zero momentum and for both families.
\end{dictionary}
\begin{physics}
Longitudinal vacua modulo null states are the internal states: the $\alpha^+$ excitations are pure
gauge, and what survives is the internal CFT state at the level fixed by the mass shell.
\end{physics}
\paragraph{Proof on paper.}  A family-$a$ polynomial is a constant plus an element of the ideal generated by
the family-$a$ variables, so $\chXVIIIK^N=(\Vnat_{N-xy}\otimes1)\oplus(\Vnat\otimes\mathbb C[a]_+)^N$; the second summand is
$\chXVIIIK\cap\chXVIIIS_K$ ($K_{-n}$ multiplies by $a_n$), i.e.\ the kernel of $\mathrm{ev}_0$, and $\mathrm{ev}_0$ is onto. On
$\Vnat\otimes\mathbb C[a]$ the form is $B^\natural(\mathrm{ev}_0u,\mathrm{ev}_0v)$, because $B_F$ pairs an $a$-monomial only with the
corresponding $b$-monomial, and the only $b$-free monomial among $a$-monomials is $1$
(\lean{vnatSectorEval_oneFamily_form}).  Equivariance: $g\otimes1$ commutes with $\mathrm{ev}_0$. \qed
\begin{status}No hypotheses; standard axioms (\leanpath{VOA/Moonshine/MoonshineCNAudit8.lean}).\end{status}

\subsection{The theorem}

\begin{theorem}[No-ghost theorem for $\Vnat$, nonzero momentum]\label{ch18:thm:noghost}
For every $(x,y)\in\Z^2\setminus\{0\}$ there is a linear isomorphism
\begin{equation*}
  \Theta_{(x,y)}:\ \chXVIIIP^1_{(x,y)}\big/\mathrm{rad}\bigl(B|_{\chXVIIIP^1_{(x,y)}}\bigr)\ \xrightarrow{\ \sim\ }\ \Vnat_{1-xy}
\end{equation*}
which is an isometry for the descended form and $B^\natural$, and intertwines the action of every
$g\in\Aut(\Vnat)$.  Here $\chXVIIIP^1_{(x,y)}$ is the weight-one physical space of the string sector $H_{(x,y)}$, and
$\mathrm{rad}=\chXVIIIP^1_{(x,y)}\cap\chXVIIIS$.  In particular $\dim\chXVIIIP^1_{(x,y)}/\mathrm{rad}=\dim\Vnat_{1-xy}$.
\end{theorem}
\begin{leancode}
def stringSectorNoGhostEquiv (x y : ℤ) (h : x ≠ 0 ∨ y ≠ 0) :
    (stringSector x y h).weightRadicalQuotient (stringSector x y h).physical 1 ≃ₗ[ℂ]
      VnatVOA.grading (1 - x * y) :=
  ((stringSector x y h).physicalKQuotientEquiv (stringSector_charge x y h)
    (stringSector_lightConeSplit x y h)).trans
      (vnatKVacuumQuotientEquiv x y (stringSectorFamily x y h) 1)
theorem stringSectorNoGhostEquiv_isometry (x y : ℤ) (h : x ≠ 0 ∨ y ≠ 0)
    (u v : (stringSector x y h).weightRadicalQuotient (stringSector x y h).physical 1) :
    natForm (stringSectorNoGhostEquiv x y h u) (stringSectorNoGhostEquiv x y h v) =
      (stringSector x y h).radicalQuotientForm (stringSector x y h).physical 1 u v := by
theorem stringSectorNoGhostEquiv_equivariant (x y : ℤ) (h : x ≠ 0 ∨ y ≠ 0)
    (g : AutVnat)
    (u : (stringSector x y h).weightRadicalQuotient (stringSector x y h).physical 1) :
    stringSectorNoGhostEquiv x y h
      ((stringSectorAction x y h).stableWeightQuotientAction _
        (stringSectorAction x y h).action_preserves_physical 1 g u) =
      GriessAut.autRestrict g (1 - x * y) (stringSectorNoGhostEquiv x y h u) := by
\end{leancode}
(\leanpath{VOA/Moonshine/MoonshineCNVnatNoGhost.lean}.)
\begin{leancode}
theorem stringSectorNoGhost_finrank (x y : ℤ) (h : x ≠ 0 ∨ y ≠ 0) :
    Module.finrank ℂ
      ((stringSector x y h).weightRadicalQuotient (stringSector x y h).physical 1) =
      Module.finrank ℂ (VnatVOA.grading (1 - x * y)) :=
\end{leancode}
(\leanpath{VOA/Moonshine/MoonshineCNVnatNoGhostConsequences.lean}.)

\begin{dictionary}
\textbf{Decoding.}  The source is the type \lean{weightRadicalQuotient physical 1}: the weight-one slice of
the physical space of the actual string sector, modulo the radical of the actual form restricted to it
(\S\ref{ch18:sec:abstractnoghost}).  The target is \lean{VnatVOA.grading (1 - x * y)}, the weight space
$\Vnat_{1-xy}$; the weight is \emph{exactly} $1-xy$ in the repository norm $2xy$ (the file docstring:
``The norm convention is \lean{2*x*y}; the matter weight is exactly \lean{1-x*y}'').  The definition is a
composition of the abstract isomorphism \lean{physicalKQuotientEquiv} (which needs $c=26$:
\lean{stringSector_charge}, a \lean{rfl}; and (F1): \lean{stringSector_lightConeSplit}) with the
longitudinal quotient of Theorem~\ref{ch18:thm:kquotient}.

\textbf{Hypotheses.}  Only $h:(x,y)\neq0$.  Every datum is concrete; the docstring: ``The actual physical
radical quotient at weight one, with no unproved data as inputs.''

\textbf{Same as / stronger than / weaker than.}  Stronger than the bare dimension statement of
Goddard--Thorn: it is an isometry and is $\Aut(\Vnat)$-equivariant.  It is about the \emph{model} sector
$\Vnat\otimes F$; the transfer to the vertex algebra $\Vnat\otimes V_{\chXVIIIII}$ is
\S\ref{ch18:sec:roots}.  No positivity: the form on $\Vnat_{1-xy}$ is \lean{natForm}, which is nondegenerate
but not claimed positive.

\textbf{Could it be vacuous?}  No.  When $xy\ge2$ both sides are $0$, and when $xy=0$ (one coordinate zero) both are
$0$ because $\Vnat_1=0$; but when $xy=1$ ($\Vnat_0=\mathbb C$) and when $xy\le-1$ ($\Vnat_{1-xy}\neq0$) the theorem is a genuine
statement about nonzero spaces, e.g.\ $\dim=196884$ at $xy=-1$.
\end{dictionary}
\begin{physics}
The physical states of the compactified critical string at momentum $r\neq0$, modulo null states, are
in bijection with the internal states at level $1-r^2/2$, with the same inner products and the same
Monster action.  This is the light-cone gauge spectrum: no ghosts, no longitudinal excitations, only
the $24$ ``transverse'' internal dimensions, here replaced by $\Vnat$.
\end{physics}
\paragraph{Proof on paper.}  Apply Theorem~\ref{ch18:thm:quotientequivs} to the string sector, using
Theorem~\ref{ch18:thm:F1} for (F1) and $c=26$ from Theorem~\ref{ch18:thm:sectorlaws}, to get
$\chXVIIIP^1/\mathrm{rad}\cong\chXVIIIK^1/\mathrm{rad}$ isometrically and equivariantly; then apply Theorem~\ref{ch18:thm:kquotient} at
$N=1$. \qed
\begin{status}
Only $(x,y)\neq0$.  \texttt{\#print axioms}: standard axioms (\leanpath{VOA/Moonshine/MoonshineCNAudit8.lean}
prints \lean{stringSectorNoGhostEquiv} and its isometry and equivariance lemmas and
\lean{stringSectorNoGhost_finrank}).
\end{status}

\begin{warning}
\textbf{Isometry is at the level of the sector form.}  The isometry of Theorem~\ref{ch18:thm:noghost} is for the
\emph{contravariant} sector form $B^\natural\otimes B_F$ (same momentum pairs with same momentum).  The natural form
on the vertex algebra $\Vnat\otimes V_{\chXVIIIII}$ (and on $\chXVIIIm$) pairs momentum $r$ with $-r$; the two are related
by momentum reversal and a cocycle sign (\S\ref{ch18:sec:form}).  A literal isometry between a sector and
the same-momentum part of the vertex algebra \emph{does not exist}: this natural first guess fails, and its failure is a
Lean theorem (Theorem~\ref{ch18:thm:nosame}).
\end{warning}

\subsection{Zero momentum}\label{ch18:sec:zero}

At $r=0$ the light-cone splitting fails, but the weight-one space is tiny: $\Vnat_1=0$ and $\Vnat_0=\mathbb C\mathbb 1$, so
$H^1_{(0,0)}=\mathbb 1\otimes F_1=\operatorname{span}(\mathbb 1\otimes a_1,\mathbb 1\otimes b_1)$.

\begin{theorem}[Zero momentum: the hyperbolic plane]\label{ch18:thm:zeromomentum}
The weight-one space of $H_{(0,0),0}$ is spanned by $\mathbb 1\otimes a_1$ and $\mathbb 1\otimes b_1$, both physical; the radical of
the form on $\chXVIIIP^1_{(0,0)}$ is zero; and $\chXVIIIP^1_{(0,0)}/\mathrm{rad}\cong\mathbb C^2$ with form
$((c_0,c_1),(d_0,d_1))\mapsto c_0d_1+c_1d_0$ and trivial $\Aut(\Vnat)$-action.
\end{theorem}
\begin{leancode}
def vnatZeroBasisVector (a : Fin 2) : Vnat ⊗[ℂ] hypFock :=
  VnatVOA.vacuum ⊗ₜ[ℂ] X (0, a)
theorem vnatZeroWeight_span :
    vnatSectorWeight 0 0 1 = Submodule.span ℂ (Set.range vnatZeroBasisVector) := by
theorem vnatZeroWeight_physical :
    vnatSectorWeight 0 0 1 ≤ (vnatStringModule 0 0 0).physical := by
\end{leancode}
(\leanpath{VOA/Moonshine/MoonshineCNVnatZeroWeight.lean}.)
\begin{leancode}
theorem vnatZeroPhysicalRadical :
    (vnatStringModule 0 0 0).weightRadical (vnatStringModule 0 0 0).physical 1 = ⊥ := by
def vnatZeroNoGhostEquiv :
    (vnatStringModule 0 0 0).weightRadicalQuotient (vnatStringModule 0 0 0).physical 1 ≃ₗ[ℂ]
      (Fin 2 → ℂ) :=
theorem vnatZeroNoGhostEquiv_isometry
    (u v : (vnatStringModule 0 0 0).weightRadicalQuotient (vnatStringModule 0 0 0).physical 1) :
    (vnatStringModule 0 0 0).radicalQuotientForm (vnatStringModule 0 0 0).physical 1 u v =
      vnatZeroNoGhostEquiv u 0 * vnatZeroNoGhostEquiv v 1 +
        vnatZeroNoGhostEquiv u 1 * vnatZeroNoGhostEquiv v 0 := by
theorem vnatZeroNoGhostEquiv_equivariant (g : AutVnat)
    (u : (vnatStringModule 0 0 0).weightRadicalQuotient (vnatStringModule 0 0 0).physical 1) :
    vnatZeroNoGhostEquiv
      ((vnatStringGroupAction 0 0 0).stableWeightQuotientAction _
        (vnatStringGroupAction 0 0 0).action_preserves_physical 1 g u) =
      vnatZeroNoGhostEquiv u := by
\end{leancode}
(\leanpath{VOA/Moonshine/MoonshineCNVnatZeroQuotient.lean}; also \lean{vnatZeroNoGhost_finrank} $=2$ in
\leanpath{VOA/Moonshine/MoonshineCNVnatNoGhostConsequences.lean}.)
\begin{dictionary}
\lean{X (0, a)} is $a_1$ for \lean{a = 0} and $b_1$ for \lean{a = 1}.  The coordinates of
\lean{vnatZeroNoGhostEquiv} are \lean{vnatZeroWeightCoordinates}: the coordinate $a$ of $u$ is
$B^\natural(\mathbb 1,\text{coefficient of }X(0,a)\text{ in }u)$, so $c_0\,\mathbb 1\otimes a_1+c_1\,\mathbb 1\otimes b_1\mapsto(c_0,c_1)$ (the
normalisation $B^\natural(\mathbb 1,\mathbb 1)=1$ is \lean{natInvForm_normalized}).  The form $c_0d_1+c_1d_0$ is the Gram
matrix $\left(\begin{smallmatrix}0&1\\1&0\end{smallmatrix}\right)$: $\mathbb 1\otimes a_1$ and $\mathbb 1\otimes b_1$ are isotropic and pair to
$1$ (\lean{vnatZeroBasisVector_form}).  No hypotheses; uses $\Vnat_1=0$ and $\Vnat_0=\mathbb C\mathbb 1$.
\end{dictionary}
\begin{physics}
At zero momentum the only weight-one physical states are the two light-cone currents $\partial X^\pm$
(the internal CFT has no currents since $\Vnat_1=0$): they span the Cartan subalgebra, with the
Lorentzian metric of $\chXVIIIII$, and the Monster acts trivially on them.
\end{physics}
\paragraph{Proof on paper.}  $\Vnat_{<0}=0$ and $\Vnat_1=0$ leave only $\mathbb 1\otimes F_1$ at weight one.  $L_n$ for $n\ge2$
maps it to negative weight, hence to $0$; $L_1(\mathbb 1\otimes a_1)=\mathbb 1\otimes L^F_1a_1$ (as $L^\natural_1\mathbb 1=0$), and by the explicit lowering derivation
$L^F_1a_1=\kappa$, $L^F_1b_1=\ell$, both $0$ at zero momentum.  The spurious space at weight
one is $L_{-1}H^0=\mathbb C L_{-1}(\mathbb 1\otimes1)=\mathbb C\,\mathbb 1\otimes(\kappa b_1+\ell a_1)=0$; so the radical, which is
$\chXVIIIP^1\cap\chXVIIIS$, is $0$.  The automorphisms fix $\mathbb 1$. \qed
\begin{status}No hypotheses; standard axioms (\leanpath{VOA/Moonshine/MoonshineCNAudit8.lean}).\end{status}

\section{The Lie algebra of physical states}\label{ch18:sec:lie}

\subsection{Physical states of a near vertex operator algebra}

The bracket is defined for an arbitrary near vertex operator algebra $V$ (Chapter~\ref{ch:L03}); the Borcherds
identity, commutator formula and skew symmetry for near-VOAs, which the proofs below use, are proved in
Chapter~\ref{ch:17} (files \lean{MoonshineCNNear*}).

\begin{definition}[Weight-one primaries, spurious primaries, the physical Lie algebra]\label{ch18:def:primary}
Let $V$ be a near-VOA with Virasoro modes $L_n$.  Put
\begin{equation*}
  P^1(V):=\{v\in V: L_0v=v,\ L_nv=0\ (n>0)\},\qquad \mathrm{Sp}(V):=\sum_{n\ge1}L_{-n}V,
\end{equation*}
$\mathrm{Null}^1(V):=P^1(V)\cap\mathrm{Sp}(V)$, and $\mathfrak g(V):=P^1(V)/\mathrm{Null}^1(V)$ with $[\bar u,\bar v]:=\overline{u_{(0)}v}$.
\end{definition}
\begin{leancode}
def primaryWeightOne : Submodule ℂ H where
  carrier := {v | V.L 0 v = v ∧ ∀ n : ℤ, 0 < n → V.L n v = 0}
noncomputable def primaryZeroProduct :
    V.primaryWeightOne →ₗ[ℂ] V.primaryWeightOne →ₗ[ℂ] V.primaryWeightOne where
noncomputable def primarySpurious : Submodule ℂ V.primaryWeightOne :=
  V.spurious.comap V.primaryWeightOne.subtype
abbrev primaryLieQuotient := V.primaryWeightOne ⧸ V.primarySpurious
noncomputable def primaryLieBracket :
    V.primaryLieQuotient →ₗ[ℂ] V.primaryLieQuotient →ₗ[ℂ] V.primaryLieQuotient :=
theorem primaryQuotient_mk_bracket (u v : V.primaryWeightOne) :
    V.primaryLieBracket (V.primarySpurious.mkQ u) (V.primarySpurious.mkQ v) =
      V.primarySpurious.mkQ (V.primaryZeroProduct u v) := rfl
noncomputable instance primaryLieRing : LieRing V.primaryLieQuotient where
noncomputable instance primaryLieAlgebra : LieAlgebra ℂ V.primaryLieQuotient where
\end{leancode}
(\leanpath{VOA/Moonshine/MoonshineCNPhysicalLie.lean}, namespace \lean{VOA.NearVOA}; \lean{V.spurious} is
\lean{⨆ n : ℕ, LinearMap.range (V.L (-(n + 1 : ℤ)))} from
\leanpath{VOA/Moonshine/MoonshineCNNearConsequences.lean}.)

\begin{dictionary}
\textbf{Decoding.}  \lean{primaryZeroProduct u v} is $u_{(0)}v$, i.e.\ \lean{V.smul u 0 v} $=$
\lean{mode (V.Y u) 0 v}, packaged as a bilinear map of $P^1$ to itself (closure is
\lean{primary_zeroMode_closed}).  \lean{primarySpurious} is $\mathrm{Sp}(V)\cap P^1$ as a subspace of $P^1$.
\lean{⧸} is the quotient, \lean{mkQ} the quotient map.  The \lean{LieRing} instance takes
\lean{primaryLieBracket} as the bracket; its axioms \lean{lie_self} and \lean{leibniz_lie} are discharged by
\lean{primary_zeroMode_self_spurious} and \lean{primary_zeroMode_leibniz}.

\textbf{No form, no central charge.}  The module docstring: ``defined without a bilinear form'', and
\lean{primaryLieQuotient}'s: ``No form or charge assumption is used.''  Borcherds defines $\chXVIIIm$ as $P^1$ modulo the
radical of the form, and Frenkel's ``$P^1/L_{-1}P^0$'' is a third common description; Remark~\ref{ch18:rem:threedescriptions}
shows that all three agree for $\Vnat\otimes V_{\chXVIIIII}$.

\textbf{Why $u_{(0)}$.}  This is the same bracket as on the weight-one space of a vertex operator algebra of CFT type, where
$[u,v]=u_{(0)}v$ makes $V_1$ a Lie algebra (Chapter~\ref{ch:L04}): the zero mode $u_{(0)}=\oint Y(u,z)\,dz$ of a weight-one
state is the charge of the current $Y(u,z)$, and $[u_{(0)},v_{(0)}]=(u_{(0)}v)_{(0)}$ by the commutator formula.  The only change
is that $V_1$ is replaced by the primaries $P^1$ (because $\Vnat\otimes V_{\chXVIIIII}$ is not of CFT type and its weight-one space is
infinite-dimensional), and one must divide by the null states, which form an ideal.
\end{dictionary}
\begin{physics}
Weight-one primaries are conserved currents; their zero modes are conserved charges, and the bracket is
the commutator of charges.  Spurious primaries (null states) are pure gauge and decouple.
\end{physics}

\begin{theorem}[$\mathfrak g(V)$ is a Lie algebra]\label{ch18:thm:primarylie}
For every near-VOA $V$: (a) $u_{(0)}P^1\subseteq P^1$ for $u\in P^1$; (b) $u_{(0)}\mathrm{Null}^1\subseteq\mathrm{Null}^1$ and $(\mathrm{Null}^1)_{(0)}P^1\subseteq\mathrm{Null}^1$;
(c) $u_{(0)}u\in\mathrm{Null}^1$; (d) $u_{(0)}(v_{(0)}w)=(u_{(0)}v)_{(0)}w+v_{(0)}(u_{(0)}w)$.  Hence the bracket is well defined,
alternating and satisfies the Jacobi identity.
\end{theorem}
\begin{leancode}
theorem primary_zeroMode_closed (u v : V.primaryWeightOne) :
    V.smul u 0 v ∈ V.primaryWeightOne := by
theorem primarySpurious_right (u v : V.primaryWeightOne) (hv : v ∈ V.primarySpurious) :
    V.primaryZeroProduct u v ∈ V.primarySpurious :=
theorem primarySpurious_left (u v : V.primaryWeightOne) (hu : u ∈ V.primarySpurious) :
    V.primaryZeroProduct u v ∈ V.primarySpurious := by
theorem primary_zeroMode_self_spurious (u : V.primaryWeightOne) :
    V.primaryZeroProduct u u ∈ V.primarySpurious := by
theorem primary_zeroMode_leibniz (u v w : V.primaryWeightOne) :
    V.primaryZeroProduct u (V.primaryZeroProduct v w) =
      V.primaryZeroProduct (V.primaryZeroProduct u v) w +
        V.primaryZeroProduct v (V.primaryZeroProduct u w) := by
\end{leancode}
(\leanpath{VOA/Moonshine/MoonshineCNPhysicalLie.lean}.)
\begin{physics}
Charges of conserved currents form a Lie algebra under commutation, and gauge (null) currents form an
ideal.
\end{physics}
\paragraph{Proof on paper.}  For $u\in P^1$ the commutator formula with $\omega$ gives
$[L_m,u_{(n)}]=-n\,u_{(m+n)}$ (because $L_0u=u$, $L_ju=0$ for $j\ge1$, and $(L_{-1}u)_{(k)}=-ku_{(k-1)}$); at
$n=0$, $u_{(0)}$ commutes with every $L_m$ (\lean{primary_zeroMode_commutes}).  Hence $u_{(0)}$ preserves $P^1$
(it preserves weights) and preserves $\mathrm{Sp}$ ($u_{(0)}L_{-n}y=L_{-n}u_{(0)}y$): (a) and the first half of (b).  Skew
symmetry at $n=0$ gives $v_{(0)}u=-u_{(0)}v+\sum_{j\ge1}(-1)^{j+1}L_{-1}^j(u_{(j)}v)/j!$, so
$u_{(0)}v+v_{(0)}u\in\mathrm{Sp}$; this gives (c) (take $v=u$: $2u_{(0)}u\in\mathrm{Sp}$) and the second half of (b).  (d) is the
commutator formula at $m=n=0$: $[u_{(0)},v_{(0)}]=(u_{(0)}v)_{(0)}$. \qed
\begin{status}No hypotheses beyond $V$; standard axioms (\leanpath{VOA/Moonshine/MoonshineCNAudit10.lean}).\end{status}

\paragraph{Functoriality.}  A linear map $f$ of $V$ commuting with all $L_n$ preserves $P^1$ and $\mathrm{Sp}$ and induces
a linear map of $\mathfrak g(V)$ (\lean{primaryQuotientMap}); an invertible such map that also preserves $u_{(0)}v$
induces a Lie algebra automorphism:
\begin{leancode}
noncomputable def primaryQuotientLieEquiv (e : H ≃ₗ[ℂ] H)
    (he : ∀ (n : ℤ) (v : H), e (V.L n v) = V.L n (e v))
    (hzero : ∀ u v : H, e (V.smul u 0 v) = V.smul (e u) 0 (e v)) :
    V.primaryLieQuotient ≃ₗ⁅ℂ⁆ V.primaryLieQuotient := by
\end{leancode}
(\leanpath{VOA/Moonshine/MoonshineCNPhysicalMaps.lean}; \lean{≃ₗ⁅ℂ⁆} is a Lie algebra isomorphism.)
\begin{physics}A global symmetry of the string background acts on the gauge algebra by automorphisms.\end{physics}

\subsection{The Monster Lie algebra}

\begin{definition}[The Monster Lie algebra]\label{ch18:def:monsterlie}
$\chXVIIIm:=\mathfrak g(\Vnat\otimes V_{\chXVIIIII})=P^1/(P^1\cap\mathrm{Sp})$ for the matter near-VOA
$\Vnat\otimes V_{\chXVIIIII}$ of Chapter~\ref{ch:17} (central charge $26$).
\end{definition}
\begin{leancode}
abbrev monsterLieAlgebra := moonshineMatterNearVOA.primaryLieQuotient
instance monsterLieRing : LieRing monsterLieAlgebra := inferInstance
instance monsterLieAlgebraComplex : LieAlgebra ℂ monsterLieAlgebra := inferInstance
def monsterLieProjection : moonshineMatterNearVOA.primaryWeightOne →ₗ[ℂ] monsterLieAlgebra :=
  moonshineMatterNearVOA.primarySpurious.mkQ
\end{leancode}
(\leanpath{VOA/Moonshine/MoonshineCNMonsterLie.lean}.)  The matter near-VOA is
\begin{leancode}
abbrev moonshineMatterSpace := Vnat ⊗[ℂ] LorentzianFockSpace
noncomputable def moonshineMatterNearVOA : NearVOA moonshineMatterSpace :=
  nearTensorVOA VnatVOA.toNearVOA lorentzianNearVOA
\end{leancode}
(\leanpath{VOA/Moonshine/MoonshineCNTensorRealization.lean}, Chapter~\ref{ch:17}; its central charge is
$26$, \lean{moonshineMatterNearVOA_centralCharge}.)
\begin{dictionary}
\lean{LorentzianFockSpace} is the lattice Fock space $\bigoplus_{r\in\chXVIIIII}S(\hat{\mathfrak h}^-)\otimes e^r$ of
$V_{\chXVIIIII}$; \lean{nearTensorVOA} is the tensor product of near-VOAs with the total grading and the conformal
vector $\omega^\natural\otimes\mathbb 1+\mathbb 1\otimes\omega_L$.  So \lean{monsterLieAlgebra} is literally $P^1/(P^1\cap\mathrm{Sp})$ of
$\Vnat\otimes V_{\chXVIIIII}$.  The \lean{inferInstance} lines say that the Lie structure is the one of
Theorem~\ref{ch18:thm:primarylie}.

\textbf{Could it be trivial?}  No: Theorems~\ref{ch18:thm:rootspaces} and \ref{ch18:thm:zerospace} below identify its
graded pieces with $\Vnat_{1-xy}$ and $\mathbb C^2$; and the bracket is nonzero, since the Cartan part acts on
$\chXVIIIm_{(x,y)}$ by the nonzero scalars of Theorem~\ref{ch18:thm:cartanaction}.
\end{dictionary}
\begin{physics}
The gauge algebra of the critical bosonic string compactified on the Monster CFT times the
two-dimensional Lorentzian torus $\mathbb R^{1,1}/\chXVIIIII$.
\end{physics}
\begin{status}Definition; \texttt{\#print axioms monsterLieAlgebra} standard (\leanpath{VOA/Moonshine/MoonshineCNAudit10.lean}).\end{status}

\subsection{The momentum grading}

\begin{theorem}[$\chXVIIIm$ is $\chXVIIIII$-graded]\label{ch18:thm:grading}
Let $\chXVIIIm_r$ be the image in $\chXVIIIm$ of the physical states of momentum $r$.  Then $\chXVIIIm=\bigoplus_{r\in\chXVIIIII}\chXVIIIm_r$
(internal direct sum) and $[\chXVIIIm_r,\chXVIIIm_s]\subseteq\chXVIIIm_{r+s}$.
\end{theorem}
\begin{leancode}
def monsterLieMomentum (α : lorentzianLattice.carrier) : Submodule ℂ monsterLieAlgebra :=
  ((moonshineMatterMomentum α).comap moonshineMatterNearVOA.primaryWeightOne.subtype).map
    monsterLieProjection
theorem monsterLieMomentum_bracket (α β : lorentzianLattice.carrier)
    (u v : monsterLieAlgebra) (hu : u ∈ monsterLieMomentum α)
    (hv : v ∈ monsterLieMomentum β) : ⁅u, v⁆ ∈ monsterLieMomentum (α + β) := by
\end{leancode}
(\leanpath{VOA/Moonshine/MoonshineCNMonsterLie.lean}.)
\begin{leancode}
theorem monsterLieMomentum_isInternal : DirectSum.IsInternal monsterLieMomentum := by
\end{leancode}
(\leanpath{VOA/Moonshine/MoonshineCNMonsterGrading.lean}.)
\begin{dictionary}
\lean{moonshineMatterMomentum α} is $\Vnat\otimes(S(\hat{\mathfrak h}^-)\otimes e^\alpha)$ (Chapter~\ref{ch:17}).  So
\lean{monsterLieMomentum α} is the image of $P^1\cap(\text{momentum }\alpha)$; the docstring of
\lean{monsterLieMomentum_isInternal}: ``The quotient is the direct sum of its momentum subspaces, not
merely their span.''
\end{dictionary}
\begin{physics}Momentum conservation: the charge of momentum $r$ maps momentum-$s$ states to momentum $r+s$.\end{physics}
\paragraph{Proof on paper.}  $\omega$ has momentum $0$ (\lean{moonshineMatterConformal_momentum}), so every $L_n$
preserves momentum and commutes with the projections $\pi_\alpha$ onto momentum components
(\lean{matterMomentumComponent_virasoro}).  Hence $\pi_\alpha$ preserves $P^1$ and $\mathrm{Sp}$ and descends to an
idempotent \lean{monsterLieMomentumComponent α} on $\chXVIIIm$ with image $\chXVIIIm_\alpha$; these idempotents are
orthogonal and sum (pointwise finitely) to the identity, which is internal directness.  Momenta add under
$u_{(0)}$ because $e^\alpha_{(n)}e^\beta\in e^{\alpha+\beta}$-sector. \qed
\begin{status}No hypotheses; standard axioms (\leanpath{VOA/Moonshine/MoonshineCNAudit10.lean}).\end{status}

\subsection{\texorpdfstring{The action of $\Aut(\Vnat)$}{The action of Aut(V-natural)}}

\begin{theorem}[$\Aut(\Vnat)$ acts on $\chXVIIIm$ by graded Lie automorphisms]\label{ch18:thm:autaction}
$g\mapsto(\bar u\mapsto\overline{(g\otimes1)u})$ is a group homomorphism $\Aut(\Vnat)\to\Aut_{\mathrm{Lie}}(\chXVIIIm)$ preserving
every $\chXVIIIm_r$.
\end{theorem}
\begin{leancode}
def moonshineMatterAut : AutVnat →* (moonshineMatterSpace ≃ₗ[ℂ] moonshineMatterSpace) where
  toFun g := g.toLinearEquiv.rTensor LorentzianFockSpace
def monsterLieAut : AutVnat →* (monsterLieAlgebra ≃ₗ⁅ℂ⁆ monsterLieAlgebra) where
  toFun g := moonshineMatterNearVOA.primaryQuotientLieEquiv (moonshineMatterAut g)
    (moonshineMatterAut_virasoro g) (fun u v => by
      simpa only [NearVOA.smul_def, mode_apply, neg_zero, zero_sub]
        using moonshineMatterAut_field g u v (-1))
theorem monsterLieAut_momentum (g : AutVnat) (α : lorentzianLattice.carrier)
    (u : monsterLieAlgebra) (hu : u ∈ monsterLieMomentum α) :
    monsterLieAut g u ∈ monsterLieMomentum α := by
\end{leancode}
(\leanpath{VOA/Moonshine/MoonshineCNMonsterLie.lean}.)
\begin{dictionary}
\lean{g.toLinearEquiv.rTensor LorentzianFockSpace} is $g\otimes1$.  The two side conditions passed to
\lean{primaryQuotientLieEquiv} are: $g\otimes1$ commutes with every $L_n$ (\lean{moonshineMatterAut_virasoro}, as $g$
fixes $\omega^\natural$) and with all fields (\lean{moonshineMatterAut_field}, as $g$ is a VOA automorphism); here only the
zero product $u_{(0)}v$ is needed, written via the raw coefficient index $-1$ by the convention
\lean{mode f n = f (-n - 1)}.  \lean{→*} is a group homomorphism (with the group structure
\lean{monsterLieAutomorphismGroup} on Lie automorphisms).
\end{dictionary}
\begin{physics}The Monster acts as global symmetry of the string background, hence on its gauge algebra.\end{physics}
\begin{status}No hypotheses; standard axioms (\leanpath{VOA/Moonshine/MoonshineCNAudit10.lean}).\end{status}

\section{Root spaces: from the sector model to the vertex algebra}\label{ch18:sec:roots}

\subsection{The sector equivalence}

Chapter~\ref{ch:17} constructs, for each momentum $\alpha$, a linear isomorphism
\lean{moonshineMatterSectorEquiv α : (Vnat ⊗[ℂ] hypFock) ≃ₗ[ℂ] moonshineMatterMomentum α}: $1_{\Vnat}$ tensored
with the change of polynomial variables from the isotropic coordinates $(a_j,b_j)$ to the complex orthonormal
frame coordinates of the lattice Fock space (\lean{hypLorentzianPolynomialEquiv}: $a_j\mapsto\frac12X_{0,j}-\frac
i2X_{1,j}$, $b_j\mapsto X_{0,j}+iX_{1,j}$), followed by insertion of $e^\alpha$.  The files
\lean{SectorFrame}, \lean{SectorOscillators}, \lean{SectorVirasoroOperators} and
\lean{SectorVirasoroComparison} of this chapter prove that this map intertwines everything relevant.

\begin{theorem}[Sector comparison]\label{ch18:thm:sectorcomparison}
For all $x,y,n,N\in\Z$, writing $\Psi=$ \lean{moonshineMatterSectorMap ![x, y]}:
(a) $\Psi\circ L_n^{\text{sector}}=L_n^{\Vnat\otimes V_{\chXVIIIII}}\circ\Psi$;
(b) $\Psi u$ has weight $N$ iff $u\in H^N_{(x,y)}$;
(c) $\Psi u$ is spurious in $\Vnat\otimes V_{\chXVIIIII}$ iff $u$ is spurious in the sector;
(d) $\Psi u$ is a weight-one primary iff $u\in H^1_{(x,y)}\cap\chXVIIIP$.
\end{theorem}
\begin{leancode}
theorem moonshineMatterSectorMap_virasoro (x y n : ℤ) (u : Vnat ⊗[ℂ] hypFock) :
    moonshineMatterSectorMap ![x, y] (vnatHypMode x y n u) =
      moonshineMatterNearVOA.L n (moonshineMatterSectorMap ![x, y] u) := by
theorem moonshineMatterSector_weight_iff (x y N : ℤ) (u : Vnat ⊗[ℂ] hypFock) :
    moonshineMatterSectorMap ![x, y] u ∈ moonshineMatterNearVOA.grading N ↔
      u ∈ vnatSectorWeight x y N := by
\end{leancode}
(\leanpath{VOA/Moonshine/MoonshineCNSectorVirasoroComparison.lean}.)
\begin{leancode}
theorem moonshineMatterSector_spurious_iff (x y : ℤ) (a : Fin 2)
    (u : Vnat ⊗[ℂ] hypFock) :
    moonshineMatterSectorMap ![x, y] u ∈ moonshineMatterNearVOA.spurious ↔
      u ∈ (vnatStringModule x y a).spurious := by
theorem moonshineMatterSector_primary_iff (x y : ℤ) (a : Fin 2)
    (u : Vnat ⊗[ℂ] hypFock) :
    moonshineMatterSectorMap ![x, y] u ∈ moonshineMatterNearVOA.primaryWeightOne ↔
      u ∈ vnatSectorWeight x y 1 ∧ u ∈ (vnatStringModule x y a).physical := by
\end{leancode}
(\leanpath{VOA/Moonshine/MoonshineCNSectorPhysicalTransfer.lean}.)
\begin{dictionary}
\textbf{What is compared.}  On the left, the explicit sector operators (\S\ref{ch18:sec:sector}); on the
right, the Virasoro modes of the actual vertex algebra $\Vnat\otimes V_{\chXVIIIII}$, whose Lorentzian part on
momentum $\alpha$ is, by \lean{lorentzianSectorVirasoro_apply}, the Heisenberg Sugawara operator plus the
momentum-coupling term $\alpha(n)$ plus $\frac12(\alpha,\alpha)\delta_{n,0}$.  (a) for $n\ge0$ is a computation on generators
(\lean{hypLorentzianPolynomial_virasoro_nonnegative}); negative modes follow from adjointness
(\lean{hypLorentzianPolynomial_virasoro}: ``Every integer mode, with negative modes recovered by the proved
bilinear adjoints'').  (b) is ``total weight $=$ matter weight $+$ level $+xy$''.  (c): the docstring
``Global negative descendants intersect a momentum sector in precisely its own descendants''; the proof uses
the left inverse \lean{moonshineMatterSectorRetract}, which is $\Psi^{-1}$ composed with the momentum projection
and intertwines the $L_n$.  The family $a$ in (c), (d) is arbitrary (the spurious and physical spaces do
not depend on $K$).
\end{dictionary}
\begin{physics}
The momentum-$r$ sector of the compactified string, written in light-cone oscillators, is exactly the
model Fock space: changing from Cartesian to null coordinates of the two light-cone bosons.
\end{physics}
\begin{status}No hypotheses; standard axioms (\leanpath{VOA/Moonshine/MoonshineCNAudit11.lean}).\end{status}

\subsection{The root spaces}

\begin{theorem}[Root spaces of $\chXVIIIm$]\label{ch18:thm:rootspaces}
For every $(x,y)$ and family $a$, the map $\chXVIIIP^1_{(x,y)}\to\chXVIIIm_{(x,y)}$, $u\mapsto\overline{\Psi u}$, is surjective with
kernel the radical; so $\chXVIIIP^1_{(x,y)}/\mathrm{rad}\cong\chXVIIIm_{(x,y)}$, $\Aut(\Vnat)$-equivariantly.  Consequently, for
$(x,y)\neq0$,
\begin{equation*}
  \chXVIIIm_{(x,y)}\ \cong\ \Vnat_{1-xy}\qquad\text{as }\Aut(\Vnat)\text{-modules}.
\end{equation*}
\end{theorem}
\begin{leancode}
theorem monsterLieSectorProjection_surjective (x y : ℤ) (a : Fin 2) :
    Function.Surjective (monsterLieSectorProjection x y a) := by
theorem monsterLieSectorProjection_ker (x y : ℤ) (a : Fin 2) :
    LinearMap.ker (monsterLieSectorProjection x y a) =
      (vnatStringModule x y a).weightRadical (vnatStringModule x y a).physical 1 := by
def monsterLieSectorQuotientEquiv (x y : ℤ) (a : Fin 2) :
    (vnatStringModule x y a).weightRadicalQuotient (vnatStringModule x y a).physical 1 ≃ₗ[ℂ]
      monsterLieMomentum ![x, y] :=
def monsterLieRootSpaceEquiv (x y : ℤ) (h : x ≠ 0 ∨ y ≠ 0) :
    monsterLieMomentum ![x, y] ≃ₗ[ℂ] VnatVOA.grading (1 - x * y) :=
  (monsterLieSectorQuotientEquiv x y (stringSectorFamily x y h)).symm.trans
    (stringSectorNoGhostEquiv x y h)
\end{leancode}
(\leanpath{VOA/Moonshine/MoonshineCNRootSpaces.lean}.)
\begin{leancode}
theorem monsterLieRootSpaceEquiv_equivariant (x y : ℤ) (h : x ≠ 0 ∨ y ≠ 0)
    (g : AutVnat) (u : monsterLieMomentum ![x, y]) :
    monsterLieRootSpaceEquiv x y h
      ⟨monsterLieAut g u, monsterLieAut_momentum g _ u u.property⟩ =
      GriessAut.autRestrict g (1 - x * y) (monsterLieRootSpaceEquiv x y h u) := by
\end{leancode}
(\leanpath{VOA/Moonshine/MoonshineCNRootSpaceAction.lean}.)
\begin{dictionary}
\textbf{Decoding.}  \lean{monsterLieSectorProjection x y a} sends a weight-one physical vector $u$ of the
sector to the class of $\Psi u$ in $\chXVIIIm_{(x,y)}$ (it is a primary by Theorem~\ref{ch18:thm:sectorcomparison}(d)).
The kernel statement uses (c) and Theorem~\ref{ch18:thm:radicals}: $\Psi u$ is spurious iff $u\in\chXVIIIP^1\cap\chXVIIIS=\mathrm{rad}$
(\lean{vnatPhysicalRadical_mem_iff}).  \lean{monsterLieRootSpaceEquiv} composes the inverse with the no-ghost
isomorphism.  The equivariance is with respect to the \emph{actual} actions: \lean{monsterLieAut g} on $\chXVIIIm$
and $g|_{\Vnat_{1-xy}}$, and the docstring says ``with no additional hypotheses''.

\textbf{Weaker than Borcherds' statement in one respect.}  The isomorphism is linear and equivariant; an
\emph{isometry} statement for it is not made at the level of $\chXVIIIm$, because the natural form on $\chXVIIIm$ pairs
$\chXVIIIm_r$ with $\chXVIIIm_{-r}$, not $\chXVIIIm_r$ with itself (\S\ref{ch18:sec:form}).  The isometry is available at the
sector level (Theorem~\ref{ch18:thm:noghost}).

\textbf{Momentum convention.}  Repository coordinates, norm $2xy$.  The Borcherds form is
\lean{monsterLieBorcherdsRootEquiv} (\S\ref{ch18:sec:monsterbd}).
\end{dictionary}
\begin{physics}
The gauge bosons of momentum $r\neq0$ are labelled by the internal states of level $1-r^2/2$; the Monster
acts on them as on $\Vnat_{1-r^2/2}$.  The real roots ($r^2=2$) give $\Vnat_0=\mathbb C$, one state; the
imaginary roots of norm $-2$ give $\Vnat_2$, the $196884$-dimensional Griess algebra.
\end{physics}
\begin{status}No hypotheses beyond $(x,y)\neq0$; standard axioms (\leanpath{VOA/Moonshine/MoonshineCNAudit11.lean}).\end{status}

\begin{theorem}[The zero root space]\label{ch18:thm:zerospace}
$\chXVIIIm_0\cong\mathbb C^2$, the isomorphism sending the class of $c_0\,\mathbb 1\otimes a_1+c_1\,\mathbb 1\otimes b_1$ (transported to
the vertex algebra) to $(c_0,c_1)$.
\end{theorem}
\begin{leancode}
def monsterLieZeroMomentumEquiv : monsterLieMomentum 0 ≃ₗ[ℂ] (Fin 2 → ℂ) :=
  monsterLieZeroMomentumTransport.symm.trans
    ((monsterLieSectorQuotientEquiv 0 0 0).symm.trans vnatZeroNoGhostEquiv)
\end{leancode}
(\leanpath{VOA/Moonshine/MoonshineCNRootSpaces.lean}; \lean{monsterLieZeroMomentumTransport} only rewrites
\lean{![0, 0]} as \lean{0}.)
\begin{physics}The Cartan subalgebra is spanned by the two light-cone momentum currents.\end{physics}
\begin{status}No hypotheses; standard axioms (\leanpath{VOA/Moonshine/MoonshineCNAudit11.lean}).\end{status}

\section{The invariant form on the Monster Lie algebra}\label{ch18:sec:form}

\subsection{Two forms and momentum reversal}

The vertex algebra $\Vnat\otimes V_{\chXVIIIII}$ carries the form \lean{moonshineMatterForm := natForm.tmul lorentzianBilinearForm}
(Chapter~\ref{ch:17}), whose lattice factor pairs $e^\alpha$ with $e^{-\alpha}$ (an ``FHL'' type invariant form).  The sector
form $B^\natural\otimes B_F$ pairs a momentum sector with itself.  These cannot be matched directly:

\begin{theorem}[No same-momentum isometry]\label{ch18:thm:nosame}
For $\alpha\neq0$: the matter form vanishes on momentum $\alpha$ against momentum $\beta$ unless $\alpha+\beta=0$; in particular
every nonzero momentum sector is totally isotropic; and there is no map $F$ from the sector to the momentum-$\alpha$
part and no nonzero scalar $c$ with $\langle Fu,Fv\rangle=c\,B(u,v)$ for all $u,v$.
\end{theorem}
\begin{leancode}
theorem moonshineMatterForm_momentum_orthogonal (α β : lorentzianLattice.carrier)
    (hαβ : α + β ≠ 0) (u v : moonshineMatterSpace)
    (hu : u ∈ moonshineMatterMomentum α) (hv : v ∈ moonshineMatterMomentum β) :
    moonshineMatterForm u v = 0 := by
theorem moonshineMatterSector_no_same_momentum_isometry (α : lorentzianLattice.carrier)
    (hα : α ≠ 0) (c : ℂ) (hc : c ≠ 0)
    (F : (Vnat ⊗[ℂ] hypFock) → moonshineMatterMomentum α) :
    ¬ (∀ u v, moonshineMatterInvariantForm.form (F u) (F v) = c * vnatHypForm u v) := by
\end{leancode}
(\leanpath{VOA/Moonshine/MoonshineCNSectorFormConvention.lean}.)
\begin{warning}
\textbf{A natural guess that fails, recorded in Lean.}  One is tempted to identify the no-ghost sector form with the
vertex-algebra form on the same momentum; \lean{moonshineMatterSector_no_same_momentum_isometry} proves this is
impossible, even for non-linear $F$ (the docstring: ``This rules out the literal same-momentum isometry even
without imposing linearity'').  Reason: the sector form is nondegenerate, the target form is identically zero
on a nonzero momentum sector.  The fix is \emph{momentum reversal}.
\end{warning}
\begin{physics}
The BPZ (Belavin--Polyakov--Zamolodchikov) two-point pairing $\langle\phi_r|\phi_s\rangle$ is nonzero only when momentum is conserved, $r+s=0$; a state of momentum
$r\neq0$ has zero BPZ overlap with every state of the same momentum.  The physicist's Hermitian norm of a state needs a
conjugation that reverses momentum.  (All forms in the Lean statements are complex \emph{bilinear}; ``Hermitian'' in these
boxes refers to the physical inner product, which is such a bilinear form composed with a real structure, and is never
used mathematically.)
\end{physics}
\begin{status}No hypotheses beyond $\alpha\neq0$, $c\neq0$; standard axioms (\leanpath{VOA/Moonshine/MoonshineCNAudit9.lean}).\end{status}

\begin{theorem}[Comparison through momentum reversal]\label{ch18:thm:reversal}
Let $\theta$ be momentum reversal on the lattice factor (\lean{moonshineMatterReversal}, $1\otimes\theta_L$ with $\theta_L$
the lift of $-1$ of Chapter~\ref{ch:17}).  For every momentum $\alpha$ (including $0$) and all $u,v$ in the sector,
\begin{equation*}
  \langle\theta\,\Psi_\alpha u,\ \Psi_\alpha v\rangle=\epsilon(\alpha)\,B(u,v),\qquad \epsilon(\alpha)=\text{\lean{chSign lorentzianLattice α}}\in\{\pm1\},
\end{equation*}
and the reversed pairing $(u,v)\mapsto\langle\theta u,v\rangle$ is nondegenerate on each momentum sector.
\end{theorem}
\begin{leancode}
theorem moonshineMatterSector_form (α : lorentzianLattice.carrier)
    (u v : Vnat ⊗[ℂ] hypFock) :
    moonshineMatterForm
      (moonshineMatterReversal (moonshineMatterSectorEquiv α u))
      (moonshineMatterSectorEquiv α v) =
      chSign lorentzianLattice α * vnatHypForm u v := by
theorem moonshineMatterSector_form_nondegenerate (α : lorentzianLattice.carrier)
    (u : moonshineMatterMomentum α)
    (hu : ∀ v : moonshineMatterMomentum α,
      moonshineMatterForm (moonshineMatterReversal u) v = 0) : u = 0 := by
\end{leancode}
(\leanpath{VOA/Moonshine/MoonshineCNSectorFormComparison.lean}.)
\begin{dictionary}
\lean{chSign lorentzianLattice α} is the cocycle sign of the lattice form (\leanpath{VOA/Moonshine/NondegLattice.lean}): in a lattice
vertex algebra the vertex operators $e^\alpha$ carry signs $\varepsilon(\alpha,\beta)=\pm1$ (a 2-cocycle, Chapter~\ref{ch:17}), and the
invariant form pairs $e^\alpha$ with $e^{-\alpha}$ with a value $\pm1$ determined by these signs; that value is $\epsilon(\alpha)$;
\lean{moonshineMatterForm_opposite_ground} computes it on ground states.  The comparison is ``with the charge
sign retained'' (file docstring).  The oscillator part uses \lean{hypLorentzianPolynomial_form}: reversal turns the
negative-sign orthonormal Fock form into the hyperbolic monomial form \lean{hypForm}.
\end{dictionary}
\begin{physics}
The BPZ (FHL) inner product pairs in-states of momentum $r$ with out-states of momentum $-r$; the Hermitian
(contravariant) inner product is obtained by composing with a conjugation that reverses momentum, here
$\theta$.  The two forms differ exactly by $\theta$ and a cocycle sign.
\end{physics}
\begin{status}No hypotheses; standard axioms (\leanpath{VOA/Moonshine/MoonshineCNAudit9.lean},
\leanpath{VOA/Moonshine/MoonshineCNAudit11.lean}).\end{status}

\begin{theorem}[The physical radical in each sector]\label{ch18:thm:rootradical}
For a weight-one primary $u$ of momentum $(x,y)$ (any $(x,y)$, including $0$): $u$ is spurious iff
$\langle\theta u,v\rangle=0$ for every weight-one primary $v$ of momentum $(x,y)$.
\end{theorem}
\begin{leancode}
theorem moonshineMatterSector_physical_radical (x y : ℤ) (u : moonshineMatterSpace)
    (hup : u ∈ moonshineMatterNearVOA.primaryWeightOne)
    (hum : u ∈ moonshineMatterMomentum ![x, y]) :
    u ∈ moonshineMatterNearVOA.spurious ↔
      ∀ v : moonshineMatterSpace, v ∈ moonshineMatterNearVOA.primaryWeightOne →
        v ∈ moonshineMatterMomentum ![x, y] →
          moonshineMatterForm (moonshineMatterReversal u) v = 0 := by
\end{leancode}
(\leanpath{VOA/Moonshine/MoonshineCNRootRadical.lean}.)
\begin{physics}A physical state is null iff it has zero (Hermitian) overlap with every physical state of the same momentum.\end{physics}
\paragraph{Proof on paper.}  Transport through $\Psi$ (Theorem~\ref{ch18:thm:sectorcomparison}) to the sector; there, by
Theorem~\ref{ch18:thm:reversal}, the reversed pairing is $\pm B$, and by Theorem~\ref{ch18:thm:radicals} the radical of $B$ on
$\chXVIIIP^1$ is $\chXVIIIP^1\cap\chXVIIIS$. \qed
\begin{status}No hypotheses; standard axioms (\leanpath{VOA/Moonshine/MoonshineCNAudit11.lean}).\end{status}

\subsection{\texorpdfstring{The form on $\chXVIIIm$}{The form on m}}

\begin{definition}[The Monster Lie algebra form]\label{ch18:def:monsterform}
$(\bar u,\bar v):=-\langle u,v\rangle$ for weight-one primaries $u,v$ of $\Vnat\otimes V_{\chXVIIIII}$.
\end{definition}
\begin{leancode}
def monsterPhysicalBilinearForm : LinearMap.BilinForm ℂ moonshineMatterNearVOA.primaryWeightOne :=
  -(moonshineMatterForm.restrict moonshineMatterNearVOA.primaryWeightOne)
def monsterLieBilinearForm : LinearMap.BilinForm ℂ monsterLieAlgebra :=
  monsterPhysicalBilinearForm.liftQ₂ moonshineMatterNearVOA.primarySpurious
    moonshineMatterNearVOA.primarySpurious
    monsterPhysicalBilinearForm_left_null monsterPhysicalBilinearForm_right_null
theorem monsterLieBilinearForm_mk (u v : moonshineMatterNearVOA.primaryWeightOne) :
    monsterLieBilinearForm (monsterLieProjection u) (monsterLieProjection v) =
      -moonshineMatterForm u v := rfl
theorem monsterLieBilinearForm_symmetric (u v : monsterLieAlgebra) :
    monsterLieBilinearForm u v = monsterLieBilinearForm v u := by
theorem monsterLieBilinearForm_aut (g : AutVnat) (u v : monsterLieAlgebra) :
    monsterLieBilinearForm (monsterLieAut g u) (monsterLieAut g v) =
      monsterLieBilinearForm u v := by
\end{leancode}
(\leanpath{VOA/Moonshine/MoonshineCNMonsterForm.lean}.)
\begin{dictionary}
\lean{liftQ₂} descends a bilinear form to the quotient in both arguments, given that the subspace is in the left
and right kernels: \lean{monsterPhysicalBilinearForm_left_null}, \lean{_right_null}, which rest on
\lean{moonshineMatterForm_spurious_primary} ($\langle L_{-n}y,v\rangle=\langle y,L_nv\rangle=0$ for a primary $v$).  The
\emph{sign} $-1$ is a normalisation, chosen so that the form on the Cartan subalgebra is the lattice form
(Theorem~\ref{ch18:thm:zeroform}; the file docstring: ``The sign is chosen for the lattice normalization on the Cartan
space'').
\end{dictionary}
\begin{physics}The Killing-type form of the gauge algebra, given by the two-point function of the currents.\end{physics}
\begin{status}No hypotheses; standard axioms (\leanpath{VOA/Moonshine/MoonshineCNAudit11.lean}).\end{status}

\begin{theorem}[Root pairing, nondegeneracy, and the radical]\label{ch18:thm:primaryradical}
(a) $(\chXVIIIm_\alpha,\chXVIIIm_\beta)=0$ unless $\alpha+\beta=0$.  (b) For every $\alpha$ (including $0$), the pairing
$\chXVIIIm_\alpha\times\chXVIIIm_{-\alpha}\to\mathbb C$ is nondegenerate on the left.  (c) The form on $\chXVIIIm$ is nondegenerate.  (d) A weight-one
primary of $\Vnat\otimes V_{\chXVIIIII}$ is spurious iff it is orthogonal to all weight-one primaries: Borcherds' definition of
$\chXVIIIm$ ($P^1$ modulo the radical) coincides with Definition~\ref{ch18:def:monsterlie}.
\end{theorem}
\begin{leancode}
theorem monsterLieBilinearForm_momentum_orthogonal (α β : lorentzianLattice.carrier)
    (hαβ : α + β ≠ 0) (u v : monsterLieAlgebra)
    (hu : u ∈ monsterLieMomentum α) (hv : v ∈ monsterLieMomentum β) :
    monsterLieBilinearForm u v = 0 := by
theorem monsterLieBilinearForm_pairing_nondegenerate (α : lorentzianLattice.carrier)
    (u : monsterLieAlgebra) (hu : u ∈ monsterLieMomentum α)
    (hB : ∀ v : monsterLieAlgebra, v ∈ monsterLieMomentum (-α) →
      monsterLieBilinearForm u v = 0) : u = 0 := by
\end{leancode}
(\leanpath{VOA/Moonshine/MoonshineCNMonsterPairing.lean}.)
\begin{leancode}
theorem monsterLieBilinearForm_nondegenerate : monsterLieBilinearForm.Nondegenerate := by
theorem moonshineMatterPrimary_radical (u : moonshineMatterNearVOA.primaryWeightOne) :
    u ∈ moonshineMatterNearVOA.primarySpurious ↔
      ∀ v : moonshineMatterNearVOA.primaryWeightOne, moonshineMatterForm u v = 0 := by
\end{leancode}
(\leanpath{VOA/Moonshine/MoonshineCNMonsterRadical.lean}.)
\begin{dictionary}
\lean{.Nondegenerate} for a bilinear form is Mathlib's \lean{LinearMap.Nondegenerate}: separating on the left
\emph{and} on the right.  (d) is the docstring's ``The original vertex-space form has exactly the physical spurious
radical''.  All four are unconditional.
\end{dictionary}
\begin{physics}Every gauge boson of momentum $r$ has a partner of momentum $-r$ with nonzero two-point function.\end{physics}
\paragraph{Proof on paper.}  (a) from Theorem~\ref{ch18:thm:nosame}.  (b): let $p\in P^1_\alpha$ pair to zero with
$\chXVIIIm_{-\alpha}$.  For any $v\in P^1_\alpha$, $\theta v\in P^1_{-\alpha}$ (\lean{moonshineMatterReversal_primary}, momentum reversed), so
$0=\langle p,\theta v\rangle=\langle\theta p,v\rangle$ ($\theta$ is self-adjoint); by Theorem~\ref{ch18:thm:rootradical}, $p$ is spurious, i.e.\
$\bar p=0$.  (c): decompose by momentum (Theorem~\ref{ch18:thm:grading}), apply (a), (b).  (d): (c) lifted to $P^1$. \qed
\begin{status}No hypotheses; standard axioms (\leanpath{VOA/Moonshine/MoonshineCNAudit11.lean}).\end{status}

\begin{remark}[Three descriptions of $\chXVIIIm$]\label{ch18:rem:threedescriptions}
Let $P^0$ be the weight-zero vectors killed by all $L_n$, $n\ge1$, and $R:=\{u\in P^1:\langle u,P^1\rangle=0\}$ the radical.  The literature
uses three quotients of $P^1$:
\begin{equation*}
  \text{(a) } P^1/(P^1\cap\mathrm{Sp})\ \text{(this chapter)},\qquad
  \text{(b) } P^1/R\ \text{(Borcherds)},\qquad
  \text{(c) } P^1/(L_{-1}P^0+R)\ \text{(Frenkel; Borcherds' ``$P^1/L_{-1}P^0$ modulo the radical'')}.
\end{equation*}
They coincide.  First, $L_{-1}P^0\subseteq P^1\cap\mathrm{Sp}$: for $p\in P^0$, $L_{-1}p$ has weight one and is spurious, and
$L_1L_{-1}p=L_{-1}L_1p+2L_0p=0$, $L_2L_{-1}p=L_{-1}L_2p+3L_1p=0$, so $L_{-1}p$ is primary (Theorem~\ref{ch18:thm:twogenerators}(b) in a
string module, or the same two-generator argument in $V$).  Second, $P^1\cap\mathrm{Sp}=R$ by Theorem~\ref{ch18:thm:primaryradical}(d).  Hence
$L_{-1}P^0+R=R=P^1\cap\mathrm{Sp}$ and (a) $=$ (b) $=$ (c).  At weight one, spurious means
$\mathrm{Sp}\cap V_1=L_{-1}V_0+L_{-2}V_{-1}$ (Theorem~\ref{ch18:thm:twogenerators}(a)); the $L_{-2}V_{-1}$ part is why (c) needs ``$+R$'':
at $c=26$ the null states $\widetilde L_{-2}y$ of \S\ref{ch18:sec:critical} are not of the form $L_{-1}P^0$ in general.  Only (a) and
(b) are Lean statements (the definition, and \lean{moonshineMatterPrimary_radical}); the inclusion $L_{-1}P^0\subseteq P^1\cap\mathrm{Sp}$
is the elementary computation just given and is not separately formalized.
\end{remark}

\begin{theorem}[The form on the Cartan subalgebra is the lattice form]\label{ch18:thm:zeroform}
For $u,v\in\chXVIIIm_0$ with coordinates $c=(c_0,c_1)$, $d=(d_0,d_1)$ under Theorem~\ref{ch18:thm:zerospace}:
$(u,v)=c_0d_1+c_1d_0$.  $\Aut(\Vnat)$ acts trivially on $\chXVIIIm_0$.
\end{theorem}
\begin{leancode}
theorem monsterLieBilinearForm_zeroMomentum (u v : monsterLieMomentum 0) :
    monsterLieBilinearForm u v =
      monsterLieZeroMomentumEquiv u 0 * monsterLieZeroMomentumEquiv v 1 +
      monsterLieZeroMomentumEquiv u 1 * monsterLieZeroMomentumEquiv v 0 := by
theorem monsterLieZeroMomentum_aut_trivial (g : AutVnat) (u : monsterLieMomentum 0) :
    monsterLieAut g u = (u : monsterLieAlgebra) := by
\end{leancode}
(\leanpath{VOA/Moonshine/MoonshineCNMonsterZeroForm.lean}.)
\begin{dictionary}
At zero momentum $\theta$ acts by $-1$ on the weight-one states $\mathbb 1\otimes h(-1)e^0$
(\lean{moonshineMatterReversal_zeroWeight}), which is why the sign $-1$ in Definition~\ref{ch18:def:monsterform} makes the
form equal to the sector form there (\lean{moonshineMatterSector_zeroWeight_form}).  In coordinates the Gram matrix is
$\left(\begin{smallmatrix}0&1\\1&0\end{smallmatrix}\right)$, the form \eqref{ch18:eq:repform} of $\chXVIIIII$ under
$(c_0,c_1)\mapsto(c_1,c_0)$ (see the next theorem for why the coordinates are swapped).
\end{dictionary}
\begin{physics}The two-point function of the momentum currents $\partial X^\pm$ is the spacetime metric.\end{physics}
\begin{status}No hypotheses; standard axioms (\leanpath{VOA/Moonshine/MoonshineCNAudit11.lean}).\end{status}

\begin{theorem}[The Cartan action]\label{ch18:thm:cartanaction}
If $h\in\chXVIIIm_0$ has coordinates $(c_0,c_1)$ and $x\in\chXVIIIm_\beta$, then $[h,x]=(c_0\beta_0+c_1\beta_1)\,x$.
\end{theorem}
\begin{leancode}
theorem monsterLie_zeroMomentum_bracket (h : monsterLieMomentum 0)
    (β : lorentzianLattice.carrier) (x : monsterLieAlgebra) (hx : x ∈ monsterLieMomentum β) :
    ⁅(h : monsterLieAlgebra), x⁆ =
      (monsterLieZeroMomentumEquiv h 0 * (β 0 : ℂ) + monsterLieZeroMomentumEquiv h 1 * (β 1 : ℂ)) •
        x := by
\end{leancode}
(\leanpath{VOA/Moonshine/MoonshineCNMonsterCartanAction.lean}.)
\begin{dictionary}
\lean{β 0}, \lean{β 1} are the two coordinates $(x,y)$ of the momentum.  The coordinate $c_0$ belongs to
$\mathbb 1\otimes a_1$, the current of the isotropic vector $w_0=(0,1)$, and $c_1$ to $\mathbb 1\otimes b_1$, the current of
$w_1=(1,0)$.  So $h$ corresponds to the lattice vector $\lambda(h)=c_0(0,1)+c_1(1,0)=(c_1,c_0)$, and
$(\lambda(h),(\beta_0,\beta_1))=c_1\beta_1+c_0\beta_0$ by \eqref{ch18:eq:repform}: the statement says $[h,x]=(\lambda(h),\beta)x$, the
expected Cartan action.  The proof evaluates the zero modes of the actual hyperbolic Heisenberg currents
(\lean{lorentzianHyperbolicCurrent_zeroMode}: ``multiplies momentum \lean{β} by \lean{β a}'').
\end{dictionary}
\begin{physics}The zero mode of the current $\partial X^\mu$ is the momentum operator $p^\mu$.\end{physics}
\begin{status}No hypotheses; standard axioms (\leanpath{VOA/Moonshine/MoonshineCNAudit17.lean}).\end{status}

\subsection{Invariance}

\begin{theorem}[Invariance of the form]\label{ch18:thm:invariance}
$([u,v],w)=(u,[v,w])$ for all $u,v,w\in\chXVIIIm$.
\end{theorem}
\begin{leancode}
theorem monsterLieBilinearForm_invariant (u v w : monsterLieAlgebra) :
    monsterLieBilinearForm ⁅u, v⁆ w = monsterLieBilinearForm u ⁅v, w⁆ :=
  monsterLieBilinearForm_invariant_of_oppositeIterate
    moonshineMatterNearVOA_oppositeIterateIdentity u v w
theorem moonshineMatterForm_primary_zeroMode_skew
    (a : moonshineMatterNearVOA.primaryWeightOne) (b c : moonshineMatterSpace) :
    moonshineMatterForm (moonshineMatterNearVOA.smul a 0 b) c =
      -moonshineMatterForm b (moonshineMatterNearVOA.smul a 0 c) := by
\end{leancode}
(\leanpath{VOA/Moonshine/MoonshineCNMonsterLieInvariance.lean}.)
\begin{leancode}
theorem monsterLieBilinearForm_invariant_of_primaryAdjoint
    (hAdj : ∀ a b c : moonshineMatterNearVOA.primaryWeightOne,
      moonshineMatterForm (moonshineMatterNearVOA.smul a 0 b) c =
        -moonshineMatterForm b (moonshineMatterNearVOA.smul a 0 c))
    (u v w : monsterLieAlgebra) :
    monsterLieBilinearForm ⁅u, v⁆ w = monsterLieBilinearForm u ⁅v, w⁆ := by
\end{leancode}
(\leanpath{VOA/Moonshine/MoonshineCNMonsterPairing.lean}.)
\begin{dictionary}
\textbf{No hypotheses.}  The headline is a one-line composition: the general reduction (invariance follows
from skew-adjointness of zero modes of weight-one primaries) applied to the proved \emph{opposite iterate
identity} of the matter near-VOA, \lean{moonshineMatterNearVOA_oppositeIterateIdentity}, which is
\lean{NearVOA.oppositeIterateIdentity_of_locallyNilpotent} (Chapter~\ref{ch:17}) fed with the local nilpotence of
$L_1$ (\lean{moonshineMatterVirasoroOne_locallyNilpotent}).  Its consequence
\lean{moonshineMatterOppositeModes_adjoint} says that every field of $\Vnat\otimes V_{\chXVIIIII}$ is adjoint to its FHL
``opposite'' field $Y(e^{zL_1}(-z^{-2})^{L_0}u,z^{-1})$.  (An \emph{FHL-invariant} form is one for which $\langle Y(u,z)a,b\rangle=
\langle a,Y(e^{zL_1}(-z^{-2})^{L_0}u,z^{-1})b\rangle$; the right-hand field is the ``opposite'' field of $u$, and the identity says that
the adjoint of $Y(u,z)$ is its opposite field.)  For a weight-one primary the opposite of $u_{(0)}$ is $-u_{(0)}$
(\lean{moonshineMatterOppositeModes_primary}).

\textbf{Structure of the proof in Lean.}  The reduction \lean{monsterLieBilinearForm_invariant_of_primaryAdjoint} takes the
skew-adjointness of zero modes as an explicit hypothesis; \lean{moonshineMatterForm_primary_zeroMode_skew} proves that hypothesis,
so the headline has none.
\end{dictionary}
\begin{physics}
Invariance of the current two-point function under the gauge algebra: $\langle[J_u,J_v]J_w\rangle=\langle J_u[J_v,J_w]\rangle$, a
consequence of the three-point function of currents being totally antisymmetric.
\end{physics}
\paragraph{Proof on paper.}  For a weight-one primary $a$, FHL invariance of the form at $z^0$ gives
$\langle a_{(0)}b,c\rangle=-\langle b,a_{(0)}c\rangle$.  Then, writing $u=\bar a$, $v=\bar b$, $w=\bar c$:
$([u,v],w)=-([v,u],w)=\langle b_{(0)}a,c\rangle=-\langle a,b_{(0)}c\rangle=(u,[v,w])$ (the overall sign of
Definition~\ref{ch18:def:monsterform} cancels). \qed
\begin{status}No hypotheses; standard axioms (\leanpath{VOA/Moonshine/MoonshineCNAudit16.lean}).\end{status}

\begin{warning}
\textbf{The VOA cutoff does not work for near-VOAs.}  The library's FHL machinery for vertex operator algebras truncates
$e^{zL_1}$ at the weight (\lean{wt.toNat + 1}); for $\Vnat\otimes V_{\chXVIIIII}$ this fails, because $L_1$ is not nilpotent of
order $\le\mathrm{wt}+1$ on states of negative or zero weight with nonzero momentum (the Lean certificate is
\lean{lorentzianVirasoroOne_weightZero_not_zero}, Chapter~\ref{ch:17}).  The correct cutoff uses the matter/oscillator
level, i.e.\ weight minus the half-norm of the momentum (\lean{moonshineMatterOppositeModes_charge_cutoff}: ``The actual
FHL transform uses matter/oscillator level, not total weight'').
\end{warning}

\section{Abstract Borcherds data}\label{ch18:sec:bd}

The rest of the chapter is pure Lie theory.  It isolates exactly the properties of $\chXVIIIm$ that Borcherds' argument
uses, as a structure \lean{BorcherdsData}, and proves the structure theorem from them.  This separation has two benefits: the Lie theory does not depend on any
vertex algebra, and the one input that is \emph{not} a property of $\chXVIIIm$ visible from its construction
(generation by simple root spaces) is an explicit field.

\subsection{What a Borcherds algebra is, and what is (not) proved here}\label{ch18:sec:bkm}

We recall the standard notions (Borcherds 1988; Kac, \emph{Infinite dimensional Lie algebras}, Chapter~11), to fix
vocabulary.  A \emph{generalized Cartan matrix} in Borcherds' sense is a real matrix $A=(a_{ij})_{i,j\in I}$, $I$ countable,
which is symmetric (in our case) and satisfies: $a_{ii}=2$ or $a_{ii}\le0$; $a_{ij}\le0$ for $i\neq j$; and $2a_{ij}/a_{ii}\in\Z$ when
$a_{ii}=2$.  An index $i$ with $a_{ii}=2$ is \emph{real}, one with $a_{ii}\le0$ \emph{imaginary}.  The \emph{generalized Kac--Moody}
(or \emph{Borcherds--Kac--Moody}) algebra $\chXVIIIg(A)$ is generated by $e_i,f_i,h_i$ ($i\in I$) subject to
\begin{equation*}
  [h_i,h_j]=0,\quad [h_i,e_j]=a_{ij}e_j,\quad [h_i,f_j]=-a_{ij}f_j,\quad [e_i,f_j]=\delta_{ij}h_i,
\end{equation*}
the \emph{Serre relations} $(\mathrm{ad}\,e_i)^{1-2a_{ij}/a_{ii}}e_j=0=(\mathrm{ad}\,f_i)^{1-2a_{ij}/a_{ii}}f_j$ for real $i$ and $j\ne i$, and
$[e_i,e_j]=0=[f_i,f_j]$ whenever $a_{ij}=0$; one may also allow several simple roots with the same image in the root lattice
(``multiplicities''), as the Monster Lie algebra requires.  Its \emph{roots} are the nonzero degrees of the grading by the root
lattice; a root $\gamma$ is \emph{real} if $(\gamma,\gamma)>0$ (it is then Weyl-conjugate to a real simple root) and \emph{imaginary} if
$(\gamma,\gamma)\le0$.  A \emph{Weyl vector} is a $\delta$ with $(\delta,\alpha_i)=\frac12(\alpha_i,\alpha_i)$ for all simple roots.

\paragraph{The Monster Lie algebra in this language} (Borcherds 1992).  The root lattice is $\chXVIIIII$ in Borcherds coordinates
with $(\cdot,\cdot)_{\mathrm{Bor}}$; the simple roots are $\alpha_r=(1,-1)$ (real, norm $2$) and $\alpha_j=(1,j)$, $j\ge1$ (imaginary, norm
$-2j$), the latter with multiplicity $c(j)=\dim\Vnat_{1+j}$.  The Cartan matrix entries are
\begin{equation*}
  (\alpha_r,\alpha_r)_{\mathrm{Bor}}=2,\qquad (\alpha_r,\alpha_j)_{\mathrm{Bor}}=1-j\le0,\qquad (\alpha_i,\alpha_j)_{\mathrm{Bor}}=-(i+j)<0,
\end{equation*}
so the conditions above hold, and the Serre relation between the real simple root and $\alpha_j$ reads
$(\mathrm{ad}\,e)^{1-(1-j)}v=(\mathrm{ad}\,e)^jv=0$ for $v\in\chXVIIIm_{\alpha_j}$: this is exactly Theorem~\ref{ch18:thm:sl2}(b) below.  The real roots of $\chXVIIIm$
are $\pm\alpha_r$ (its Weyl group is $\{1,s\}$ with $s(m,n)=(n,m)$), all other roots $(m,n)$, $mn\ge1$, are imaginary, and
$\delta=(1,0)$ is a Weyl vector.

\paragraph{What is and is not proved.}  Borcherds proves that $\chXVIIIm$ \emph{is} the generalized Kac--Moody algebra of this Cartan
matrix, using his characterisation theorem (which needs a positive-definite contravariant form) and the denominator identity.
\textbf{This chapter does not prove that.}  The structure \lean{BorcherdsData} below is a weaker, purpose-built list of
properties: a $\Z^2$-grading, a Cartan subalgebra acting by the form, an invariant form pairing opposite root spaces, the root
support, one-dimensional real root spaces, and generation by the simple root spaces.  It does not assert a presentation by
generators and relations, and no positivity.  From it we prove exactly what Chapter~\ref{ch:19} needs: the resolution, $T(U)\cong\mathcal
\mathcal U(\chXVIIIu^+)$ and the freeness of $\chXVIIIu^+$.  ``Borcherds--Kac--Moody structure'' in the chapter title refers to these statements.

\subsection{Rank-two root arithmetic}

All roots are in Borcherds coordinates $\gamma=(m,n)\in\Z^2$ with the form $(\xi,\mu)_{\mathrm{Bor}}=-(\xi_1\mu_2+\mu_1\xi_2)$.

\begin{definition}[Roots]\label{ch18:def:roots}
$\alpha_r=(1,-1)$; $\alpha_j=(1,j)$; the positive cone $C=\{(m,n):m,n\ge1\}$; positive roots $\Phi^+=\{\alpha_r\}\cup C$;
\emph{root support}: $\gamma=0$ or $\pm\gamma\in\Phi^+$; simple roots $\alpha_r$ and $\alpha_j$ ($j\ge1$); height $m+n$; the
Casimir scalar $Q(\gamma)=(\gamma,\gamma)_{\mathrm{Bor}}-2(\delta,\gamma)_{\mathrm{Bor}}$ with Weyl vector $\delta=(1,0)$.
\end{definition}
\begin{leancode}
def borcherdsRealRoot : ℤ × ℤ := (1, -1)
def borcherdsImaginaryRoot (j : ℤ) : ℤ × ℤ := (1, j)
def borcherdsPositiveCone (γ : ℤ × ℤ) : Prop := 1 ≤ γ.1 ∧ 1 ≤ γ.2
def borcherdsPositiveRoot (γ : ℤ × ℤ) : Prop :=
  γ = borcherdsRealRoot ∨ borcherdsPositiveCone γ
def borcherdsRootSupport (γ : ℤ × ℤ) : Prop :=
  γ = 0 ∨ borcherdsPositiveRoot γ ∨ borcherdsPositiveRoot (-γ)
def borcherdsSimpleRoot (γ : ℤ × ℤ) : Prop :=
  γ = borcherdsRealRoot ∨ ∃ j : ℤ, 1 ≤ j ∧ γ = borcherdsImaginaryRoot j
def borcherdsRootHeight (γ : ℤ × ℤ) : ℤ := γ.1 + γ.2
def borcherdsWeightVector (γ : ℤ × ℤ) : ℂ × ℂ := (γ.1, γ.2)
def borcherdsCartanPairing (ξ μ : ℂ × ℂ) : ℂ := -(ξ.1 * μ.2 + μ.1 * ξ.2)
def borcherdsCasimirScalar (γ : ℤ × ℤ) : ℂ :=
  borcherdsCartanPairing (borcherdsWeightVector γ) (borcherdsWeightVector γ) -
    2 * borcherdsCartanPairing (1, 0) (borcherdsWeightVector γ)
theorem borcherdsCasimirScalar_eq (γ : ℤ × ℤ) :
    borcherdsCasimirScalar γ = -2 * (γ.2 : ℂ) * ((γ.1 : ℂ) - 1) := by
theorem borcherdsCasimirScalar_zero_iff {γ : ℤ × ℤ} (hγ : borcherdsPositiveCone γ) :
    borcherdsCasimirScalar γ = 0 ↔ γ.1 = 1 := by
theorem borcherdsSimpleRoot_sub_not_positive {α β : ℤ × ℤ}
    (hα : borcherdsSimpleRoot α) (hβ : borcherdsPositiveRoot β) :
    ¬ borcherdsPositiveRoot (α - β) := by
\end{leancode}
(\leanpath{VOA/Moonshine/MoonshineCNBorcherdsRoots.lean}.)
\begin{dictionary}
Literal transcriptions.  \lean{γ.1}, \lean{γ.2} are the coordinates $m,n$.  $\alpha_r^2=-(1\cdot(-1)+1\cdot(-1))=2$,
$\alpha_j^2=-2j$.  The key arithmetic: $Q(m,n)=-2mn+2n=-2n(m-1)$, which on the positive cone vanishes exactly on the
line $m=1$.  \textbf{$\delta$ is a Weyl vector}: $2(\delta,\alpha)_{\mathrm{Bor}}=(\alpha,\alpha)_{\mathrm{Bor}}$ for every simple root ($2\cdot1=2$ for
$\alpha_r$, $2\cdot(-j)=-2j$ for $\alpha_j$), which is what makes the Casimir central below.  Also proved here:
\lean{borcherdsPositiveRoot_add_simple}, \lean{borcherdsPositiveCone_height_finite}, \lean{borcherdsRootSupport_neg}.
\end{dictionary}
\begin{physics}
Roots are the allowed momenta of the gauge bosons: by the no-ghost theorem, $\chXVIIIm_{(m,n)}\cong\Vnat_{1+mn}$ is nonzero
only if $mn=-1$ (the real roots $\pm\alpha_r$, tachyons) or $mn\ge1$ (massive states balanced by momentum).  The
Weyl vector $\delta=(1,0)$ is the lightlike vector against which all simple roots have ``height'' $\frac12$ their norm.
\end{physics}
\begin{status}No hypotheses; standard axioms (\leanpath{VOA/Moonshine/MoonshineCNAudit12.lean}).\end{status}

\subsection{The structure}

\begin{definition}[Borcherds datum]\label{ch18:def:bd}
A \emph{Borcherds datum} on a complex Lie algebra $\chXVIIIg$ consists of
\begin{enumerate}[label=(BD\arabic*)]
\item a grading $\chXVIIIg=\bigoplus_{\gamma\in\Z^2}\chXVIIIg_\gamma$ with finite-dimensional pieces and $[\chXVIIIg_\gamma,\chXVIIIg_\delta]\subseteq\chXVIIIg_{\gamma+\delta}$;
\item an injective linear map $\iota:\mathbb C^2\to\chXVIIIg$ with image $\chXVIIIg_0$ and $[\iota\xi,x]=(\xi,\gamma)_Bx$ for $x\in\chXVIIIg_\gamma$;
\item a symmetric invariant bilinear form with $(\chXVIIIg_\gamma,\chXVIIIg_\delta)=0$ unless $\gamma+\delta=0$, $(\iota\xi,\iota\mu)=(\xi,\mu)_{\mathrm{Bor}}$, and
  $\chXVIIIg_\gamma\times\chXVIIIg_{-\gamma}\to\mathbb C$ nondegenerate (left) for every $\gamma$;
\item root support: $\chXVIIIg_\gamma\neq0\Rightarrow\gamma$ is supported; $\dim\chXVIIIg_{\pm\alpha_r}=1$;
\item generation: $\chXVIIIg$ is generated as a Lie algebra by $\chXVIIIg_0$, $\chXVIIIg_{\pm\alpha_r}$ and the $\chXVIIIg_{\pm\alpha_j}$;
\end{enumerate}
and, optionally, (BD6) a group acting by graded, form-preserving Lie automorphisms, trivially on $\chXVIIIg_0$ and $\chXVIIIg_{\pm\alpha_r}$.
\end{definition}
\begin{leancode}
def borcherdsSimpleGenerators (wt : ℤ × ℤ → Submodule ℂ g) : Set g :=
  {x | x ∈ wt 0 ∨ ∃ α, borcherdsSimpleRoot α ∧ (x ∈ wt α ∨ x ∈ wt (-α))}
structure BorcherdsData (g : Type*) [LieRing g] [LieAlgebra ℂ g] where
  weightSpace : ℤ × ℤ → Submodule ℂ g
  weight_internal : DirectSum.IsInternal weightSpace
  weight_finite : ∀ γ, FiniteDimensional ℂ (weightSpace γ)
  bracket_weight : ∀ {γ δ} {x y : g},
    x ∈ weightSpace γ → y ∈ weightSpace δ → ⁅x, y⁆ ∈ weightSpace (γ + δ)
  cartanMap : (ℂ × ℂ) →ₗ[ℂ] g
  cartan_injective : Function.Injective cartanMap
  cartan_range : LinearMap.range cartanMap = weightSpace 0
  cartan_action : ∀ (ξ : ℂ × ℂ) (γ : ℤ × ℤ) (x : g), x ∈ weightSpace γ →
    ⁅cartanMap ξ, x⁆ = borcherdsCartanPairing ξ (borcherdsWeightVector γ) • x
  form : LinearMap.BilinForm ℂ g
  form_symmetric : ∀ x y, form x y = form y x
  form_invariant : ∀ x y z, form ⁅x, y⁆ z = form x ⁅y, z⁆
  form_orthogonal : ∀ γ δ, γ + δ ≠ 0 → ∀ x ∈ weightSpace γ,
    ∀ y ∈ weightSpace δ, form x y = 0
  form_cartan : ∀ ξ μ, form (cartanMap ξ) (cartanMap μ) = borcherdsCartanPairing ξ μ
  form_nondegenerate : ∀ γ, ∀ x ∈ weightSpace γ,
    (∀ y ∈ weightSpace (-γ), form x y = 0) → x = 0
  root_support : ∀ γ, weightSpace γ ≠ ⊥ → borcherdsRootSupport γ
  real_dimension : Module.finrank ℂ (weightSpace borcherdsRealRoot) = 1 ∧
    Module.finrank ℂ (weightSpace (-borcherdsRealRoot)) = 1
  generated : LieSubalgebra.lieSpan ℂ g (borcherdsSimpleGenerators weightSpace) = ⊤
structure GAction (G : Type*) [Group G] where
  action : G →* (g ≃ₗ[ℂ] g)
  preserves_lie : ∀ σ x y, action σ ⁅x, y⁆ = ⁅action σ x, action σ y⁆
  graded : ∀ σ γ x, x ∈ D.weightSpace γ → action σ x ∈ D.weightSpace γ
  preserves_form : ∀ σ x y, D.form (action σ x) (action σ y) = D.form x y
  fixes_cartan : ∀ σ x, x ∈ D.weightSpace 0 → action σ x = x
  fixes_real : ∀ σ x,
    x ∈ D.weightSpace borcherdsRealRoot ∨ x ∈ D.weightSpace (-borcherdsRealRoot) →
      action σ x = x
\end{leancode}
(\leanpath{VOA/Moonshine/MoonshineCNBorcherdsData.lean}; \lean{GAction} is \lean{BorcherdsData.GAction} with
\lean{D : BorcherdsData g} a section variable.)
\begin{dictionary}
\textbf{Field by field.}  (BD1) \lean{weight_internal}, \lean{weight_finite}, \lean{bracket_weight}; (BD2)
\lean{cartanMap}, \lean{cartan_injective}, \lean{cartan_range}, \lean{cartan_action}; (BD3) \lean{form},
\lean{form_symmetric}, \lean{form_invariant}, \lean{form_orthogonal}, \lean{form_cartan}, \lean{form_nondegenerate}; (BD4)
\lean{root_support}, \lean{real_dimension}; (BD5) \lean{generated}.  \lean{LieSubalgebra.lieSpan ℂ g S} is the Lie
subalgebra generated by $S$; \lean{= ⊤} says it is all of $\chXVIIIg$.

\textbf{The docstring is explicit}: ``Generation is an input here; no assertion is made that the Monster Lie
algebra has it.''  The Lean field \lean{form_nondegenerate} is left-nondegeneracy of the pairing
$\chXVIIIg_\gamma\times\chXVIIIg_{-\gamma}$; by symmetry and finite dimension it is a perfect pairing
(\lean{rootPairingEquiv}, \S\ref{ch18:sec:casimir}).

\textbf{Conventions.}  $\iota$ is a map from $\mathbb C^2$ with the Borcherds pairing, i.e.\ from $\chXVIIIII\otimes\mathbb C$.
\textbf{No positivity} (no Hermitian form, no positive-definiteness) is part of the structure: the arguments below use only
invariance, symmetry and nondegeneracy of the opposite-root pairings.

\textbf{Is it satisfiable?}  Yes: $\chXVIIIm$ is an instance (Theorem~\ref{ch18:thm:monsterbd}, given generation, which is
proved in Chapter~\ref{ch:20}).
\end{dictionary}
\begin{physics}
The abstract properties of a gauge algebra of a compactified string with two-dimensional ``light-cone'' momentum
lattice: momentum grading, Cartan currents, a Killing form from the current two-point function, and the
spectrum of allowed momenta.
\end{physics}
\begin{status}Definitions; standard axioms (\leanpath{VOA/Moonshine/MoonshineCNAudit12.lean}).\end{status}

\begin{theorem}[First consequences]\label{ch18:thm:bdbasic}
In every Borcherds datum: (a) for $x\in\chXVIIIg_\gamma$, $y\in\chXVIIIg_{-\gamma}$, $[x,y]=(x,y)\,\iota(\gamma)$; (b) there are $e\in\chXVIIIg_{\alpha_r}$,
$f\in\chXVIIIg_{-\alpha_r}$ with $(e,f)=1$; for any such pair, $(h,e,f)$ with $h=\iota(1,-1)$ is an $\mathfrak{sl}_2$-triple:
$[e,f]=h$, $[h,e]=2e$, $[h,f]=-2f$; (c) $[f,v]=0$ for $v\in\chXVIIIg_{\alpha_j}$, $j\ge1$.
\end{theorem}
\begin{leancode}
theorem bracket_eq_form_smul_cartan (γ : ℤ × ℤ) (x y : g)
    (hx : x ∈ D.weightSpace γ) (hy : y ∈ D.weightSpace (-γ)) :
    ⁅x, y⁆ = D.form x y • D.cartanMap (borcherdsWeightVector γ) := by
theorem exists_normalized_real_pair : ∃ e f : g,
    e ∈ D.weightSpace borcherdsRealRoot ∧ f ∈ D.weightSpace (-borcherdsRealRoot) ∧
      D.form e f = 1 := by
theorem real_sl2_relations (e f : g)
    (he : e ∈ D.weightSpace borcherdsRealRoot)
    (hf : f ∈ D.weightSpace (-borcherdsRealRoot)) (hef : D.form e f = 1) :
    ⁅e, f⁆ = D.cartanMap (1, -1) ∧
      ⁅D.cartanMap (1, -1), e⁆ = (2 : ℂ) • e ∧
      ⁅D.cartanMap (1, -1), f⁆ = (-2 : ℂ) • f := by
theorem simple_lowering_vanishes (j : ℤ) (hj : 1 ≤ j) (f v : g)
    (hf : f ∈ D.weightSpace (-borcherdsRealRoot))
    (hv : v ∈ D.weightSpace (borcherdsImaginaryRoot j)) : ⁅f, v⁆ = 0 := by
\end{leancode}
(\leanpath{VOA/Moonshine/MoonshineCNBorcherdsData.lean}.)
\begin{dictionary}
The normalized pair $(e,f)$ appears as explicit arguments \lean{(e f : g) (he) (hf) (hef : D.form e f = 1)} in
all later theorems; \lean{exists_normalized_real_pair} shows such arguments always exist (docstring: ``no choice
of such a pair is an extra input''), so these hypotheses are never restrictive.
\end{dictionary}
\begin{physics}The real root vectors and their commutator form an $\mathfrak{su}(2)$-type (``W-boson'') gauge subalgebra.\end{physics}
\paragraph{Proof on paper.}  (a) $[x,y]\in\chXVIIIg_0=\iota(\mathbb C^2)$ and $([x,y],\iota\lambda)=(x,[y,\iota\lambda])=(\lambda,\gamma)_{\mathrm{Bor}}(x,y)=((x,y)\iota\gamma,\iota\lambda)$;
the form is nondegenerate on $\chXVIIIg_0$.  (b) $\chXVIIIg_{\pm\alpha_r}$ are lines paired nondegenerately; (a) gives $[e,f]=\iota(\alpha_r)$, and
$(\alpha_r,\alpha_r)_{\mathrm{Bor}}=2$.  (c) $[f,v]\in\chXVIIIg_{(0,j+1)}$, which is unsupported, hence $0$. \qed
\begin{status}No hypotheses beyond the datum; standard axioms (\leanpath{VOA/Moonshine/MoonshineCNAudit12.lean}).\end{status}

\section{The Monster Lie algebra is a Borcherds datum}\label{ch18:sec:monsterbd}

\paragraph{Where the Cartan map comes from.}  By Theorems~\ref{ch18:thm:zerospace} and \ref{ch18:thm:cartanaction}, an element $h\in\chXVIIIm_0$
with no-ghost coordinates $(c_0,c_1)$ is $c_0$ times the current of $w_0=(0,1)$ plus $c_1$ times the current of $w_1=(1,0)$, i.e.\ it
corresponds to the lattice vector $\lambda(h)=(c_1,c_0)$, and it acts on repository momentum $\beta$ by $(\lambda(h),\beta)=c_0\beta_0+c_1\beta_1$.  A
Borcherds root $(m,n)$ has repository momentum $(m,-n)$, so $h$ acts on it by $c_0m-c_1n$.  We want $\iota(\xi)$ to act by
$(\xi,(m,n))_{\mathrm{Bor}}=-\xi_1n-\xi_2m$; matching coefficients gives $c_0=-\xi_2$, $c_1=\xi_1$.  This explains the coordinates
$(-\xi_2,\xi_1)$ in the next definition.

\begin{definition}[Borcherds coordinates on $\chXVIIIm$]\label{ch18:def:monstercoords}
$\chXVIIIm_{(m,n)}^{\mathrm{Bor}}:=\chXVIIIm_{(m,-n)}$ (repository momentum), and $\iota(\xi):=$ the element of $\chXVIIIm_0$ with no-ghost coordinates
$(-\xi_2,\xi_1)$.
\end{definition}
\begin{leancode}
def borcherdsMomentum (γ : ℤ × ℤ) : lorentzianLattice.carrier := ![γ.1, -γ.2]
def monsterLieBorcherdsWeight (γ : ℤ × ℤ) : Submodule ℂ monsterLieAlgebra :=
  monsterLieMomentum (borcherdsMomentum γ)
def borcherdsCartanCoordinates : (ℂ × ℂ) →ₗ[ℂ] (Fin 2 → ℂ) where
  toFun ξ := ![-ξ.2, ξ.1]
def monsterLieCartanMap : (ℂ × ℂ) →ₗ[ℂ] monsterLieAlgebra :=
  (monsterLieMomentum 0).subtype ∘ₗ monsterLieZeroMomentumEquiv.symm.toLinearMap ∘ₗ
    borcherdsCartanCoordinates
\end{leancode}
(\leanpath{VOA/Moonshine/MoonshineCNMonsterBorcherdsData.lean}.)
\begin{dictionary}
Check of the Cartan action: $\iota(\xi)$ has coordinates $(c_0,c_1)=(-\xi_2,\xi_1)$; on $\chXVIIIm^{\mathrm{Bor}}_{(m,n)}=\chXVIIIm_{(m,-n)}$
Theorem~\ref{ch18:thm:cartanaction} gives the scalar $c_0m+c_1(-n)=-\xi_2m-\xi_1n=(\xi,(m,n))_{\mathrm{Bor}}$.  Check of the form: by
Theorem~\ref{ch18:thm:zeroform}, $(\iota\xi,\iota\mu)=(-\xi_2)\mu_1+\xi_1(-\mu_2)=(\xi,\mu)_{\mathrm{Bor}}$.  These are the Lean theorems
\lean{monsterLieCartanMap_action} and \lean{monsterLieCartanMap_form}.
\end{dictionary}
\begin{physics}
Borcherds coordinates $(m,n)$ are light-cone momenta with the sign of the second one flipped, so that the
positive cone, where the massive gauge bosons live, is $m,n\ge1$; $\iota(\xi)$ is the momentum current in the direction $\xi$.
\end{physics}

\begin{theorem}[Borcherds-coordinate root spaces]\label{ch18:thm:borcherdsroot}
For $\gamma=(m,n)\neq0$, $\chXVIIIm^{\mathrm{Bor}}_\gamma\cong\Vnat_{1+mn}$ as vector spaces (a linear isomorphism; its $\Aut(\Vnat)$-equivariance is
the repository-coordinate statement \lean{monsterLieRootSpaceEquiv_equivariant}, since the map is \lean{monsterLieRootSpaceEquiv m (-n)}
followed by a relabelling of the weight); $\chXVIIIm^{\mathrm{Bor}}_\gamma\neq0$ implies that $\gamma$ is supported (using $\Vnat_{<0}=0$,
$\Vnat_1=0$); $\dim\chXVIIIm^{\mathrm{Bor}}_{\pm\alpha_r}=1$.
\end{theorem}
\begin{leancode}
def monsterLieBorcherdsRootEquiv (γ : ℤ × ℤ) (hγ : γ ≠ 0) :
    monsterLieBorcherdsWeight γ ≃ₗ[ℂ] VnatVOA.grading (1 + γ.1 * γ.2) :=
theorem monsterLieBorcherdsWeight_support (γ : ℤ × ℤ) (hne : monsterLieBorcherdsWeight γ ≠ ⊥) :
    borcherdsRootSupport γ := by
theorem monsterLieBorcherdsWeight_real_finrank (γ : ℤ × ℤ)
    (hγ : γ = borcherdsRealRoot ∨ γ = -borcherdsRealRoot) :
    Module.finrank ℂ (monsterLieBorcherdsWeight γ) = 1 := by
\end{leancode}
(\leanpath{VOA/Moonshine/MoonshineCNMonsterBorcherdsData.lean}.)
\begin{dictionary}
\lean{monsterLieBorcherdsRootEquiv} is \lean{monsterLieRootSpaceEquiv m (-n)} composed with the identity
$1-m(-n)=1+mn$: this is the bridging lemma of the sign warning in \S\ref{ch18:sec:conventions}.  Support: if
$\Vnat_{1+mn}\neq0$ then $1+mn\ge0$ and $1+mn\ne1$, i.e.\ $mn\ge-1$, $mn\ne0$, and integer arithmetic
(\lean{borcherdsRootSupport_of_product}) gives $\gamma\in\{\pm\alpha_r\}\cup C\cup(-C)$.  Real roots: $1+mn=0$, and $\dim\Vnat_0=1$
(\lean{vnat_finrank_zero}).
\end{dictionary}
\begin{physics}
The tachyon ($\Vnat_0$) at the real roots; nothing on the axes $m=0$ or $n=0$ except the Cartan currents (since
$\Vnat_1=0$: no internal currents); everything else lives in the cones $\pm C$.
\end{physics}
\begin{status}No hypotheses; standard axioms (\leanpath{VOA/Moonshine/MoonshineCNAudit17.lean}).\end{status}

\begin{definition}[Simple-root generation]\label{ch18:def:hgen}
\lean{MonsterLieSimpleRootGeneration}: $\chXVIIIm$ is generated as a Lie algebra by $\chXVIIIm_0$, $\chXVIIIm^{\mathrm{Bor}}_{\pm\alpha_r}$ and the
$\chXVIIIm^{\mathrm{Bor}}_{\pm\alpha_j}$, $j\ge1$.
\end{definition}
\begin{leancode}
def MonsterLieSimpleRootGeneration : Prop :=
  LieSubalgebra.lieSpan ℂ monsterLieAlgebra (borcherdsSimpleGenerators monsterLieBorcherdsWeight) =
    ⊤
\end{leancode}
(\leanpath{VOA/Moonshine/MoonshineCNMonsterBorcherdsData.lean}.)
\begin{physics}
All gauge bosons of the Monster string arise as commutators of the momentum currents, the tachyon pair and the simple-root
states $\Vnat_{1+j}\otimes e^{(1,-j)}$: there are no further ``elementary'' gauge bosons.  Physically this is a statement about
the full spectrum, which is why it is equivalent to the modular product formula for $j(p)-j(q)$.
\end{physics}

\begin{theorem}[The Monster Lie algebra as Borcherds datum]\label{ch18:thm:monsterbd}
Given \lean{MonsterLieSimpleRootGeneration}, $\chXVIIIm$ with the Borcherds-coordinate grading, $\iota$ and the form of
\S\ref{ch18:sec:form} is a Borcherds datum, and $\Aut(\Vnat)$ acts on it as in (BD6).
\end{theorem}
\begin{leancode}
def monsterLieBorcherdsData (hgen : MonsterLieSimpleRootGeneration) :
    BorcherdsData monsterLieAlgebra where
  weightSpace := monsterLieBorcherdsWeight
  weight_internal := monsterLieBorcherdsWeight_isInternal
  weight_finite := monsterLieBorcherdsWeight_finite
  ...
  cartanMap := monsterLieCartanMap
  cartan_injective := monsterLieCartanMap_injective
  cartan_range := monsterLieCartanMap_range
  cartan_action := monsterLieCartanMap_action
  form := monsterLieBilinearForm
  form_symmetric := monsterLieBilinearForm_symmetric
  form_invariant := monsterLieBilinearForm_invariant
  ...
  form_cartan := monsterLieCartanMap_form
  ...
  root_support := monsterLieBorcherdsWeight_support
  ...
  generated := hgen
def monsterLieBorcherdsAction (hgen : MonsterLieSimpleRootGeneration) :
    (monsterLieBorcherdsData hgen).GAction AutVnat where
  action := monsterLieLinearAction
  preserves_lie σ x y := (monsterLieAut σ).map_lie x y
  graded σ γ x hx := monsterLieAut_momentum σ _ x hx
  preserves_form σ x y := monsterLieBilinearForm_aut σ x y
  ...
\end{leancode}
(\leanpath{VOA/Moonshine/MoonshineCNMonsterBorcherdsData.lean}; the lines marked \texttt{...} elide the fields \lean{bracket_weight}, \lean{form_orthogonal},
\lean{form_nondegenerate}, \lean{real_dimension}, \lean{fixes_cartan}, \lean{fixes_real}, whose proofs are a few
lines each and reduce to \lean{monsterLieMomentum_bracket}, \lean{monsterLieBilinearForm_momentum_orthogonal},
\lean{monsterLieBilinearForm_pairing_nondegenerate}, \lean{monsterLieBorcherdsWeight_real_finrank},
\lean{monsterLieZeroMomentum_aut_trivial} and \lean{monsterLieAut_fixes_real}.)
\begin{dictionary}
\textbf{The one hypothesis.}  \lean{hgen} is used for exactly one field, \lean{generated := hgen}.  Every other field
is a theorem proved in this chapter.  The file docstring: ``Every field of the data is proved here from the
vertex-algebra construction, except generation \dots{} kept as the explicit hypothesis
\lean{MonsterLieSimpleRootGeneration}; it is equivalent, given the other fields, to the product formula for
$j(p)-j(q)$, which is not yet in Lean.''  \textbf{That last clause is out of date} (see the next paragraph).

\textbf{(BD6) for $\Aut(\Vnat)$.}  \lean{fixes_real}: $\chXVIIIm_{\pm\alpha_r}\cong\Vnat_0=\mathbb C\mathbb 1$ and automorphisms fix $\mathbb 1$
(\lean{monsterLieAut_fixes_real}); \lean{fixes_cartan}: Theorem~\ref{ch18:thm:zeroform}.
\end{dictionary}

\paragraph{The hypothesis is discharged.}  In the release the product formula is in Lean
(\lean{M8.knz_product}, Chapter~\ref{ch:19}) and generation is a theorem with no hypotheses:
\begin{leancode}
theorem monsterLieSimpleRootGeneration : MonsterLieSimpleRootGeneration :=
  monsterLiePreData.lieSpan_eq_top_of_KNZ monsterLiePreData_KNZ
\end{leancode}
(\leanpath{VOA/Moonshine/MoonshineLRGenerationMonster.lean}, \lean{VOA.Moonshine.ConwayNorton.M8.monsterLieSimpleRootGeneration},
documented in Chapter~\ref{ch:20}; printed in \leanpath{VOA/Moonshine/MoonshineCNFinalAudit.lean} with standard axioms.)
Hence \lean{monsterLieBorcherdsData M8.monsterLieSimpleRootGeneration} is an unconditional Borcherds datum, and every
consequence stated below for $\chXVIIIm$ ``given \lean{hgen}'' holds unconditionally after this substitution.
\begin{physics}
The gauge algebra of the Monster string has exactly the root data of a rank-two Borcherds algebra with one real
simple root and imaginary simple roots $(1,j)$ of multiplicity $\dim\Vnat_{j+1}$, and the Monster acts preserving all of it.
\end{physics}
\begin{status}
Conditional on \lean{hgen : MonsterLieSimpleRootGeneration} in this file; that hypothesis is a theorem with no
hypotheses (\lean{M8.monsterLieSimpleRootGeneration}).  \texttt{\#print axioms}: standard
(\leanpath{VOA/Moonshine/MoonshineCNAudit17.lean} for the instance, \leanpath{VOA/Moonshine/MoonshineCNFinalAudit.lean} for
the discharge).
\end{status}

\begin{warning}
\textbf{Generation cannot be proved Lie-theoretically here; it is equivalent to the product formula.}  Given everything else,
``$\chXVIIIm$ is generated by its simple root spaces'' $\Longleftrightarrow$ the untwisted denominator identity
$\prod_{m,n\ge1}(1-p^mq^n)^{c(mn)}=1-\sum_{a,b\ge1}c(a+b-1)p^aq^b$.  Deriving the product identity \emph{from} $\chXVIIIm$ is therefore
circular.  The Casimir argument of \S\ref{ch18:sec:resolution} excludes \emph{relations} among the known simple root
vectors, not \emph{extra} simple roots $(m,n)$ with $m\ge2$ (for those, $\delta$ is not known to be a Weyl vector).  The
library proves the product identity analytically from the modularity of $J$ (Chapter~\ref{ch:19}) and then deduces
generation by a dimension sandwich (Chapter~\ref{ch:20}).
\end{warning}

\section{Real-root strings and the generating space}\label{ch18:sec:strings}

Fix a Borcherds datum $D$ on $\chXVIIIg$ and a normalized pair $(e,f)$.

\subsection{\texorpdfstring{The $\mathfrak{sl}_2$-strings of the imaginary simple roots}{The sl2-strings of the imaginary simple roots}}

\begin{theorem}[Real-root strings]\label{ch18:thm:sl2}
Let $j\ge1$ and $v\in\chXVIIIg_{\alpha_j}$.  (a) $(h,e,f)$ is an $\mathfrak{sl}_2$-triple and $[h,v]=(1-j)v$; (b) $(\mathrm{ad}\,e)^kv\in\chXVIIIg_{(1+k,j-k)}$ and
$(\mathrm{ad}\,e)^jv=0$; (c) $[f,(\mathrm{ad}\,e)^{k+1}v]=(k+1)(j-k-1)(\mathrm{ad}\,e)^kv$; (d) for $0\le k<j$, $(\mathrm{ad}\,e)^k$ is injective on
$\chXVIIIg_{\alpha_j}$.
\end{theorem}
\begin{leancode}
theorem real_sl2_triple (e f : g)
    (he : e ∈ D.weightSpace borcherdsRealRoot)
    (hf : f ∈ D.weightSpace (-borcherdsRealRoot)) (hef : D.form e f = 1) :
    IsSl2Triple (D.cartanMap (1, -1)) e f := by
theorem adReal_pow_simple_vanishes (e : g) (he : e ∈ D.weightSpace borcherdsRealRoot)
    (j : ℕ) (v : g) (hv : v ∈ D.weightSpace (borcherdsImaginaryRoot (j : ℤ))) :
    (LieModule.toEnd ℂ g g e ^ j) v = 0 := by
theorem real_lowering_string (e f : g)
    (he : e ∈ D.weightSpace borcherdsRealRoot)
    (hf : f ∈ D.weightSpace (-borcherdsRealRoot)) (hef : D.form e f = 1)
    (j : ℤ) (hj : 1 ≤ j) (v : g)
    (hv : v ∈ D.weightSpace (borcherdsImaginaryRoot j)) (k : ℕ) :
    ⁅f, (LieModule.toEnd ℂ g g e ^ (k + 1)) v⁆ =
      (((k : ℂ) + 1) * ((j : ℂ) - (k : ℂ) - 1)) •
        (LieModule.toEnd ℂ g g e ^ k) v := by
theorem adReal_pow_injective (e f : g)
    (he : e ∈ D.weightSpace borcherdsRealRoot)
    (hf : f ∈ D.weightSpace (-borcherdsRealRoot)) (hef : D.form e f = 1)
    (j : ℤ) (hj : 1 ≤ j) (k : ℕ) (hk : (k : ℤ) < j) :
    Function.Injective (fun v : D.weightSpace (borcherdsImaginaryRoot j) =>
      (LieModule.toEnd ℂ g g e ^ k) (v : g)) := by
\end{leancode}
(\leanpath{VOA/Moonshine/MoonshineCNBorcherdsSl2.lean}.)
\begin{dictionary}
\lean{LieModule.toEnd ℂ g g e} is $\mathrm{ad}\,e$; \lean{^ k} its $k$-th power.  \lean{IsSl2Triple h e f} is Mathlib's
$\mathfrak{sl}_2$-triple predicate.  (c) is the lowest-weight formula with $v$ a lowest-weight vector of weight $-(j-1)$.
\end{dictionary}
\begin{physics}
Each imaginary simple root state $v$ is the bottom of a $j$-dimensional multiplet of the W-boson $\mathfrak{sl}_2$ (spin
$\frac{j-1}2$): $v,[e,v],\dots,(\mathrm{ad}\,e)^{j-1}v$.
\end{physics}
\paragraph{Proof on paper.}  (a) by Theorem~\ref{ch18:thm:bdbasic}; $(\alpha_r,\alpha_j)_{\mathrm{Bor}}=-(1\cdot j+1\cdot(-1))=1-j$.  (b) $(1+j,0)$ is
unsupported.  (c) induction with $[f,e]=-h$ and $[f,v]=0$.  (d) apply $f$ repeatedly: the coefficients $(k+1)(j-k-1)$
are nonzero for $k+1<j$. \qed
\begin{status}No hypotheses beyond the datum and the pair; standard axioms (\leanpath{VOA/Moonshine/MoonshineCNAudit12.lean}).\end{status}

\begin{definition}[The generating space $U$]\label{ch18:def:U}
$U_{j,k}:=(\mathrm{ad}\,e)^k\chXVIIIg_{\alpha_j}$ and $U:=\sum_{j\ge1}\sum_{0\le k<j}U_{j,k}$.
\end{definition}
\begin{leancode}
def generatorString (e : g) (j : ℤ) (k : ℕ) : Submodule ℂ g :=
  (D.weightSpace (borcherdsImaginaryRoot j)).map (LieModule.toEnd ℂ g g e ^ k)
def genSpace (e : g) : Submodule ℂ g :=
  ⨆ (j : ℤ) (_ : 1 ≤ j) (k : ℕ) (_ : (k : ℤ) < j), D.generatorString e j k
\end{leancode}
(\leanpath{VOA/Moonshine/MoonshineCNBorcherdsSl2.lean}.)
\begin{physics}
$U$ is the span of all states obtained from simple-root states by acting with the W-boson raising operator, before the
multiplet terminates.
\end{physics}

\begin{theorem}[Structure of $U$]\label{ch18:thm:Ustructure}
(a) $U_{j,k}\subseteq\chXVIIIg_{(1+k,j-k)}$, $(\mathrm{ad}\,e)^k:\chXVIIIg_{\alpha_j}\to U_{j,k}$ is an isomorphism, $U\cap\chXVIIIg_{(1+k,j-k)}=U_{j,k}$, and
$U=\bigoplus_{j,k}U_{j,k}$; so $U_{(a,b)}\cong\chXVIIIg_{\alpha_{a+b-1}}$ for $a,b\ge1$.  (b) $U\subseteq\chXVIIIu^+:=\bigoplus_{\gamma\in C}\chXVIIIg_\gamma$.  (c) $U$
is stable under $\iota(\mathbb C^2)$, $\mathrm{ad}\,e$, $\mathrm{ad}\,f$, hence is a module for the Levi subalgebra
$\chXVIIIl=\chXVIIIg_{-\alpha_r}\oplus\chXVIIIg_0\oplus\chXVIIIg_{\alpha_r}$.  (d) $[\chXVIIIu^-,U]\subseteq\chXVIIIp^-$, where $\chXVIIIu^-=\bigoplus_{-\gamma\in C}\chXVIIIg_\gamma$ and
$\chXVIIIp^-=\bigoplus_{\gamma\notin C}\chXVIIIg_\gamma$.  (e) $U$ is stable under every symmetry (BD6).
\end{theorem}
\begin{leancode}
def generatorStringEquiv (e f : g)
    (he : e ∈ D.weightSpace borcherdsRealRoot)
    (hf : f ∈ D.weightSpace (-borcherdsRealRoot)) (hef : D.form e f = 1)
    (j : ℤ) (hj : 1 ≤ j) (k : ℕ) (hk : (k : ℤ) < j) :
    D.weightSpace (borcherdsImaginaryRoot j) ≃ₗ[ℂ] D.generatorString e j k :=
theorem genSpace_weight_intersection (e : g) (he : e ∈ D.weightSpace borcherdsRealRoot)
    (i : borcherdsStringIndex) :
    D.genSpace e ⊓ D.weightSpace (borcherdsStringDegree i) =
      D.generatorString e i.val.1 i.val.2 := by
def genSpaceDirectSumEquiv (e : g) (he : e ∈ D.weightSpace borcherdsRealRoot) :
    (⨁ i : borcherdsStringIndex, D.generatorString e i.val.1 i.val.2) ≃ₗ[ℂ] D.genSpace e :=
theorem genSpace_levi_stable (e f : g)
    (he : e ∈ D.weightSpace borcherdsRealRoot)
    (hf : f ∈ D.weightSpace (-borcherdsRealRoot)) (hef : D.form e f = 1)
    {x u : g} (hx : x ∈ D.leviPart) (hu : u ∈ D.genSpace e) :
    ⁅x, u⁆ ∈ D.genSpace e := by
theorem negative_bracket_genSpace (e : g) (he : e ∈ D.weightSpace borcherdsRealRoot)
    {y u : g} (hy : y ∈ D.negativePart) (hu : u ∈ D.genSpace e) :
    ⁅y, u⁆ ∈ D.parabolicPart :=
\end{leancode}
(\leanpath{VOA/Moonshine/MoonshineCNBorcherdsSl2.lean}, \leanpath{VOA/Moonshine/MoonshineCNBorcherdsGeneratorWeights.lean},
\leanpath{VOA/Moonshine/MoonshineCNBorcherdsStringSum.lean}, \leanpath{VOA/Moonshine/MoonshineCNBorcherdsStringAction.lean};
equivariance \lean{genSpace_equivariant}.)
\begin{dictionary}
\lean{borcherdsStringIndex} is $\{(j,k):1\le j,\ 0\le k<j\}$ and \lean{borcherdsStringDegree (j,k) = (1+k, j-k)}, an
injective map (distinct strings positions have distinct degrees, which is why the sum is direct).
\lean{⨁} is the external direct sum.  The subalgebras \lean{positivePart} ($\chXVIIIu^+$), \lean{negativePart} ($\chXVIIIu^-$),
\lean{leviPart} ($\chXVIIIl$), \lean{parabolicPart} ($\chXVIIIp^-$) are defined next.
\end{dictionary}
\begin{physics}$U$ collects all the W-boson multiplets built on the imaginary simple root states.\end{physics}
\begin{status}No hypotheses beyond the datum and the pair; standard axioms (\leanpath{VOA/Moonshine/MoonshineCNAudit12.lean},
\leanpath{VOA/Moonshine/MoonshineCNAudit14.lean}).\end{status}

\subsection{Parabolic decomposition}

The names come from the finite-dimensional theory: a \emph{parabolic} subalgebra contains a Cartan subalgebra and all
negative root spaces (here: everything outside the positive cone), its \emph{Levi} subalgebra is the reductive part
(here $\chXVIIIl\cong\mathfrak{gl}_2$, spanned by the Cartan and the real root lines), and the \emph{nilradical} is the remaining ideal
(here $\chXVIIIu^-$).

\begin{definition}[Positive part, Levi, negative part, opposite parabolic]\label{ch18:def:parabolic}
For a predicate $P$ on degrees let $\chXVIIIg_P:=\sum_{P(\gamma)}\chXVIIIg_\gamma$.  Then $\chXVIIIu^+=\chXVIIIg_{C}$, $\chXVIIIp^-=\chXVIIIg_{\neg C}$, $\chXVIIIu^-=\chXVIIIg_{-C}$,
$\chXVIIIl=\chXVIIIg_{\{0,\pm\alpha_r\}}$; all are Lie subalgebras.
\end{definition}
\begin{leancode}
def weightSpan (P : ℤ × ℤ → Prop) : Submodule ℂ g :=
  ⨆ γ, ⨆ (_ : P γ), D.weightSpace γ
def positivePart : LieSubalgebra ℂ g where
  __ := D.weightSpan borcherdsPositiveCone
def parabolicPart : LieSubalgebra ℂ g where
  __ := D.weightSpan (fun γ => ¬ borcherdsPositiveCone γ)
theorem positive_parabolic_isCompl :
    IsCompl D.positivePart.toSubmodule D.parabolicPart.toSubmodule := by
\end{leancode}
(\leanpath{VOA/Moonshine/MoonshineCNBorcherdsParabolic.lean}.)
\begin{leancode}
def negativePart : LieSubalgebra ℂ g where
  __ := D.weightSpan borcherdsNegativeCone
def leviPart : LieSubalgebra ℂ g where
  __ := D.weightSpan borcherdsLeviDegree
def negativeParabolicIdeal : LieIdeal ℂ D.parabolicPart where
theorem parabolic_eq_levi_sup_negative : D.parabolicPart.toSubmodule =
    D.leviPart.toSubmodule ⊔ D.negativePart.toSubmodule := by
\end{leancode}
(\leanpath{VOA/Moonshine/MoonshineCNBorcherdsLevi.lean}.)
\begin{dictionary}
\textbf{Important: \lean{positivePart} is $\chXVIIIu^+$, the positive cone $C=\{m,n\ge1\}$; it does not contain the real root.}
The ``positive part'' $\mathfrak n^+=\chXVIIIg_{\alpha_r}\oplus\chXVIIIu^+$ of Kac--Moody theory (Borcherds' $E$) is not given a name.  $\chXVIIIp^-$
contains $\chXVIIIg_0$, both real root lines and $\chXVIIIu^-$ (the support forbids anything else outside $C$), and
$\chXVIIIu^-$ is an ideal of $\chXVIIIp^-$ (\lean{negativeParabolicIdeal}).  The complement $\chXVIIIg=\chXVIIIu^+\oplus\chXVIIIp^-$ is the input of relative PBW.
The Levi projection along $\chXVIIIu^-$ is a Lie homomorphism $\chXVIIIp^-\to\chXVIIIl$ (\lean{parabolicLeviProjection},
\leanpath{VOA/Moonshine/MoonshineCNBorcherdsParabolicAction.lean}, using the general
\lean{leviProjectionOfComplement} of \leanpath{VOA/Moonshine/MoonshineCNLieProjection.lean}), and $\chXVIIIp^-$ acts on $U$ through
it, $\chXVIIIu^-$ acting by $0$ (\lean{genSpaceParabolicRep}).
\end{dictionary}
\begin{physics}A light-cone-like splitting of the gauge algebra into ``creation'' ($C$) and ``annihilation plus zero modes'' parts.\end{physics}
\begin{status}No hypotheses beyond the datum; standard axioms (\leanpath{VOA/Moonshine/MoonshineCNAudit12.lean}).\end{status}

\subsection{Generation of the positive cone}

\begin{theorem}[$\chXVIIIu^+$ is generated by $U$]\label{ch18:thm:generation}
$\chXVIIIu^+$ equals the Lie subalgebra generated by $U$; consequently the augmentation ideal of $\mathcal U(\chXVIIIu^+)$ is the left ideal
$\mathcal U(\chXVIIIu^+)\cdot U$.
\end{theorem}
\begin{leancode}
theorem positivePart_eq_lieSpan_genSpace (e f : g)
    (he : e ∈ D.weightSpace borcherdsRealRoot)
    (hf : f ∈ D.weightSpace (-borcherdsRealRoot)) (hef : D.form e f = 1) :
    D.positivePart = LieSubalgebra.lieSpan ℂ g (D.genSpace e : Set g) := by
\end{leancode}
(\leanpath{VOA/Moonshine/MoonshineCNBorcherdsGeneration.lean}.)
\begin{leancode}
theorem positive_augmentationIdeal_eq (e f : g)
    (he : e ∈ D.weightSpace borcherdsRealRoot)
    (hf : f ∈ D.weightSpace (-borcherdsRealRoot)) (hef : D.form e f = 1) :
    LinearMap.ker (envelopingAugmentation (g := D.positivePart)).toLinearMap =
      (envelopingGeneratorLeftIdeal (D.positiveGeneratorSpace e)).restrictScalars ℂ :=
\end{leancode}
(\leanpath{VOA/Moonshine/MoonshineCNBorcherdsAugmentation.lean}, where \lean{envelopingAugmentation} is the algebra map
$\mathcal U(L)\to\mathbb C$ killing $L$, \lean{envelopingGeneratorLeftIdeal S} the left ideal $\mathcal U(L)S$, and
\lean{positiveGeneratorSpace e} is $U$ viewed inside $\chXVIIIu^+$.)
\begin{dictionary}
Uses (BD5) through the datum.  The augmentation statement holds for any Lie algebra and any Lie-generating subspace
(\lean{enveloping_augmentationIdeal_eq}: ``No PBW assumption is needed'').
\end{dictionary}
\begin{physics}All gauge bosons with momenta in the positive cone arise as nested commutators of the W-multiplets of the imaginary simple roots.\end{physics}
\paragraph{Proof on paper.}  The sum $S=\mathrm{Lie}(U)+\chXVIIIp^-$ is stable under $\mathrm{ad}$ of every generator of $\chXVIIIg$
(\lean{lieSpan_parabolic_bracket_mem}): $\chXVIIIp^-$ normalizes $\mathrm{Lie}(U)$ modulo $\chXVIIIp^-$ because $\chXVIIIl$ preserves $U$ and
$[\chXVIIIu^-,U]\subseteq\chXVIIIp^-$ (Theorem~\ref{ch18:thm:Ustructure}(c),(d)), and $e,\chXVIIIg_{\alpha_j}$ lie in $\chXVIIIp^-+U$.  So $S$ is a subalgebra containing
all simple generators, hence $S=\chXVIIIg$ by (BD5); comparing the $C$-graded parts gives $\chXVIIIu^+=\mathrm{Lie}(U)$.  For the augmentation
ideal: every element of $\chXVIIIu^+$ is a Lie polynomial without constant term in $U$, hence an associative one, hence in $\mathcal U(\chXVIIIu^+)U$. \qed
\begin{status}No hypotheses beyond the datum and the pair; standard axioms (\leanpath{VOA/Moonshine/MoonshineCNAudit14.lean}).\end{status}

\section{Generalized Verma modules and relative PBW}\label{ch18:sec:verma}

\subsection{Induced modules}

For a Lie subalgebra $\chXVIIIp\subseteq\chXVIIIg$ and a $\chXVIIIp$-module $W$ (a Lie homomorphism $\sigma:\chXVIIIp\to\End W$; Lean calls it \lean{ρ}), the induced module is
$\mathrm{Ind}(W)=\mathcal U(\chXVIIIg)\otimes_{U(\chXVIIIp)}W$, realised as the quotient of $\mathcal U(\chXVIIIg)\otimes W$ by the span of the balancing relations
$yx\otimes w-y\otimes\sigma(x)w$:
\begin{leancode}
abbrev relativeInducedSpace :=
  (UniversalEnvelopingAlgebra ℂ g ⊗[ℂ] U) ⧸ PBWTheorem.balancing p ρ
def relativeInducedRep : g →ₗ⁅ℂ⁆ End ℂ (relativeInducedSpace p ρ) :=
\end{leancode}
(\leanpath{VOA/Moonshine/MoonshineCNRelativeInduction.lean}; \lean{PBWTheorem.balancing} and \lean{PBWTheorem.bottom}, the
map $w\mapsto[1\otimes w]$, are in \leanpath{VOA/Foundations/PBWTheorem/Induced.lean}, Chapter~\ref{ch:L02}.  Here the type
variable named \lean{U} is the module $W$, not the generating space.)

\begin{theorem}[Relative PBW freeness]\label{ch18:thm:relpbw}
Let $\chXVIIIg=\chXVIIIp\oplus\mathfrak c$ as vector spaces, $(c_a)_{a\in A}$ a basis of $\mathfrak c$ indexed by a linearly ordered set.  Then
$\mathrm{Ind}(W)\cong\bigoplus_{\text{multisets }m\text{ in }A}W$, the class $[c_{a_1}\cdots c_{a_k}\otimes w]$ ($a_1\le\dots\le a_k$) corresponding to $w$ in
slot $m=\{a_1,\dots,a_k\}$.  If $\mathfrak c=\mathfrak q$ is a Lie subalgebra, then multiplication
$\mathcal U(\mathfrak q)\otimes W\to \mathrm{Ind}(W)$, $a\otimes w\mapsto[a\otimes w]$, is bijective.
\end{theorem}
\begin{leancode}
def relativePBWEquiv : relativeInducedSpace p ρ ≃ₗ[ℂ] (Multiset A →₀ U) :=
\end{leancode}
(\leanpath{VOA/Moonshine/MoonshineCNRelativePBW.lean}, with section variables \lean{(hc : IsCompl p.toSubmodule c)
(bc : Basis A ℂ c)}, \lean{[LinearOrder A]}.)
\begin{leancode}
theorem relativeVermaMultiplication_bijective :
    Function.Bijective (relativeVermaMultiplication p q ρ) := by
theorem parabolicVermaMultiplication_bijective :
    Function.Bijective (D.parabolicVermaMultiplication ρ) := by
\end{leancode}
(\leanpath{VOA/Moonshine/MoonshineCNVermaMultiplication.lean}, with \lean{(hc : IsCompl p.toSubmodule q.toSubmodule)}; the
second is the case $\chXVIIIp=\chXVIIIp^-$, $\mathfrak q=\chXVIIIu^+$, using \lean{positive_parabolic_isCompl}.)
\begin{dictionary}
\lean{Multiset A →₀ U} is the space of finitely supported functions from multisets of basis indices to $W$, i.e.\
$\bigoplus_mW$.  The proof uses the repository's relative PBW representation \lean{PBWTheorem.rep} (Chapter~\ref{ch:L02},
the construction that replaced the former axiom \lean{VOA.pbw}), so the result is axiom-clean.  No finite-dimensionality is
assumed; $A$ is any linear order (a well-order exists by choice).
\end{dictionary}
\begin{physics}
Induced representations are ``Fock spaces'': states are obtained from a vacuum multiplet $W$ by applying creation operators
from the complement, freely and in a fixed order.
\end{physics}
\begin{status}No hypotheses beyond complementarity; standard axioms (\leanpath{VOA/Moonshine/MoonshineCNAudit14.lean}).\end{status}

\subsection{Weight gradings of induced modules}

The Casimir argument needs the induced modules to be graded by Cartan weights, bounded below in height.  The library calls
such a grading \emph{admissible}:
\begin{leancode}
structure BorcherdsAdmissible (D : BorcherdsData g)
    (ρ : g →ₗ⁅ℂ⁆ Module.End ℂ M) where
  weightSpace : ℤ × ℤ → Submodule ℂ M
  weight_internal : DirectSum.IsInternal weightSpace
  action_weight : ∀ {α μ} {x : g} {v : M},
    x ∈ D.weightSpace α → v ∈ weightSpace μ → ρ x v ∈ weightSpace (μ + α)
  cartan_action : ∀ (ξ : ℂ × ℂ) {μ} {v : M}, v ∈ weightSpace μ →
    ρ (D.cartanMap ξ) v = borcherdsCartanPairing ξ (borcherdsWeightVector μ) • v
  lowerHeight : ℤ
  below_eq_bot : ∀ μ, borcherdsRootHeight μ < lowerHeight → weightSpace μ = ⊥
\end{leancode}
(\leanpath{VOA/Moonshine/MoonshineCNBorcherdsRootCutoff.lean}; the docstring: ``No finite-dimensional module-slice assumption is
included or inferred''.)  The files \leanpath{VOA/Moonshine/MoonshineCNBorcherdsCartanWeights.lean} and
\leanpath{VOA/Moonshine/MoonshineCNBorcherdsVermaWeights.lean} construct admissible gradings by simultaneous Cartan eigenspaces
(\lean{cartanEigenspace}, \lean{cartanAdmissible}): eigenspaces of distinct weights are independent
(\lean{cartanEigenspace_iSupIndep}), root vectors shift them (\lean{cartanEigenspace_action}), and an induced module is
generated from its bottom by $\chXVIIIu^+$ (\lean{inducedSpace_eq_top_of_bottom}), so its weights lie in $\{0\}\cup C$ (for $\mathrm{Ind}(\mathbb C)$,
\lean{trivialVermaAdmissible}) or in $C$ (for $\mathrm{Ind}(U)$, \lean{generatorVermaAdmissible}).  The kernel of a Cartan-intertwining map
out of such a module is spanned by weight vectors (\lean{cartanEigenspace_kernel_graded}).
\begin{physics}States of a Fock-like representation carry a definite momentum, bounded below in ``light-cone energy'' (height).\end{physics}

\section{The Casimir operator}\label{ch18:sec:casimir}

\subsection{Dual bases and root contractions}

By (BD3) the pairing $\chXVIIIg_\alpha\times\chXVIIIg_{-\alpha}\to\mathbb C$ is perfect (\lean{rootPairingEquiv}, in
\leanpath{VOA/Moonshine/MoonshineCNBorcherdsRootDuality.lean}).  So for every basis $(b_i)$ of $\chXVIIIg_\alpha$ there is a dual basis
$(b^i)$ of $\chXVIIIg_{-\alpha}$ (\lean{rootOppositeBasis}) with $(b_i,b^j)=\delta_{ij}$, and every contraction
$\sum_iB(b_i,b^i)$ by a bilinear map $B$ is independent of the basis (\lean{rootContraction_basis_independent}: it is the
image of the canonical element of $\chXVIIIg_\alpha\otimes\chXVIIIg_{-\alpha}$).  For a representation $\rho$ of $\chXVIIIg$ on $M$ put
$T_\alpha:=\sum_i\rho(b_i)\rho(b^i)$:
\begin{leancode}
def rootCasimirTerm (α : ℤ × ℤ) : End ℂ M :=
  ∑ i, rootActionContraction ρ (D.rootFiniteBasis α i)
    (D.rootOppositeBasis (D.rootFiniteBasis α) i)
theorem rootCasimirTerm_eq_basis {ι : Type*} [Fintype ι] (α : ℤ × ℤ)
    (b : Basis ι ℂ (D.weightSpace α)) :
    D.rootCasimirTerm ρ α = ∑ i, ρ (b i) * ρ (D.rootOppositeBasis b i) :=
\end{leancode}
(\leanpath{VOA/Moonshine/MoonshineCNBorcherdsCasimir.lean}; \lean{rootActionContraction ρ x y = ρ x * ρ y}.)

\subsection{The operator}

\begin{definition}[Casimir operator]\label{ch18:def:casimir}
On an admissible module $M$,
\begin{equation*}
  \Omega:=\underbrace{-\rho(h_1)\rho(h_2)-\rho(h_2)\rho(h_1)-2\rho(h_1)}_{\text{Cartan part}}+2\sum_{\alpha\in\Phi^+}T_\alpha,\qquad
  h_1=\iota(1,0),\ h_2=\iota(0,1),
\end{equation*}
the sum over positive roots being finite on each vector because $\rho(\chXVIIIg_{-\alpha})$ kills $M_\mu$ once $\mathrm{ht}\,\alpha$ exceeds
$\mathrm{ht}\,\mu-h_0$.
\end{definition}
\begin{leancode}
def cartanCasimirOperator : End ℂ M :=
  -(ρ (D.cartanMap (1, 0)) * ρ (D.cartanMap (0, 1))) -
    ρ (D.cartanMap (0, 1)) * ρ (D.cartanMap (1, 0)) -
      (2 : ℂ) • ρ (D.cartanMap (1, 0))
def positiveCasimirOperator : End ℂ M :=
  borcherdsLocallyFiniteRootSum (D.rootCasimirTerm ρ) A.rootCasimirTerms_finiteSupport
def casimirOperator : End ℂ M :=
  D.cartanCasimirOperator ρ + (2 : ℂ) • A.positiveCasimirOperator
\end{leancode}
(\leanpath{VOA/Moonshine/MoonshineCNBorcherdsCasimir.lean}; \lean{A : BorcherdsAdmissible D ρ};
\lean{borcherdsLocallyFiniteRootSum} from \leanpath{VOA/Moonshine/MoonshineCNBorcherdsRootCutoff.lean} evaluates the root
sum at each vector, where only finitely many terms are nonzero.)
\begin{dictionary}
\textbf{The Cartan part.}  On $M_\mu$, $\mu=(m,n)$, $\rho(h_1)$ acts by $(h_1,\mu)_{\mathrm{Bor}}=-n$ and $\rho(h_2)$ by $-m$, so the Cartan part
acts by $-2mn+2n=(\mu,\mu)_{\mathrm{Bor}}-2(\delta,\mu)_{\mathrm{Bor}}=Q(\mu)$ (\lean{cartanCasimirOperator_on_weight}).  The first two terms are the
contraction of the dual bases $\{h_1,h_2\}$, $\{-h_2,-h_1\}$ of $\chXVIIIg_0$ for $(\cdot,\cdot)_{\mathrm{Bor}}$; the third is $-2\rho(\iota\delta)$ with
$\delta=(1,0)$.  \textbf{Positive roots include $\alpha_r$} (height $0$; the root-cutoff docstring: ``the height-zero real root
is retained'').  This is Kac's Casimir (Kac, \emph{Infinite dimensional Lie algebras}, §2.5), with two differences: our modules
are \emph{lowest}-weight type (weights bounded below in height, and the positive-root terms are written $e_\alpha e_{-\alpha}$ so that
they kill lowest-weight vectors), whence the sign $-2(\delta,\cdot)$ instead of $+2(\delta,\cdot)$; and there are infinitely many simple
roots, which is why the root sum is only locally finite.
\end{dictionary}
\begin{physics}
The quadratic Casimir of the gauge algebra: a ``mass-squared'' operator, $Q(\mu)$ on a lowest-weight state of momentum $\mu$.
\end{physics}

\begin{theorem}[Lowest vectors and centrality]\label{ch18:thm:casimir}
On every admissible $M$: (a) if $v\in M_\mu$ is killed by $\rho(\chXVIIIg_{-\alpha})$ for every $\alpha\in\Phi^+$, then $\Omega v=Q(\mu)v$; and if
moreover $\mu\in C$, $v\neq0$ and $\Omega v=0$, then $\mu_1=1$; (b) $\Omega$ commutes with $\rho(z)$ for every $z\in\chXVIIIg$.
\end{theorem}
\begin{leancode}
theorem casimirOperator_lowest_vector {μ : ℤ × ℤ} {v : M}
    (hv : v ∈ A.weightSpace μ)
    (hlow : ∀ α, borcherdsPositiveRoot α → ∀ x ∈ D.weightSpace (-α), ρ x v = 0) :
    A.casimirOperator v = borcherdsCasimirScalar μ • v := by
theorem lowest_casimir_zero_first_coordinate {μ : ℤ × ℤ} {v : M}
    (hv : v ∈ A.weightSpace μ) (hv0 : v ≠ 0) (hμ : borcherdsPositiveCone μ)
    (hlow : ∀ α, borcherdsPositiveRoot α → ∀ x ∈ D.weightSpace (-α), ρ x v = 0)
    (hΩ : A.casimirOperator v = 0) : μ.1 = 1 := by
\end{leancode}
(\leanpath{VOA/Moonshine/MoonshineCNBorcherdsCasimir.lean}.)
\begin{leancode}
theorem casimirOperator_commutes (z : g) (v : M) :
    A.casimirOperator (ρ z v) = ρ z (A.casimirOperator v) := by
\end{leancode}
(\leanpath{VOA/Moonshine/MoonshineCNBorcherdsCasimirCentral.lean}.)
\begin{dictionary}
(a) needs no centrality (docstring: ``with no centrality hypothesis'').  (b) is for all $z$, by (BD5): the docstring
``BD5 generation extends the checked Cartan and simple-root commutators to all of g''.  Here (BD5) is a field of the
\emph{abstract} datum $D$; for the Monster it is \lean{hgen}.
\end{dictionary}
\begin{physics}The Casimir is a gauge-invariant mass operator; on lowest-weight states it is computed from the momentum alone.\end{physics}
\paragraph{Proof on paper.}  (a) $T_\alpha v=\sum_i\rho(b_i)\rho(b^i)v=0$ since $b^i\in\chXVIIIg_{-\alpha}$.  The second claim:
$Q(\mu)=-2n(m-1)$ with $n\ge1$.  (b) The set of $z$ with $[\Omega,\rho(z)]=0$ is a Lie subalgebra containing $\chXVIIIg_0$ ($\Omega$ preserves
weights), so by (BD5) it suffices to take $z\in\chXVIIIg_{\pm\alpha_i}$ for simple $\alpha_i$.  The dual-basis identity
(\lean{rootBracketContraction}, from invariance and nondegeneracy) gives
$[T_\alpha,\rho(z)]=A_\alpha(z)-A_{\alpha+\gamma}(z)$ with $A_\beta(z)=\sum_i\rho(b_i^{(\beta)})\rho([b^{i}_{(\beta)},z])$ for $z\in\chXVIIIg_\gamma$
(\lean{rootCasimirTerm_commutator}); summing over $\alpha\in\Phi^+$ telescopes (\lean{positiveCasimirOperator_commutator_finsum}) because
a simple root minus a positive root is never positive (\lean{borcherdsSimpleRoot_sub_not_positive}), leaving one boundary term:
$-\rho(z)\rho(\iota\alpha_i)$ for $z\in\chXVIIIg_{\alpha_i}$ and $\rho(\iota\alpha_i)\rho(z)$ for $z\in\chXVIIIg_{-\alpha_i}$
(\lean{positiveCasimirOperator_commutator_positive_simple}, \lean{_negative_simple}).  On $v\in M_\mu$ and $z\in\chXVIIIg_{\alpha_i}$:
\begin{equation*}
  [\Omega,\rho(z)]v=\bigl[Q(\mu+\alpha_i)-Q(\mu)-2(\alpha_i,\mu)_{\mathrm{Bor}}\bigr]\rho(z)v
  =\bigl[(\alpha_i,\alpha_i)_{\mathrm{Bor}}-2(\delta,\alpha_i)_{\mathrm{Bor}}\bigr]\rho(z)v=0,
\end{equation*}
because $\delta$ is a Weyl vector; similarly for $\chXVIIIg_{-\alpha_i}$ (\lean{borcherdsCasimirScalar_positive_simple_step},
\lean{_negative_simple_step}). \qed
\begin{status}No hypotheses beyond the datum and admissibility; standard axioms (\leanpath{VOA/Moonshine/MoonshineCNAudit15.lean}).
The supporting files \lean{BorcherdsRootDuality}, \lean{BorcherdsRootCutoff}, \lean{BorcherdsCasimirCommutator},
\lean{BorcherdsCasimirTelescoping}, \lean{BorcherdsCasimirBoundary} contain exactly the steps named in the proof; their finite-support
lemmas (\lean{rootCasimirTerms_finiteSupport}, \lean{shiftedCasimirBoundary_finiteSupport},
\lean{lowering_common_height_bound}) make the infinite root sums honest finite sums at each vector.\end{status}

\begin{theorem}[Casimir on induced modules]\label{ch18:thm:casimirinduced}
On an admissibly graded induced module, if $\Omega$ acts on the bottom $1\otimes W$ by a scalar $c$, it acts by $c$ everywhere.
\end{theorem}
\begin{leancode}
theorem inducedCasimir_scalar_of_bottom (D : BorcherdsData g)
    (A : BorcherdsAdmissible D (relativeInducedRep p σ)) (c : ℂ)
    (hbottom : ∀ w : W,
      A.casimirOperator (PBWTheorem.bottom p σ w) = c • PBWTheorem.bottom p σ w) :
    A.casimirOperator = c • LinearMap.id := by
\end{leancode}
(\leanpath{VOA/Moonshine/MoonshineCNBorcherdsCasimirInduction.lean}.)
\begin{physics}A gauge-invariant operator is constant on an irreducible-like multiplet generated from a vacuum.\end{physics}
\paragraph{Proof on paper.}  $\mathrm{Ind}(W)=\mathcal U(\chXVIIIg)(1\otimes W)$ and $\Omega$ commutes with $\mathcal U(\chXVIIIg)$ (Theorem~\ref{ch18:thm:casimir}(b)). \qed
\begin{status}Standard axioms (\leanpath{VOA/Moonshine/MoonshineCNAudit15.lean}).\end{status}

\section{The length-one resolution}\label{ch18:sec:resolution}

\begin{theorem}[Length-one resolution]\label{ch18:thm:resolution}
Let $D$ be a Borcherds datum with a normalized pair $(e,f)$.  Let $\mathrm{Ind}(\mathbb C)$ be induced from the trivial $\chXVIIIp^-$-module and $\mathrm{Ind}(U)$
from $U$ (with $\chXVIIIp^-$ acting through its Levi quotient).  The map $\Phi:\mathrm{Ind}(U)\to \mathrm{Ind}(\mathbb C)$, $[a\otimes u]\mapsto[au\otimes1]$, and the
augmentation $\epsilon:\mathrm{Ind}(\mathbb C)\to\mathbb C$ form an exact sequence of vector spaces
\begin{equation*}
  0\longrightarrow \mathrm{Ind}(U)\xrightarrow{\ \Phi\ }\mathrm{Ind}(\mathbb C)\xrightarrow{\ \epsilon\ }\mathbb C\longrightarrow0 .
\end{equation*}
Moreover $\Phi$ is a map of $\chXVIIIg$-modules (\lean{generatorVermaResolution_equivariant}).
Equivalently, right multiplication $\mathcal U(\chXVIIIu^+)\otimes U\to \mathcal U(\chXVIIIu^+)$ is injective with image the augmentation ideal.
\end{theorem}
\begin{leancode}
def generatorVermaResolution :
    relativeInducedSpace D.parabolicPart (D.genSpaceParabolicRep e f he hf hef) →ₗ[ℂ]
      relativeInducedSpace D.parabolicPart D.parabolicTrivialRep :=
theorem generatorVermaResolution_injective :
    Function.Injective (D.generatorVermaResolution e f he hf hef) := by
theorem generatorVermaResolution_shortExact :
    Function.Injective (D.generatorVermaResolution e f he hf hef) ∧
      Function.Exact (D.generatorVermaResolution e f he hf hef) D.trivialVermaAugmentation ∧
      Function.Surjective D.trivialVermaAugmentation := by
theorem positiveGeneratorMultiplication_shortExact :
    Function.Injective (envelopingGeneratorMultiplication (D.positiveGeneratorSpace e)) ∧
      Function.Exact (envelopingGeneratorMultiplication (D.positiveGeneratorSpace e))
        (envelopingAugmentation (g := D.positivePart)).toLinearMap ∧
      Function.Surjective (envelopingAugmentation (g := D.positivePart)) :=
\end{leancode}
(\leanpath{VOA/Moonshine/MoonshineCNBorcherdsResolution.lean}; section variables \lean{(e f : g) (he) (hf) (hef : D.form e f = 1)};
\lean{envelopingGeneratorMultiplication S} is $a\otimes u\mapsto a\,\iota(u)$, \leanpath{VOA/Moonshine/MoonshineCNBorcherdsResolutionRange.lean}.)
\begin{dictionary}
\lean{Function.Exact f g} means $\mathrm{im}\,f=\ker g$.  \lean{parabolicTrivialRep} is the zero representation on $\mathbb C$;
\lean{genSpaceParabolicRep e f he hf hef} is the $\chXVIIIp^-$-action on $U$ through the Levi projection.
\lean{trivialVermaAugmentation} is read off in PBW coordinates (\lean{trivialVermaAugmentation_verma}: it is $\epsilon(a)$ on
$[a\otimes1]$).  \textbf{Hypotheses}: the datum (including (BD5)) and the normalized pair, which always exists.
\end{dictionary}
\begin{physics}
In homological terms: $H_k(\chXVIIIu^+;\mathbb C)$ vanishes for $k\ge2$ and $H_1(\chXVIIIu^+;\mathbb C)\cong U$: the gauge algebra's positive cone has generators
$U$ and no relations.  In string terms, the Casimir (mass) of a would-be relation would have to vanish, which forces it to sit on the
line $m=1$ where there are no relations.
\end{physics}
\paragraph{Proof on paper.}  \emph{$\Phi$ is well defined.}  By the universal property of induction, a $\chXVIIIg$-map out of
$\mathrm{Ind}(U)$ is the same as a $\chXVIIIp^-$-map out of $U$.  The map $u\mapsto un_0=[u\otimes1]$ is $\chXVIIIp^-$-equivariant: for $p\in\chXVIIIp^-$,
$p\cdot un_0=[p,u]n_0+u\,p\,n_0=[p,u]n_0$, and $[p,u]$ differs from the action of $p$ on $U$ (through the Levi projection) by an element
of $\chXVIIIp^-$, which kills $n_0$ (\lean{genSpaceParabolicRep_sub_mem}; this uses $[\chXVIIIu^-,U]\subseteq\chXVIIIp^-$, Theorem~\ref{ch18:thm:Ustructure}(d)).
\emph{Image.}  $\mathrm{im}\,\Phi=\mathcal U(\chXVIIIg)(Un_0)=\mathcal U(\chXVIIIu^+)\mathcal U(\chXVIIIp^-)(Un_0)=\mathcal U(\chXVIIIu^+)Un_0$, which is
$\mathcal U(\chXVIIIu^+)_+n_0=\ker\epsilon$ by Theorem~\ref{ch18:thm:generation} and relative PBW (Theorem~\ref{ch18:thm:relpbw}).
\emph{Casimir vanishes on $\mathrm{Ind}(U)$.}  $\mathrm{Ind}(U)$ is generated from $1\otimes U$, and $1\otimes U$ from $1\otimes\bigoplus_j\chXVIIIg_{\alpha_j}$ by the Levi
action; $1\otimes v$, $v\in\chXVIIIg_{\alpha_j}$, has weight $\alpha_j$ and is killed by $\chXVIIIu^-$ (which acts by $0$ on $U$) and by $f$
($[f,v]=0$); so $\Omega(1\otimes v)=Q(1,j)(1\otimes v)=0$ (\lean{generatorVerma_casimir_bottom}, \lean{generatorVerma_casimir_eq_zero}).
\emph{Injectivity.}  $Z:=\ker\Phi$ is a graded submodule (\lean{cartanEigenspace_kernel_graded}) with weights in $C$.  If $Z\neq0$,
take the least height $h$ with $Z_{(h)}:=\bigoplus_{\mathrm{ht}\,\mu=h}Z_\mu\ne0$.  It is finite-dimensional (finitely many $\mu\in C$ of height $h$,
each weight space of $\mathrm{Ind}(U)\cong\mathcal \mathcal U(\chXVIIIu^+)\otimes U$ finite-dimensional).  Every lowering operator from $\chXVIIIg_{-\alpha}$ with $\alpha\in C$ lowers
the height, so kills $Z_{(h)}$ (the target is $Z_{(<h)}=0$).  The remaining lowering operator $f$ keeps the height and lowers the first
coordinate by one; since only finitely many first coordinates occur at height $h$, a power of $f$ vanishes on $Z_{(h)}$, so
$\ker f\cap Z_{(h)}\ne0$.  This kernel is stable under the Cartan ($[\iota\xi,f]$ is a multiple of $f$), so it contains a nonzero weight
vector $k$, of weight $\mu\in C$, killed by \emph{all} lowering operators.  By Theorem~\ref{ch18:thm:casimir}(a), $0=\Omega k=Q(\mu)k$, so
$\mu_1=1$ (\lean{submodule_eq_bot_of_casimir_zero}).  But on the line $m=1$ the module $\mathrm{Ind}(U)$ has only bottom vectors: $\mathrm{Ind}(U)=\mathcal
\mathcal U(\chXVIIIu^+)\otimes U$ (relative PBW), every homogeneous element of $\mathcal \mathcal U(\chXVIIIu^+)$ of nonzero degree raises the first coordinate by
at least one, and $U$ has first coordinates $\ge1$; so a weight vector with first coordinate $1$ is $1\otimes u$ with $u\in U_{(1,n)}=
\chXVIIIg_{\alpha_n}$ (\lean{generatorVerma_firstLine_bottom}).  On such vectors $\Phi$ is $u\mapsto un_0$, which is injective because
$\chXVIIIu^+\to\mathcal \mathcal U(\chXVIIIu^+)\cong \mathrm{Ind}(\mathbb C)$ is injective (\lean{generatorVermaResolution_firstLine}).  Contradiction. \qed
\begin{status}No hypotheses beyond the datum and the pair; standard axioms (\leanpath{VOA/Moonshine/MoonshineCNAudit17.lean}).\end{status}

\section{Freeness: \texorpdfstring{$T(U)\cong \mathcal U(\mathfrak u^+)$}{T(U) = U(u+)} and \texorpdfstring{$\mathfrak u^+$}{u+} is free on \texorpdfstring{$U$}{U}}\label{ch18:sec:free}

\subsection{The tensor algebra}

\begin{theorem}[$T(U)\cong \mathcal U(\chXVIIIu^+)$]\label{ch18:thm:tensor}
For any Lie algebra $L$ and subspace $S$: if right multiplication $\mathcal U(L)\otimes S\to \mathcal U(L)$ is injective, then $T(S)\to \mathcal U(L)$ is injective;
if $S$ generates $L$, $T(S)\to \mathcal U(L)$ is surjective.  For a Borcherds datum, $T(U)\to \mathcal U(\chXVIIIu^+)$ is an algebra isomorphism, natural in
every symmetry (BD6).
\end{theorem}
\begin{leancode}
theorem tensorGeneratorToEnveloping_injective_of_multiplication
    (hμ : Function.Injective (envelopingGeneratorMultiplication S)) :
    Function.Injective (tensorGeneratorToEnveloping S) := by
theorem positiveTensorGeneratorToEnveloping_bijective :
    Function.Bijective (tensorGeneratorToEnveloping (D.positiveGeneratorSpace e)) :=
def positiveTensorEnvelopingEquiv :
    TensorAlgebra ℂ (D.positiveGeneratorSpace e) ≃ₐ[ℂ]
      UniversalEnvelopingAlgebra ℂ D.positivePart :=
theorem positiveTensorEnvelopingEquiv_equivariant (σ : G)
    (t : TensorAlgebra ℂ (D.positiveGeneratorSpace e)) :
    D.positiveTensorEnvelopingEquiv e f he hf hef
        (tensorGeneratorMap _ ((D.positivePartAction A σ).toLinearMap.restrict
          (p := D.positiveGeneratorSpace e) (q := D.positiveGeneratorSpace e)
          (D.positiveGenerator_action_mem e he A σ)) t) =
      envelopingMap (D.positivePartAction A σ)
        (D.positiveTensorEnvelopingEquiv e f he hf hef t) :=
\end{leancode}
(\leanpath{VOA/Moonshine/MoonshineCNBorcherdsTensorEnveloping.lean}; \lean{tensorGeneratorToEnveloping S} is the algebra map
$T(S)\to \mathcal U(L)$ extending $S\hookrightarrow L\to \mathcal U(L)$, \leanpath{VOA/Moonshine/MoonshineCNBorcherdsResolutionRange.lean}.)
\begin{dictionary}
\lean{≃ₐ[ℂ]} is an algebra isomorphism.  The equivariance says: for $\sigma$ in a (BD6) group, applying $\sigma$ on generators and then the
isomorphism equals applying the isomorphism and then $\mathcal U(\sigma|_{\chXVIIIu^+})$.
\end{dictionary}
\begin{physics}Products of positive-cone gauge charges are freely generated by the $U$-charges: no relations at the level of the enveloping algebra.\end{physics}
\paragraph{Proof on paper.}  Filter $T(S)$ by degree (\lean{tensorDegreeFiltration}).  If $\tau(t)=0$ with $t\ne0$ of minimal degree,
write $t=c+\sum_it_iu_i$ with $c$ a scalar; applying the augmentation gives $c=0$, and then $\tau(t)=\mu(\sum\tau(t_i)\otimes u_i)=0$ forces
$\sum\tau(t_i)\otimes u_i=0$ by injectivity of $\mu$; taking the $u_i$ linearly independent (part of a basis of $S$), this forces every
$\tau(t_i)=0$, hence every $t_i=0$ by induction on the degree, hence $t=0$.  For Borcherds data,
injectivity of $\mu$ is Theorem~\ref{ch18:thm:resolution} and surjectivity of $\tau$ is Theorem~\ref{ch18:thm:generation}. \qed
\begin{status}No hypotheses beyond the datum and the pair; standard axioms (\leanpath{VOA/Moonshine/MoonshineCNAudit17.lean}).\end{status}

\subsection{The free Lie algebra}

\begin{theorem}[$\chXVIIIu^+$ is free on $U$]\label{ch18:thm:freelie}
Let $(u_i)_{i\in I}$ be the chosen basis of $U$ (\lean{Basis.ofVectorSpace}; by the general lemmas, any basis would do).  The Lie
algebra map from the free Lie algebra on $I$ to $\chXVIIIu^+$ sending the $i$-th generator to $u_i$ is an isomorphism, intertwining every
symmetry (BD6).
\end{theorem}
\begin{leancode}
def freeLieGeneratorMap : FreeLieAlgebra ℂ ι →ₗ⁅ℂ⁆ L :=
  FreeLieAlgebra.lift ℂ (fun i => (b i : L))
theorem freeLieGeneratorMap_surjective (hS : LieSubalgebra.lieSpan ℂ L (S : Set L) = ⊤) :
    Function.Surjective (freeLieGeneratorMap S b) := by
theorem freeLieGeneratorMap_injective
    (hτ : Function.Injective (tensorGeneratorToEnveloping S)) :
    Function.Injective (freeLieGeneratorMap S b) := by
def positiveGeneratorBasis :
    Basis (Basis.ofVectorSpaceIndex ℂ (D.positiveGeneratorSpace e)) ℂ
      (D.positiveGeneratorSpace e) :=
  Basis.ofVectorSpace ℂ (D.positiveGeneratorSpace e)
theorem positivePartFreeLie_bijective :
    Function.Bijective (freeLieGeneratorMap (D.positiveGeneratorSpace e)
      (D.positiveGeneratorBasis e)) :=
def positivePartFreeLieEquiv :
    FreeLieAlgebra ℂ (Basis.ofVectorSpaceIndex ℂ (D.positiveGeneratorSpace e)) ≃ₗ⁅ℂ⁆
      D.positivePart :=
theorem positivePartFreeLieEquiv_equivariant {G : Type*} [Group G] (A : D.GAction G) (σ : G)
    (x : FreeLieAlgebra ℂ (Basis.ofVectorSpaceIndex ℂ (D.positiveGeneratorSpace e))) :
    D.positivePartFreeLieEquiv e f he hf hef
        (freeLieInducedMap _ (D.positiveGeneratorBasis e) (D.positivePartAction A σ)
          (D.positiveGenerator_action_mem e he A σ) x) =
      D.positivePartAction A σ (D.positivePartFreeLieEquiv e f he hf hef x) :=
\end{leancode}
(\leanpath{VOA/Moonshine/MoonshineCNBorcherdsFreeLie.lean}; \lean{S : Submodule ℂ L}, \lean{b : Basis ι ℂ S} section variables.)
\begin{dictionary}
\textbf{Exactly what is proved about freeness.}  The Lie algebra that is shown to be free is \lean{D.positivePart}, i.e.\
$\chXVIIIu^+=\bigoplus_{m,n\ge1}\chXVIIIg_{(m,n)}$ (Borcherds coordinates): the positive cone, which is the positive part
$E=\chXVIIIg_{\alpha_r}\oplus\chXVIIIu^+$ \emph{minus the real simple root}.  The free generating space is $U$, the span of all
$(\mathrm{ad}\,e)^k\chXVIIIg_{\alpha_j}$, $0\le k<j$, \emph{not} $\bigoplus_j\chXVIIIg_{\alpha_j}$ alone.  The free Lie algebra is Mathlib's
\lean{FreeLieAlgebra ℂ ι} on the index set of an arbitrary basis of $U$ (\lean{Basis.ofVectorSpace}, chosen by
\lean{Classical.choice}); since a free Lie algebra on a basis of $U$ is the free Lie algebra on the vector space $U$, the choice is
immaterial, and the equivariance statement shows the isomorphism is compatible with symmetries via the induced map
\lean{freeLieInducedMap}.

\textbf{Injectivity of $L\to \mathcal U(L)$.}  \lean{freeLieGeneratorMap_injective} uses that the free Lie algebra embeds in its enveloping
algebra (\lean{envelopingEmbedding_injective}, Chapter~\ref{ch:17}, the PBW theorem) and Mathlib's
\lean{FreeLieAlgebra.universalEnvelopingEquivFreeAlgebra} with \lean{TensorAlgebra.equivFreeAlgebra}.

\textbf{Relation to the informal statement.}  One often reads ``$E$ minus the real root is a free Lie algebra''.  The Lean statement
is the precise version: $\chXVIIIu^+$ is free on $U$.  The variant ``the subalgebra generated by $\bigoplus_j\chXVIIIg_{\alpha_j}$ is free on it''
(Jurisich's form) is
\emph{not} formalized; it follows on paper from Theorem~\ref{ch18:thm:freelie} because $\bigoplus_j\chXVIIIg_{\alpha_j}=\bigoplus_jU_{j,0}$ is a graded
direct summand of $U$.
\end{dictionary}
\begin{physics}
The positive-cone gauge bosons are all nested commutators of the W-multiplets of the imaginary simple root states, with no relations
among them: a free Lie algebra, as for the ``Monster Lie algebra minus its $\mathfrak{gl}_2$''.
\end{physics}
\paragraph{Proof on paper.}  Surjective because $U$ generates $\chXVIIIu^+$ (Theorem~\ref{ch18:thm:generation}).  Injective: the composite
$\mathrm{Free}(I)\to\chXVIIIu^+\to \mathcal U(\chXVIIIu^+)$ equals $\mathrm{Free}(I)\hookrightarrow \mathcal U(\mathrm{Free}(I))\cong\mathrm{FreeAlg}(I)\cong T(U)\xrightarrow{\tau}\mathcal U(\chXVIIIu^+)$,
which is injective by PBW for the free Lie algebra and Theorem~\ref{ch18:thm:tensor}. \qed
\begin{status}No hypotheses beyond the datum and the pair; standard axioms (\leanpath{VOA/Moonshine/MoonshineCNAudit17.lean}).\end{status}

\begin{warning}
\textbf{$E$ itself is not free.}  $E=\bigoplus_{m>0}\chXVIIIg_{(m,n)}=\chXVIIIg_{\alpha_r}\oplus\chXVIIIu^+$ satisfies the relations $(\mathrm{ad}\,e)^j\chXVIIIg_{\alpha_j}=0$
(Theorem~\ref{ch18:thm:sl2}(b)); e.g.\ $[e,v]\in\chXVIIIg_{(2,0)}=0$ for $v\in\chXVIIIg_{\alpha_1}$.  And freeness on $\bigoplus_j\chXVIIIg_{\alpha_j}$ alone, although true for
the subalgebra it generates, would not give the denominator identity: its graded dimension is $1/(1-\sum_jc(j)pq^j)$, not the product side.  The Lean statement with generators $U$ is the one Chapter~\ref{ch:19} needs.
\end{warning}

\section{Consequences for the Monster Lie algebra}\label{ch18:sec:monsterconsequences}

\begin{theorem}[Structure of the Monster Lie algebra]\label{ch18:thm:monsterstructure}
Given simple-root generation (a theorem, \S\ref{ch18:sec:monsterbd}), and any $e\in\chXVIIIm^{\mathrm{Bor}}_{\alpha_r}$, $f\in\chXVIIIm^{\mathrm{Bor}}_{-\alpha_r}$ with $(e,f)=1$:
(a) $\dim U_{j,k}=\dim\Vnat_{1+j}$ for $0\le k<j$; (b) $0\to \mathrm{Ind}(U)\to \mathrm{Ind}(\mathbb C)\to\mathbb C\to0$ is exact; (c) $T(U)\to \mathcal U(\chXVIIIu^+)$ is bijective;
(d) $\chXVIIIu^+=\bigoplus_{m,n\ge1}\chXVIIIm^{\mathrm{Bor}}_{(m,n)}$ is free on $U$.
\end{theorem}
\begin{leancode}
theorem monsterLie_generatorString_finrank (j : ℤ) (hj : 1 ≤ j) (k : ℕ) (hk : (k : ℤ) < j) :
    Module.finrank ℂ ((monsterLieBorcherdsData hgen).generatorString e j k) =
      Module.finrank ℂ (VnatVOA.grading (1 + j)) := by
theorem monsterLie_generatorVermaResolution_shortExact :
    Function.Injective ((monsterLieBorcherdsData hgen).generatorVermaResolution e f he hf hef) ∧
      Function.Exact ((monsterLieBorcherdsData hgen).generatorVermaResolution e f he hf hef)
        (monsterLieBorcherdsData hgen).trivialVermaAugmentation ∧
      Function.Surjective (monsterLieBorcherdsData hgen).trivialVermaAugmentation :=
theorem monsterLie_positiveTensorEnveloping_bijective :
    Function.Bijective (tensorGeneratorToEnveloping
      ((monsterLieBorcherdsData hgen).positiveGeneratorSpace e)) :=
theorem monsterLie_positivePartFreeLie_bijective :
    Function.Bijective (freeLieGeneratorMap
      ((monsterLieBorcherdsData hgen).positiveGeneratorSpace e)
      ((monsterLieBorcherdsData hgen).positiveGeneratorBasis e)) :=
\end{leancode}
(\leanpath{VOA/Moonshine/MoonshineCNMonsterBorcherdsData.lean}; section variables
\lean{(hgen : MonsterLieSimpleRootGeneration) (e f : monsterLieAlgebra) (he : e ∈ monsterLieBorcherdsWeight borcherdsRealRoot)
(hf : f ∈ monsterLieBorcherdsWeight (-borcherdsRealRoot)) (hef : monsterLieBilinearForm e f = 1)}.)
\begin{dictionary}
Instances of Theorems~\ref{ch18:thm:Ustructure}, \ref{ch18:thm:resolution}, \ref{ch18:thm:tensor}, \ref{ch18:thm:freelie} at
\lean{monsterLieBorcherdsData hgen}.  (a) combines \lean{generatorString_finrank} with \lean{monsterLieBorcherdsRootEquiv} at $\gamma=\alpha_j$:
$1+1\cdot j=1+j$.  So $U_{(a,b)}\cong\Vnat_{a+b}$, $\dim U_{(a,b)}=c(a+b-1)$, as needed for the product side $1-\sum c(a+b-1)p^aq^b$.
\end{dictionary}
\begin{physics}The Monster string's gauge algebra has a free positive cone, generated by the W-multiplets of the states $\Vnat_{1+j}\otimes e^{(1,-j)}$.\end{physics}
\begin{status}
Conditional on \lean{hgen} in this file; unconditional after substituting \lean{M8.monsterLieSimpleRootGeneration}.  Standard axioms
(\leanpath{VOA/Moonshine/MoonshineCNAudit17.lean}).
\end{status}

\section{What the chapter proves, and what it does not}\label{ch18:sec:scope}

\begin{itemize}
\item \textbf{Proved, no hypotheses:} the abstract no-ghost theorem (under its stated hypotheses $c=26$ and (F1)); the concrete
  light-cone splitting; the $\Vnat$ no-ghost isometry at every nonzero momentum and the zero-momentum hyperbolic plane; the Lie algebra
  $\chXVIIIm$, its $\chXVIIIII$-grading, root spaces $\Vnat_{1-xy}$ / $\mathbb C^2$, $\Aut(\Vnat)$-action, invariant nondegenerate form with lattice
  normalisation on $\chXVIIIm_0$, Cartan action; the whole Lie theory of Borcherds data.
\item \textbf{Proved, conditional on \lean{hgen}} in the files of this chapter, \textbf{discharged} in Chapter~\ref{ch:20}: $\chXVIIIm$ is a
  Borcherds datum; its resolution, $T(U)\cong \mathcal U(\chXVIIIu^+)$, freeness.
\item \textbf{Not formalized:} the BRST cohomology description; positivity (unitarity) of the form on $\Vnat$ or $\chXVIIIm$; that $\chXVIIIm$ is a
  generalized Kac--Moody algebra in Borcherds' axiomatic sense, or any Weyl--Kac--Borcherds character formula as a Lie-theoretic
  theorem (the denominator identity is proved analytically in Chapter~\ref{ch:19}); Jurisich's freeness on $\bigoplus_j\chXVIIIm_{\alpha_j}$; an
  isometry statement for $\chXVIIIm_r\cong\Vnat_{1-r^2/2}$ at the Lie-algebra level (only at the sector level).
\item \textbf{Trust model:} every headline is \texttt{\#print axioms}-clean with the three standard axioms; the former project axiom
  \lean{VOA.pbw} is a theorem and is not an axiom of any declaration here.
\end{itemize}

\paragraph{Coverage of the assigned files.}  The 75 files assigned to this chapter are covered as follows.
\emph{Individually (headlines with Lean quotations):} \texttt{StringModule}, \texttt{StringOrthogonals}, \texttt{StringGenerators},
\texttt{StringCritical}, \texttt{StringLongitudinal}, \texttt{StringGrading}, \texttt{StringTransverse}, \texttt{StringPhysical},
\texttt{StringQuotient}, \texttt{StringComparison}, \texttt{StringStraightening}, \texttt{VnatSectorOperators}, \texttt{VnatSectorLaws},
\texttt{VnatSectorGrading}, \texttt{VnatLightCone}, \texttt{VnatStringModule}, \texttt{VnatFormAction}, \texttt{VnatSectorAction},
\texttt{VnatSectorEvaluation}, \texttt{VnatLongitudinalQuotient}, \texttt{VnatNoGhost}, \texttt{VnatNoGhostConsequences},
\texttt{VnatZeroWeight}, \texttt{VnatZeroQuotient}, \texttt{PhysicalLie}, \texttt{PhysicalMaps}, \texttt{MonsterLie}, \texttt{MonsterGrading},
\texttt{SectorVirasoroComparison}, \texttt{SectorPhysicalTransfer}, \texttt{RootSpaces}, \texttt{RootSpaceAction}, \texttt{RootRadical},
\texttt{SectorFormConvention}, \texttt{SectorFormComparison}, \texttt{MonsterForm}, \texttt{MonsterPairing}, \texttt{MonsterRadical},
\texttt{MonsterZeroForm}, \texttt{MonsterCartanAction}, \texttt{MonsterLieInvariance}, \texttt{BorcherdsRoots}, \texttt{BorcherdsData},
\texttt{MonsterBorcherdsData}, \texttt{BorcherdsSl2}, \texttt{BorcherdsGeneratorWeights}, \texttt{BorcherdsStringSum},
\texttt{BorcherdsStringAction}, \texttt{BorcherdsParabolic}, \texttt{BorcherdsLevi}, \texttt{BorcherdsGeneration},
\texttt{BorcherdsAugmentation}, \texttt{RelativeInduction}, \texttt{RelativePBW}, \texttt{VermaMultiplication}, \texttt{BorcherdsRootCutoff},
\texttt{BorcherdsCasimir}, \texttt{BorcherdsCasimirCentral}, \texttt{BorcherdsCasimirInduction}, \texttt{BorcherdsResolution},
\texttt{BorcherdsResolutionRange}, \texttt{BorcherdsTensorEnveloping}, \texttt{BorcherdsFreeLie} (all with prefix \texttt{MoonshineCN}).
\emph{Collectively (technical support, described where used):} \texttt{StringSymmetry} (equivariance of the quotient maps,
\S\ref{ch18:sec:abstractnoghost}); \texttt{SectorFrame}, \texttt{SectorOscillators}, \texttt{SectorVirasoroOperators} (the frame change and the
explicit lattice-sector Virasoro operators behind Theorem~\ref{ch18:thm:sectorcomparison}); \texttt{LieProjection},
\texttt{BorcherdsParabolicAction} (Levi projection and the $\chXVIIIp^-$-action on $U$, \S\ref{ch18:sec:strings});
\texttt{BorcherdsCartanWeights}, \texttt{BorcherdsVermaWeights} (admissible gradings, \S\ref{ch18:sec:verma}); \texttt{BorcherdsRootDuality},
\texttt{BorcherdsCasimirCommutator}, \texttt{BorcherdsCasimirTelescoping}, \texttt{BorcherdsCasimirBoundary} (the Casimir computation,
\S\ref{ch18:sec:casimir}).

\section{Referee's checklist}\label{ch18:sec:checklist}

\begin{longtable}{p{0.30\textwidth}p{0.33\textwidth}p{0.18\textwidth}p{0.1\textwidth}}
\toprule
declaration & intended meaning & hypotheses & check \\ \midrule
\endhead
\lean{StringModule} & abstract string module (S1)--(S4) & — (structure) & Def~\ref{ch18:def:stringmodule}\\
\lean{spurious}, \lean{physical}, \lean{kVacuum}, \lean{transverse}, \lean{LightConeSplit} & $\chXVIIIS,\chXVIIIP,\chXVIIIK,\chXVIIIT$, (F1) & — & Def~\ref{ch18:def:derived}\\
\lean{orthogonal_spurious_eq} etc. & $(\chXVIIIS^N)^\perp=\chXVIIIP^N$, etc. & none & Thm~\ref{ch18:thm:orthogonals}\\
\lean{physicalOne_radical_eq} & $\mathrm{rad}\,\chXVIIIP^1=\chXVIIIP^1\cap\chXVIIIS$ & none & Thm~\ref{ch18:thm:radicals}\\
\lean{spurious_eq_sup_two} & $\chXVIIIS=L_{-1}H+L_{-2}H$ & none & Thm~\ref{ch18:thm:twogenerators}\\
\lean{stringNegativeWord_mem_sortedSpan} & Virasoro straightening & neg.\ commutators & Thm~\ref{ch18:thm:straightening}\\
\lean{tildeTwo_commutator_weight_neg_one} & defect $c/2-13$ & weight $-1$ & Thm~\ref{ch18:thm:critical}\\
\lean{criticalSpurious_L_two_tilde} & $\widetilde L_2\chXVIIIS^1\subseteq\chXVIIIS$ & $c=26$ & Thm~\ref{ch18:thm:critical}\\
\lean{physical_kSpurious_isCompl}, \lean{kVacuum_eq_transverse_sup_radical}, \lean{transverse_nondegenerate} & (A1)--(A4) & (F1) & Thm~\ref{ch18:thm:A14}\\
\lean{physicalOne_eq_transverse_sup_null} & (A5) $\chXVIIIP^1=\chXVIIIT^1\oplus\mathrm{null}$ & $c=26$, (F1) & Thm~\ref{ch18:thm:A5}\\
\lean{transversePhysicalQuotientEquiv}, \lean{physicalKQuotientEquiv} & $\chXVIIIT^1\cong\chXVIIIP^1/\mathrm{rad}\cong\chXVIIIK^1/\mathrm{rad}$ isometric, equivariant & $c=26$, (F1) & Thm~\ref{ch18:thm:quotientequivs}\\
\lean{vnatHypMode_virasoro} & sector Virasoro, $c=26$ & none & Thm~\ref{ch18:thm:sectorlaws}\\
\lean{vnatStringModule}, \lean{stringSector} & concrete string sector & none & Def~\ref{ch18:def:vnatstringmodule}\\
\lean{vnatHypTensor_lightConeSplit}, \lean{stringSector_lightConeSplit} & (F1) for the sector & $(x,y)\neq0$ & Thm~\ref{ch18:thm:F1}\\
\lean{vnatAut_preserves_natForm} & $\Aut\Vnat$ preserves $B^\natural$ & none & Thm~\ref{ch18:thm:sectoraction}\\
\lean{vnatKVacuumQuotientEquiv} & $\chXVIIIK^N/\mathrm{rad}\cong\Vnat_{N-xy}$ & none & Thm~\ref{ch18:thm:kquotient}\\
\lean{stringSectorNoGhostEquiv} (+ isometry, equivariant) & \textbf{no-ghost}: $\chXVIIIP^1/\mathrm{rad}\cong\Vnat_{1-xy}$ & $(x,y)\neq0$ & Thm~\ref{ch18:thm:noghost}\\
\lean{vnatZeroNoGhostEquiv} & zero momentum: $\mathbb C^2$, hyperbolic & none & Thm~\ref{ch18:thm:zeromomentum}\\
\lean{primaryLieAlgebra} & $P^1/(P^1\cap\mathrm{Sp})$ is a Lie algebra & none & Thm~\ref{ch18:thm:primarylie}\\
\lean{monsterLieAlgebra} & the Monster Lie algebra & — (def.) & Def~\ref{ch18:def:monsterlie}\\
\lean{monsterLieMomentum_isInternal} & $\chXVIIIII$-grading & none & Thm~\ref{ch18:thm:grading}\\
\lean{monsterLieAut} & $\Aut\Vnat\to\Aut_{\mathrm{Lie}}\chXVIIIm$ & none & Thm~\ref{ch18:thm:autaction}\\
\lean{moonshineMatterSector_spurious_iff}, \lean{_primary_iff} & sector $=$ vertex-algebra momentum sector & none & Thm~\ref{ch18:thm:sectorcomparison}\\
\lean{monsterLieRootSpaceEquiv} (+ equivariant) & $\chXVIIIm_{(x,y)}\cong\Vnat_{1-xy}$ & $(x,y)\neq0$ & Thm~\ref{ch18:thm:rootspaces}\\
\lean{monsterLieZeroMomentumEquiv} & $\chXVIIIm_0\cong\mathbb C^2$ & none & Thm~\ref{ch18:thm:zerospace}\\
\lean{moonshineMatterSector_no_same_momentum_isometry} & dead end certificate & $\alpha\ne0$, $c\ne0$ & Thm~\ref{ch18:thm:nosame}\\
\lean{moonshineMatterSector_form} & reversal comparison & none & Thm~\ref{ch18:thm:reversal}\\
\lean{monsterLieBilinearForm} & form on $\chXVIIIm$ & — (def.) & Def~\ref{ch18:def:monsterform}\\
\lean{monsterLieBilinearForm_nondegenerate}, \lean{moonshineMatterPrimary_radical} & nondegenerate; radical $=$ null & none & Thm~\ref{ch18:thm:primaryradical}\\
\lean{monsterLieBilinearForm_zeroMomentum} & lattice form on $\chXVIIIm_0$ & none & Thm~\ref{ch18:thm:zeroform}\\
\lean{monsterLie_zeroMomentum_bracket} & Cartan action $(\lambda(h),\beta)$ & none & Thm~\ref{ch18:thm:cartanaction}\\
\lean{monsterLieBilinearForm_invariant} & invariance & none & Thm~\ref{ch18:thm:invariance}\\
\lean{BorcherdsData} & (BD1)--(BD5) & — (structure) & Def~\ref{ch18:def:bd}\\
\lean{real_sl2_relations} & real $\mathfrak{sl}_2$ & datum, pair & Thm~\ref{ch18:thm:bdbasic}\\
\lean{monsterLieBorcherdsRootEquiv} & $\chXVIIIm^{\mathrm{Bor}}_{(m,n)}\cong\Vnat_{1+mn}$ & $\gamma\ne0$ & Thm~\ref{ch18:thm:borcherdsroot}\\
\lean{monsterLieBorcherdsData} & $\chXVIIIm$ is a Borcherds datum & \lean{hgen} (discharged, Ch.~\ref{ch:20}) & Thm~\ref{ch18:thm:monsterbd}\\
\lean{adReal_pow_injective}, \lean{genSpaceDirectSumEquiv} & real-root strings, $U=\bigoplus U_{j,k}$ & datum, pair & Thms~\ref{ch18:thm:sl2}, \ref{ch18:thm:Ustructure}\\
\lean{positivePart_eq_lieSpan_genSpace} & $\chXVIIIu^+=\mathrm{Lie}(U)$ & datum, pair & Thm~\ref{ch18:thm:generation}\\
\lean{relativePBWEquiv}, \lean{parabolicVermaMultiplication_bijective} & relative PBW & complement & Thm~\ref{ch18:thm:relpbw}\\
\lean{casimirOperator_lowest_vector}, \lean{casimirOperator_commutes} & Casimir & admissible & Thm~\ref{ch18:thm:casimir}\\
\lean{generatorVermaResolution_shortExact} & $0\to \mathrm{Ind}(U)\to \mathrm{Ind}(\mathbb C)\to\mathbb C\to0$ & datum, pair & Thm~\ref{ch18:thm:resolution}\\
\lean{positiveTensorGeneratorToEnveloping_bijective} & $T(U)\cong \mathcal U(\chXVIIIu^+)$ & datum, pair & Thm~\ref{ch18:thm:tensor}\\
\lean{positivePartFreeLie_bijective} & $\chXVIIIu^+$ free on $U$ & datum, pair & Thm~\ref{ch18:thm:freelie}\\
\lean{monsterLie_positivePartFreeLie_bijective} etc. & the above for $\chXVIIIm$ & \lean{hgen}, pair & Thm~\ref{ch18:thm:monsterstructure}\\
\bottomrule
\end{longtable}

\begin{status}
Every row is \texttt{\#print axioms}-clean (standard axioms, no \lean{sorryAx}) in the audit modules
\leanpath{VOA/Moonshine/MoonshineCNAudit4.lean}--\leanpath{VOA/Moonshine/MoonshineCNAudit17.lean}.  ``Pair'' means a normalized real pair
$(e,f)$, which always exists (\lean{exists_normalized_real_pair}).  The only non-structural hypothesis in the chapter, \lean{hgen}, is discharged
by \lean{M8.monsterLieSimpleRootGeneration} (no hypotheses; \leanpath{VOA/Moonshine/MoonshineCNFinalAudit.lean}).  These statements were
\emph{not} re-built for this chapter: the outputs are those reported by the maintainers' targeted builds of the audit modules, and
no aggregate build of the whole library is claimed.
\end{status}

\section{Notes and references}\label{ch18:sec:notes}

The no-ghost theorem is due to P.~Goddard and C.~B.~Thorn, \emph{Compatibility of the dual Pomeron with unitarity and the absence of
ghosts in the dual resonance model}, Phys.\ Lett.\ B 40 (1972) 235--238; R.~C.~Brower, \emph{Spectrum-generating algebra and no-ghost
theorem for the dual model}, Phys.\ Rev.\ D 6 (1972) 1655--1662, gave the independent argument via DDF operators.  The BRST
treatment is in M.~Kato and K.~Ogawa, \emph{Covariant quantization of string based on BRS invariance}, Nucl.\ Phys.\ B 212 (1983) 443--460,
and in I.~B.~Frenkel, H.~Garland and G.~J.~Zuckerman, \emph{Semi-infinite cohomology and string theory}, Proc.\ Natl.\ Acad.\ Sci.\ USA 83
(1986) 8442--8446.  The physical background is in M.~B.~Green, J.~H.~Schwarz and E.~Witten, \emph{Superstring Theory}, Vol.~1, Cambridge
University Press, 1987, Chapter~2, and J.~Polchinski, \emph{String Theory}, Vol.~1, Cambridge University Press, 1998, Chapters~4 and~6.

Generalized Kac--Moody (Borcherds) algebras were introduced by R.~E.~Borcherds, \emph{Generalized Kac--Moody algebras}, J.\ Algebra 115
(1988) 501--512.  The Monster Lie algebra, its construction from $\Vnat\otimes V_{\chXVIIIII}$ via the no-ghost theorem, the denominator
identity and the proof of the Conway--Norton conjecture are in R.~E.~Borcherds, \emph{Monstrous moonshine and monstrous Lie superalgebras},
Invent.\ Math.\ 109 (1992) 405--444; §5 there sketches Goddard--Thorn's proof in the form used here, and §§6--7 the Lie theory.  The
freeness of the subalgebra generated by the imaginary simple root spaces is due to E.~Jurisich, \emph{Generalized Kac--Moody Lie algebras,
free Lie algebras and the structure of the Monster Lie algebra}, J.\ Pure Appl.\ Algebra 126 (1998) 233--266; see also E.~Jurisich,
J.~Lepowsky and R.~L.~Wilson, \emph{Realizations of the Monster Lie algebra}, Selecta Math.\ (N.S.) 1 (1995) 129--161.  A careful expository
account is R.~W.~Gebert, \emph{Introduction to vertex algebras, Borcherds algebras and the Monster Lie algebra}, Internat.\ J.\ Modern Phys.\
A 8 (1993) 5441--5503 (arXiv:hep-th/9308151).  The Casimir operator and its use for resolutions follow V.~G.~Kac, \emph{Infinite dimensional
Lie algebras}, 3rd ed., Cambridge University Press, 1990, §§2.4--2.6 and §9.11 (and the Garland--Lepowsky/Kostant homology theory that
lies behind it).  Invariant forms on vertex algebras: I.~B.~Frenkel, Y.-Z.~Huang and J.~Lepowsky, \emph{On axiomatic approaches to vertex
operator algebras and modules}, Mem.\ Amer.\ Math.\ Soc.\ 104 (1993), no.~494.

The development followed the authors' working notes, which are not part of the release: the no-ghost argument
restructures Goddard--Thorn into the single hypothesis (F1), proved by a leading-term argument that avoids the Kac
determinant, and the Borcherds-data Lie theory uses a Casimir route to the length-one resolution (with the recorded dead
end ``generation $\Leftrightarrow$ product formula'').  Those notes are unrefereed; where they and the Lean source differ (e.g.\ the
light-cone charges are $(\kappa,\ell)=(x,y)$ in Lean rather than $(y,x)$; the radical identification holds without $c=26$), this chapter
follows the source.
